\documentclass[12pt]{book}

% ------------------------------
% PAQUETES DE IDIOMA Y ENCODING
% ------------------------------
\usepackage[utf8]{inputenc}

% ------------------------------
% PAQUETES MATEMÁTICOS
% ------------------------------
\usepackage{amsmath}
\usepackage{amsthm}
\usepackage{amssymb}
\usepackage{amsfonts}
\usepackage{amsbsy}
\usepackage{amstext}
\usepackage{mathrsfs}
\usepackage[mathscr]{euscript}
\usepackage{mathabx}

% ------------------------------
% PAQUETES DE DIBUJO Y DIAGRAMAS
% ------------------------------
\usepackage{tikz}
\usetikzlibrary{positioning}
\usepackage{tikz-cd}
\usepackage[all]{xy}

% ------------------------------
% PAQUETES DE FORMATO Y GRÁFICOS
% ------------------------------
\usepackage{graphicx}
\usepackage{float}
\usepackage{multirow}
\usepackage{rotating}
\usepackage{anysize}
\usepackage[left=2cm,right=2cm,top=2cm,bottom=2cm]{geometry}
\usepackage{newcent}
\usepackage{tcolorbox}
\usepackage{eqparbox}

% ------------------------------
% OTROS PAQUETES ÚTILES
% ------------------------------
\PassOptionsToPackage{normalem}{ulem}
\usepackage{ulem}
\usepackage{enumerate}

% ------------------------------
% CONFIGURACIONES DE CAPÍTULOS Y NOMBRES
% ------------------------------
\renewcommand{\chaptername}{Capítulo}
\renewcommand{\bibname}{Bibliografía}
\renewcommand{\contentsname}{Índice general}

% ------------------------------
% ENTORNOS DE TEOREMAS
% ------------------------------
\newtheorem{theorem}{Teorema}[section]
\newtheorem{proposition}[theorem]{Proposición}
\newtheorem{lemma}[theorem]{Lema}
\newtheorem{corollary}[theorem]{Corolario}
\newtheorem{definition}[theorem]{Definición}
\newtheorem{example}{Ejemplo}
\newtheorem{remark}[theorem]{Observación}
\newtheorem{remarks}[theorem]{Observaciones}

% ------------------------------
% COMANDOS GENERALES
% ------------------------------
\newcommand{\ra}{\rightarrow}
\newcommand{\rla}{\leftrightarrow}
\newcommand{\lra}{\leftrightarrow}
\newcommand{\Ra}{\Rightarrow}
\newcommand{\we}{\wedge}
\newcommand{\sq}{\square}
\newcommand{\up}[1]{[{#1})}
\newcommand{\down}[1]{({#1}]}
\newcommand{\downc}[1]{({#1}]^{c}}
\newcommand{\Cl}{\mathrm{Cl}}
\newcommand{\Int}{\mathrm{int}}
\newcommand{\Fg}{\mathnormal{F_g}}
\newcommand{\Ig}{\mathnormal{I_g}}

% ------------------------------
% COMANDOS DE ÁLGEBRAS Y OPERADORES
% ------------------------------
\newcommand{\Con}{\mathrm{Con}}
\newcommand{\Fil}{\mathrm{Fil}}
\newcommand{\F}{\mathrm{F}}
\newcommand{\Fo}{\mathrm{F}^{o}}
\newcommand{\Fom}{\mathrm{F}^{o}_{m}}
\newcommand{\Eo}{\mathrm{P1}}
\newcommand{\Et}{\mathrm{P2}}
\newcommand{\Eth}{\mathrm{P3}}
\newcommand{\Q}{\mathrm{P}}
\newcommand{\K}{\mathrm{K}}
\newcommand{\C}{\mathrm{C}}
\newcommand{\V}{\mathrm{V}}
\newcommand{\I}{\mathbb{I}}
\newcommand{\Pro}{\mathbb{P}}
\newcommand{\Sub}{\mathbb{S}}
\newcommand{\Do}{\mathrm{Do}}
\newcommand{\Hom}{\mathrm{Hom}}
\newcommand{\Id}{\mathrm{Id}}

% ------------------------------
% COMANDOS DE SISTEMAS LÓGICOS
% ------------------------------
\newcommand{\vd}{\vdash}
\newcommand{\dv}{\dashv}
\newcommand{\vD}{\vDash}
\newcommand{\Gs}{\mid\thicksim}
\newcommand{\GS}{\mid\thickapprox}

% ------------------------------
% CATEGORÍAS Y FUNTORES
% ------------------------------
\newcommand{\BDL}{\mathsf{BDL}}
\newcommand{\PS}{\mathsf{PS}}
\newcommand{\Spec}{\mathsf{Spec}}
\newcommand{\SSpec}{\mathsf{pSpec}}
\newcommand{\CB}{\mathsf{BS}}
\newcommand{\CP}{\mathsf{SRLS}}
\newcommand{\PWH}{\mathsf{SRLS}}
\newcommand{\CanOp}{\mathsf{CanOp}}
\newcommand{\StarCan}{\mathsf{StarCan}}
\newcommand{\Ock}{\mathsf{Ockham}}

% ------------------------------
% FUNTORES Y VARIANTES
% ------------------------------
\newcommand{\X}{\mathfrak{X}}
\newcommand{\hX}{\hat{\mathrm{X}}}
\newcommand{\D}{\mathfrak{D}}
\newcommand{\hD}{\hat{\mathrm{D}}}
\newcommand{\Up}{\mathbb{UP}}
\newcommand{\UpX}{\mathrm{Up}(X)}
\newcommand{\tp}{\tau^{\prime}}
\newcommand{\G}{\mathrm{G}}
\newcommand{\DS}{\mathrm{D}_{\mathrm{S}}}

% ------------------------------
% OTROS COMANDOS MATEMÁTICOS
% ------------------------------
\newcommand{\KO}{\mathrm{KO}}
\newcommand{\Upi}{\mathrm{Up}_i}
\newcommand{\mWG}{\mathrm{MWG}}
\newcommand{\Xom}{\mathrm{X}^{o}_{m}}
\newcommand{\WH}{\mathrm{WH}}
\newcommand{\RWH}{\mathrm{RWH}}
\newcommand{\SRL}{\mathrm{SRL}}
\newcommand{\Hey}{\mathrm{Hey}}
\newcommand{\KRWH}{\mathrm{KRWH}}
\newcommand{\PKRWH}{\mathrm{PKRWH}}
\newcommand{\KSRL}{\mathrm{KSRL}}
\newcommand{\PKSRL}{\mathrm{PKSRL}}
\newcommand{\ID}{\mathrm{M}}
\newcommand{\hsD}{\hat{\mathscr{D}}}
\newcommand{\hsX}{\hat{\mathscr{X}}}

% ------------------------------
% FORMATO FINAL
% ------------------------------
\makeatother






\begin{document}





\begin{titlepage}
\begin{center}
\includegraphics[scale=0.1]{Imágenes/Logo Unicen.png}
\end{center}
\begin{center}
\begin{LARGE}
\textbf{\sffamily{UNIVERSIDAD NACIONAL DEL CENTRO DE LA PROVINCIA DE BUENOS AIRES}} \\
\end{LARGE}
\end{center}
\vspace*{0.07in}
\begin{center}
\begin{large}
\textbf{\sffamily{FACULTAD DE CIENCIAS EXACTAS}}\\
\end{large}
\end{center}
\vspace*{0.20in}
\begin{center}
\begin{large}
\textbf{\sffamily{DEPARTAMENTO DE MATEMÁTICA}}\\
\end{large}
\end{center}
\vspace*{0.3in}

\begin{center}
\rule{130mm}{0.1mm}\\
\end{center}

\vspace*{0.5in}
\begin{center}
\begin{large}
\underline{\textbf{\sffamily{Trabajo de Tesis de Licenciatura}}}\\
\vspace{1cm}
\sffamily{Dualidades para retículos distributivos acotados con operadores normales}
\end{large}
\end{center}
\vspace*{0.3in}

\begin{center}
\rule{130mm}{0.1mm}\\
\end{center}

\vspace*{0.1in}
\begin{center}
\begin{large}
\textbf{\sffamily{Tesista:}} \sffamily{Ayrton Porto} \\
\vspace{0.3cm}
\textbf{\sffamily{Director:}} \sffamily{Dr. William Javier Zuluaga Botero} \\
\vspace{0.3cm}
\textbf{\sffamily{Codirector:}} \sffamily{Dr. Agustín Leonel Nagy} \\
\vspace{0.3cm}
\textbf{\sffamily{Año:}} $2025$ 
\end{large}
\end{center}
\end{titlepage}




\chapter*{Agradecimientos}

Este trabajo no lo hice solo, así que quiero aprovechar este espacio para agradecer a quienes me acompañaron durante estos años.
\vspace{0.2cm}

A mis directores, el Dr. William Zuluaga y el Dr. Agustín Nagy, por la paciencia infinita para orientarme. Desde que toqué la puerta para ver en que trabajaban, le han dedicado tiempo a este proceso que fue escribir la tesis.
\vspace{0.2cm}

A mi familia, que estuvo firme en todo momento. Gracias por el apoyo constante, por entender mis tiempos raros y por levantarme el ánimo cuando parecía que no avanzaba más. Si llegué hasta acá, es en gran parte por ustedes.
\vspace{0.2cm}

A Martha, mi novia, por estar ahí de la forma más simple y más linda posible. Por la compañía, por las palabras justas y por hacer que este proceso fuera mucho menos caótico. También agradecerle a su familia que me ha acompañado. 
\vspace{0.2cm}

A mis amigos, por las charlas, los descansos necesarios y por aparecer cuando yo estaba enterrado entre PDFs, chocolatadas y LaTeX. Gracias por recordarme que existía el mundo exterior.
\vspace{0.2cm}


Al Departamento de Matemática por la formación, el apoyo y el buen trato durante todo este recorrido.
\vspace{0.2cm}

A todos ustedes, gracias.


\tableofcontents


\chapter*{Introducción general}


El estudio de la lógica puede abordarse esencialmente desde dos perspectivas complementarias. Por un lado, el \emph{enfoque sintáctico}, que analiza las reglas deductivas y las propiedades formales de un sistema lógico sin recurrir a interpretaciones externas. Por otro lado, el \emph{enfoque semántico}, cuyo objetivo es asociar las fórmulas bien formadas con estructuras que les otorguen significado.
\vspace{0.2cm}

Dentro de este enfoque semántico, los retículos distributivos acotados ocupan un lugar central, ya que aparecen como reductos de diversas álgebras que sirven como semánticas algebraicas de múltiples sistemas lógicos. Asimismo, las representaciones topológicas de estructuras algebraicas desempeñan un papel fundamental en la comprensión de la lógica. Las distintas representaciones disponibles para las álgebras que sirven como semánticas algebraicas de una lógica dada permiten, en muchos casos, desarrollar semánticas relacionales formuladas en términos de conjuntos y de relaciones definidas sobre ellos. 
\vspace{0.2cm}



El objetivo del presente trabajo es desarrollar un estudio sobre ciertos operadores unarios definidos sobre retículos distributivos acotados desde dos enfoques diferentes. Por un lado, se desarrolla una dualidad de tipo Priestley para retículos distributivos acotados dotados de los operadores unarios mencionados, generalizando la dualidad de Priestley clásica. Por el otro, se estudia la noción de extensión canónica como herramienta para establecer representaciones sobre álgebras de conjuntos para retículos distributivos acotados dotados de operadores unarios que satisfacen ciertas propiedades.



Este trabajo está compuesto por cuatro capítulos, cuyos contenidos se detallan a continuación:

\begin{enumerate}[\normalfont 1)]

\item El primer capítulo presenta las nociones preliminares necesarias para el desarrollo de la tesis. Se introducen definiciones básicas y resultados clásicos sobre conjuntos ordenados, teoría de retículos, álgebra universal, topología y teoría de categorías. Asimismo, se recordará la dualidad de Priestley para retículos distributivos acotados, la cual será fundamental para nuestro trabajo.


\item El segundo capítulo introduce un marco general de \emph{dualidad tipo Priestley} para retículos distributivos acotados equipados con operadores unarios normales. Se definen las categorías algebraicas y topológicas adecuadas, se construyen funtores duales entre ellas, y se analizan fenómenos de simetría que permiten 


\item El tercer capítulo está dedicado al estudio de la \emph{extensión canónica} para retículos distributivos acotados dotados de operadores unarios normales. En este capítulo se analizan sus propiedades básicas, se describe su comportamiento sobre elementos abiertos y cerrados, y se establece cómo esta construcción permite obtener una representación relacional de los operadores. Para cada tipo de operador considerado, se obtiene una representación relacional que coincide con la obtenida mediante las dualidades del capítulo anterior, mostrando así la conexión entre ambos enfoques.


\item En el cuarto capítulo se presentan \emph{aplicaciones} de las dualidades desarrolladas. Se estudia como primera aplicación en particular el caso de las álgebras de Ockham como subcategorías de las categorías trabajadas en el capítulo \ref{Ch 3} y cómo ciertas propiedades ecuacionales de los operadores estudiados (preservación de uniones, idempotencia, conmutatividad) se
traducen en propiedades relacionales de los espacios obtenidos (clausura, reflexividad, transitividad).


\end{enumerate}

\chapter{Preliminares}

En el presente capítulo vamos a introducir las nociones elementales, la notación y los resultados básicos que serán utilizados a los largo de este trabajo. Teniendo en cuenta que el presente capítulo tiene un carácter introductorio, no serán demostrados cada uno de los resultados citados, sino que remitiremos al lector interesado a los diversos trabajos científicos y libros que se citan en cada sección.

\section{Conjuntos, relaciones y órdenes}

Comenzaremos por introducir las definiciones y la notación elemental sobre teoría de conjuntos, relaciones binarias y relaciones de orden.

\vspace{0.2cm}

Sea \( X \) un conjunto no vacío. Denotaremos por \( \mathcal{P}(X) \) 
al conjunto de todos los subconjuntos de \( X \). Para \( U,V \in \mathcal{P}(X) \) escribimos \( U \cap V \) y \( U \cup V \) para la intersección y unión, respectivamente. Asimismo, 
\[
U^{c} = \{ x \in X : x \notin U \} \hspace{0.2cm} \text{ y } \hspace{0.2cm}
U \setminus V = \{ x \in X : x \in U \text{ y } x \notin V \},
\]
para indicar el complemento y la diferencia de conjuntos, respectivamente.

\vspace{0.2cm}

Sea \( \{X_i\}_{i=1}^{n} \) una familia finita de conjuntos. Indicaremos su producto cartesiano de la siguiente manera:
\[
\prod_{i=1}^{n} X_i = \{ (x_1,\dots,x_n) : x_i \in X_i \text{ para cada } 1 \le i \le n \}.
\]
En particular, si \( X_i = X \) para todo \( i \), escribimos \( X^n \) en lugar de
\( \prod_{i=1}^{n} X \).

\vspace{0.2cm}

Sean \( X,Y \) conjuntos. Una \emph{relación binaria} \( R \) definida en el producto \( X \times Y \) es un subconjunto \( R \subseteq X \times Y \). En particular, si \( X = Y \) diremos que \( R \) es una relación binaria definida sobre \( X \).
Asimismo,  dados  \( U \subseteq X \) y \( V \subseteq Y \) definimos
\[
R(U) = \{ y \in Y : (x,y) \in R \text{ para algún } x \in U \} \hspace{0.1cm} \text{ y } \hspace{0.1cm}
R^{-1}(V) = \{ x \in X : R(x) \cap V \neq \emptyset \}.
\]

En particular, si $U = \{x\}$ y $V= \{y\}$ escribiremos $R(x)$ y $R^{-1}(y)$  en lugar de $R(\{x\})$ y $R^{-1}(\{y\})$, respectivamente. 
\vspace{0.2cm}

Sean $R,S$ dos relaciones definidas sobre $X$. La \emph{composición} de $R$ y $S$ es una nueva relación binaria que se define de la siguiente manera:
\[
R \circ S = \{ (x,z) \in X \times X : \exists\, y \in X,\ (x,y) \in R \text{ y } (y,z) \in S \}.
\]

\medskip

Sea $R$ una relacion binaria definida sobre un conjunto $X$. Diremos que $R$ es un \emph{orden parcial} si es reflexiva, antisimétrica y transitiva. Cuando no haya ambigüedad, denotaremos el orden parcial por \( \leq \) y escribiremos \( x \leq y \) para \( (x,y)\, \in \,\ \leq \).
Llamaremos conjunto parcialmente ordenado a cualquier par \( (X,\leq) \), donde $X$ es un conjunto no vacío y $\leq$ es una relacion de orden parcial definida sobre $X$).

El orden es \emph{total} si para todo \( x,y \in X \) se cumple \( x \le y \) o \( y \le x \).

\medskip

Sea \( (X,\leq) \) un conjunto parcialmente ordenado y \( U \subseteq X \).
Decimos que \( U \) es un \emph{conjunto creciente} (o \emph{upset}) si
para todo \( x,y \in X \), de \( x \in U \) y \( x \le y \) se sigue \( y \in U \).
Denotamos por  $\mathrm{Up}(X)$  el conjunto de todos los subconjuntos crecientes de \( X \).
Dualmente, \( V \subseteq X \) es un \emph{conjunto decreciente} (o \emph{downset}) si
de \( x \in V \) y \( y \le x \) se sigue \( y \in V \); su familia se denota \( \Do(X) \).

Para \( U \subseteq X \) definimos el \emph{creciente} y el \emph{decreciente generados} por \( U \):
\[
[U) = \{ y \in X : \exists\, u \in U,\ u \le y \},
\qquad
(U] = \{ y \in X : \exists\, u \in U,\ y \le u \}.
\]
 En particular, si $U = \{x\}$ escribiremos $[x)$ y $(x]$ en lugar de $[\{x\})$ y $(\{x\}]$.

\medskip

Una \emph{función} \( f \colon X \to Y \) es una relación \( f \subseteq X \times Y \)
tal que para cada \( x \in X \) existe un único \( y \in Y \) con \( (x,y) \in f \). La denotaremos como \( f \colon X \to Y \). Si \( U \subseteq X \) indicaremos la imagen del conjunto \( U \) por medio de \( f \) como 
\[
f[U] = \{ y \in Y : \exists\, x \in U,\ f(x) = y \},
\]
Si \( V \subseteq Y \) indicaremos la imagen inversa de \( V \) por medio de \( f \) como
\[
f^{-1}[V] = \{ x \in X : f(x) \in V \}.
\]

Sean \( (X,\leq) \) e \( (Y,\leq) \) dos conjuntos parcialmente ordenados y
\( f \colon X \to Y \) una función. La función $f$ se dirá:
\begin{itemize}
  \item \emph{monótona} si \( a \le b \) implica \( f(a) \le f(b) \);
  \item \emph{antimonótona}  si \( a \le b \) implica \( f(b) \le f(a) \);
  \item \emph{isomorfismo de orden} si es biyectiva y \( a \le b \iff f(a) \le f(b) \);
  \item \emph{isomorfismo de orden dual} si es biyectiva y \( a \le b \iff f(b) \le f(a) \).
\end{itemize}


\section{Álgebra universal}

En esta sección recordaremos algunas nociones básicas de álgebra universal que serán de gran utilidad en los capítulos centrales de este trabajo. Para un tratamiento más detallado de los conceptos mencionados, remitimos al lector los textos clásicos \cite{BD,BS}.

\begin{definition}\label{def:lenguaje-algebraico}
Un \emph{lenguaje algebraico} (o \emph{tipo algebraico}) es un par \( (\mathcal{L},\mathrm{ar}) \),
donde \( \mathcal{L} \) es un conjunto de símbolos y \( \mathrm{ar}\colon \mathcal{L}\to\mathbb{N}_{0} \) es una función, llamada \emph{aridad}, que asigna a cada símbolo \( f\in\mathcal{L} \) un número entero \( \mathrm{ar}(f) \).
Aquí \( \mathbb{N}_0=\{0,1,2,\dots\} \) denota el conjunto de los enteros naturales que incluyen al cero.
\end{definition}


\begin{definition}\label{def:algebra}
Un \emph{álgebra} \( \mathbf{A} \) de tipo \( \mathcal{L} \) es un par
\( \mathbf{A}=(A,\mathcal{L}) \), donde \( A \) es un conjunto no vacío, el cual llamaremos \emph{universo},
y para cada símbolo \( f\in\mathcal{L} \) de aridad \( n \),
se especifica una operación (u operador) \( f^{\mathbf{A}}\colon A^{n}\to A \).
\end{definition}

Sea $\mathbf{A} = (A,\mathcal{L})$ un álgebra de tipo $\mathcal{L}$. En caso de que el tipo algebraico sea claro por el contexto, denotaremos simplemente el álgebra por \( \mathbf{A} \) y su universo por \( A \) en lugar de $\mathbf{A} = (A,\mathcal{L})$.

\begin{definition}\label{def:subalgebra}
Sean \( \mathbf{A},\mathbf{B} \) álgebras de tipo \( \mathcal{L} \).
Diremos que \( \mathbf{B} \) es una \emph{subálgebra} de \( \mathbf{A} \)
si \( B\subseteq A \) y para toda operación \( f\in\mathcal{L} \) y todo
\( (b_{1},\dots,b_{n})\in B^{n} \) se cumple
\[
f^{\mathbf{B}}(b_{1},\dots,b_{n}) = f^{\mathbf{A}}(b_{1},\dots,b_{n}).
\]
Equivalentemente, \( B \) es un \emph{subuniverso} de \( A \) si es cerrado bajo las operaciones del lenguaje.
En tal caso escribimos \( \mathbf{B}\leq \mathbf{A} \).
\end{definition}

\begin{definition}\label{def:homomorfismo}
Una función \( \alpha\colon A\to B \) es un \emph{homomorfismo} de álgebras
si para toda operación \( f\in\mathcal{L} \) de aridad \( n \) y todo
\( a_{1},\dots,a_{n}\in A \) se cumple
\[
\alpha(f^{\mathbf{A}}(a_{1},\dots,a_{n})) = f^{\mathbf{B}}(\alpha(a_{1}),\dots,\alpha(a_{n})).
\]
Denotaremos por \( \Hom(\mathbf{A},\mathbf{B}) \) el conjunto de todos los homomorfismos.  
El homomorfismo identidad de \( \mathbf{A} \) se denota \( \Id_{\mathbf{A}} \).  

Diremos que:
\begin{itemize}
  \item \(\alpha\) es un \emph{monomorfismo} si es inyectiva; en tal caso se escribe
  \( \alpha\colon \mathbf{A} \hookrightarrow \mathbf{B} \).
  \item \(\alpha\) es un \emph{epimorfismo} si es sobreyectiva.
  \item \(\alpha\) es un \emph{isomorfismo} si es biyectiva, en cuyo caso
  \( \mathbf{A}\cong\mathbf{B} \).
  \item \(\alpha\) es un \emph{endomorfismo} si \( \mathbf{A}=\mathbf{B} \).
\end{itemize}
\end{definition}

\begin{definition}\label{def:producto-algebras}
Sea \( I \) un conjunto de índices y \( \{\mathbf{A}_{i}\}_{i\in I} \) una familia de álgebras del mismo tipo.
Su \emph{producto} \( \prod_{i\in I}\mathbf{A}_{i} \) es el álgebra de universo
\( \prod_{i\in I}A_{i} \), donde las operaciones se definen coordenada a coordenada:
\[
f^{\prod \mathbf{A}_{i}}(a_{1},\dots,a_{n})(i) = 
f^{\mathbf{A}_{i}}(a_{1}(i),\dots,a_{n}(i)).
\]
Para cada \( j\in I \), la proyección sobre la \( j \)-ésima coordenada
\(\pi_{j}\colon \prod_{i\in I}\mathbf{A}_{i}\to \mathbf{A}_{j}\) es un homomorfismo.
\end{definition}

\section{Topología}

En esta sección recordaremos las nociones topológicas básicas que serán necesarias en los capítulos siguientes.
Para una exposición detallada remitimos a \cite{Munkres}.

\begin{definition}\label{def:topologia}
Sea \( X \) un conjunto no vacío.  
Una \emph{topología} sobre \( X \) es una familia \( \tau \subseteq \mathcal{P}(X) \)
que cumple:
\begin{enumerate}
  \item \( \emptyset,\, X \in \tau \);
  \item si \( \{U_i\}_{i\in I} \subseteq \tau \), entonces \( \bigcup_{i\in I} U_i \in \tau \);
  \item si \( U_1,\dots,U_k \in \tau \), entonces \( \bigcap_{i=1}^{k} U_i \in \tau \).
\end{enumerate}
El par \( (X,\tau) \) se llama \emph{espacio topológico}.
\end{definition}

Sean $(X,\tau)$ un espacio topológico y $U$ un subconjunto de $X$. Diremos que $U$ es un conjunto $\tau$-abierto, o simplemente abierto, si  \( U \in \tau \).  Un subconjunto \( V\subseteq X \) se dirá \emph{cerrado} si \( V^{c}\) es abierto. Un conjunto que es simultáneamente abierto y cerrado se denomina \emph{clopen}.

\begin{definition}\label{def:clausura-interior}
Sea \( (X,\tau) \) un espacio topológico y \( A\subseteq X \).
La \emph{clausura} de \( A \) en \( (X,\tau) \) se define como
\[
Cl_{\tau}(A)
   = \bigcap \{V \subseteq X : A\subseteq V,\ V \text{ es cerrado}\},
\]
y el \emph{interior} de \( A \) como
\[
Int_{\tau}(A)
   = \bigcup \{O\in\tau : O\subseteq A\}.
\]
Cuando no haya necesidad de hacer
referencia concreta a $\tau$, escribiremos simplemente \( Cl(A) \) e \( Int(A) \) en lugar de $Cl_{\tau}(A)$ y $Int_{\tau}(A)$, respectivamente.
\end{definition}

\begin{definition}\label{def:base-generadora}
Sea $X$ un conjunto no vacío. Una familia $\mathcal{B}\subseteq\mathcal{P}(X)$
se dice \emph{base} para una topología en $X$ si satisface:

\begin{enumerate}
  \item $\displaystyle \bigcup_{B\in\mathcal{B}} B = X$;
  \item para cualesquiera $B_1,B_2\in\mathcal{B}$ y $x\in B_1\cap B_2$,
        existe $B_3\in\mathcal{B}$ tal que $x\in B_3\subseteq B_1\cap B_2$.
\end{enumerate}

La \emph{topología generada por $\mathcal{B}$} es la familia de todas
las uniones arbitrarias de elementos de $\mathcal{B}$.
\end{definition}

\begin{definition}\label{def:base-espacio}
Sea $(X,\tau)$ un espacio topológico. Una familia
$\mathcal{B}\subseteq\tau$ se dice \emph{base de $(X,\tau)$} si cumple:

\begin{enumerate}
  \item $\mathcal{B}\subseteq\tau$;
  \item para todo $U\in\tau$ y todo $x\in U$,
        existe $B\in\mathcal{B}$ tal que $x\in B\subseteq U$.
\end{enumerate}
\end{definition}



\begin{definition}\label{def:subbase}
Una \emph{subbase} para una topología \( \tau \) sobre \( X \)
es una familia \( \mathscr{S}\subseteq\mathcal{P}(X) \)
tal que \( \bigcup_{S\in\mathscr{S}} S = X \).
La topología generada por \( \mathscr{S} \) es la de todos los conjuntos \( U\subseteq X \)
que cumplen:  
para cada \( x\in U \) existen \( S_1,\dots,S_n\in\mathscr{S} \)
con \( x\in S_1\cap\dots\cap S_n\subseteq U \).
\end{definition}

\begin{definition}\label{def:cubrimiento-abierto}
Sea $(X,\tau)$ un espacio topológico.  
Una familia $\{U_i\}_{i\in I}\subseteq\tau$ se denomina 
\emph{cubrimiento abierto} de $X$ si
\[
X = \bigcup_{i\in I} U_i.
\]
\end{definition}

\begin{definition}\label{def:compacidad}
Un espacio topológico $(X,\tau)$ es \emph{compacto} si 
para todo cubrimiento abierto $\{U_i\}_{i\in I}$ de $X$ 
existe un subcubrimiento finito, es decir,
\[
X = U_{i_1}\cup\cdots\cup U_{i_n}
\quad\text{para algunos}\quad i_1,\dots,i_n\in I.
\]
\end{definition}

\begin{definition}\label{def:compacidad-subespacio}
Sea $(X,\tau)$ un espacio topológico y $V\subseteq X$ un subespacio
con la topología inducida $\tau_V$.  
Decimos que $V$ es \emph{compacto} si $(V,\tau_V)$ es compacto.
\end{definition}

\medskip

Notemos que, para un espacio topológico $(X,\tau)$, hemos definido la noción de compacidad para el espacio completo.  
Sin embargo, el siguiente lema proporciona un criterio equivalente para determinar la compacidad de un subconjunto cualquiera de $X$.

\begin{lemma}[{\cite[Lema~26.1]{Munkres}}]\label{lem:compacto-subespacio-Munkres}
Sean $(X,\tau)$ un espacio topológico y $V \subseteq X$. Las siguientes afirmaciones son equivalentes:
\begin{enumerate}
    \item $(V,\tau_{V})$ es un subespacio compacto.
    \item Para todo cubrimiento $\{U_{i}\}_{i \in I}$ de $V$ por $\tau$-abiertos, 
    existe un subcubrimiento finito $U_{i_1},\dots,U_{i_n}$ de $\tau$-abiertos de $V$.
\end{enumerate}
\end{lemma}

\begin{definition}\label{def:espacios-T}
Sea \( (X,\tau) \) un espacio topológico.  
Diremos que:
\begin{itemize}
  \item \( (X,\tau) \) es \emph{\(T_0\)} si para cualesquiera \( x\ne y\in X \)
        existe \( U\in\tau \) tal que \( x\in U,\, y\notin U \)
        o bien \( y\in U,\, x\notin U \);
  \item \( (X,\tau) \) es \emph{\(T_1\)} si para cualesquiera \( x\ne y\in X \)
        existen abiertos \( U,V\in\tau \) con \( x\in U,\ y\notin U \)
        y \( y\in V,\ x\notin V \);
  \item \( (X,\tau) \) es \emph{\(T_2\)} (o \emph{Hausdorff})
        si para todo \( x\ne y\in X \) existen \( U,V\in\tau \)
        con \( x\in U \), \( y\in V \) y \( U\cap V=\emptyset \).
\end{itemize}
\end{definition}

\begin{definition}\label{def:continuidad-homeomorfismo}
Sean $(X,\tau)$ y $(Y,\tau')$ espacios topológicos y 
$f\colon X\to Y$ una función. Diremos que $f$ es \emph{continua} si, para todo abierto 
$U\in\tau'$, la preimagen $f^{-1}[U]$ es un abierto de $(X,\tau)$.

Por otro lado, diremos que $f$ es un \emph{homeomorfismo} si es 
biyectiva, continua y su inversa $f^{-1}$ también es continua.  
En tal caso, los espacios $(X,\tau)$ y $(Y,\tau')$ se dicen 
\emph{homeomorfos}.
\end{definition}



\section{Categorías}

Recordaremos brevemente las nociones básicas de teoría de categorías que serán utilizadas a lo largo de este trabajo.  
Para una exposición más general referimos al lector interesado a los siguientes textos \cite{Adamek}.

\begin{definition}\label{def:categoria}
Una \emph{categoría} \( \mathsf{C} \) consta de los siguientes datos:
\begin{enumerate}[\normalfont (1)]
  \item una colección denominada \( Ob(\mathsf{C}) \), cuyos elementos se llaman \emph{objetos};
  \item una colección denominada \( Fle(\mathsf{C}) \), cuyos elementos se llaman \emph{morfismos} o \emph{flechas};
  \item dos aplicaciones \( Dom,Cod \colon Fle(\mathsf{C}) \to Ob(\mathsf{C}) \), que asignan a cada morfismo su \emph{dominio} y su \emph{codominio};
  \item una operación de composición que, dada una pareja de morfismos \( f,g \) tales que \( Cod(f)=Dom(g) \), produce un nuevo morfismo \( g\circ f \) con \( Dom(g\circ f)=Dom(f) \) y \( Cod(g\circ f)=Cod(g) \);
  \item para cada objeto \( x \) existe un morfismo identidad \( \Id_{x}\colon x\to x \);
  \item la composición es asociativa y las identidades actúan como neutros, es decir,
  \[
  h\circ(g\circ f)=(h\circ g)\circ f,
  \qquad
  \Id_{Cod(f)}\circ f=f=f\circ \Id_{Dom(f)}.
  \]
\end{enumerate}
\end{definition}

Si \( f \) es un morfismo de \( \mathsf{C} \) con dominio \( x=Dom(f) \) y codominio \( y=Cod(f) \), se escribe \( f\colon x\to y \).  
El conjunto de todos los morfismos de \( \mathsf{C} \) que van de \( x \) a \( y \) se denota \( \mathsf{C}(x,y) \).


\begin{definition}
Sea \( \mathsf{C} \) una categoría y \( f\colon x\to y \) un morfismo.  
Se dice que \( f \) es un \emph{isomorfismo} si existe un morfismo \( g\colon y\to x \) tal que
\[
g\circ f=\Id_{x},\qquad f\circ g=\Id_{y}.
\]
En tal caso los objetos \( x \) e \( y \) se dicen \emph{isomorfos} en \( \mathsf{C} \).
\end{definition}

\begin{definition}
A toda categoría \( \mathsf{C} \) se le asocia su \emph{categoría opuesta} \( \mathsf{C}^{op} \), que posee los mismos objetos y los mismos morfismos, pero con las direcciones invertidas.  
Formalmente,
\[
\mathsf{C}^{op}(x,y) = \mathsf{C}(y,x),
\qquad
g^{op}\circ f^{op} = (f\circ g)^{op}.
\]
\end{definition}

\begin{definition}
Sea $\mathsf{C}$ una categoría.  
Una \emph{subcategoría} $\mathsf{D}$ de $\mathsf{C}$ consta de:

\begin{enumerate}[\normalfont (a)]
  \item una colección de objetos $\mathrm{Ob}(\mathsf{D})$, donde cada objeto de $\mathsf{D}$ es también un objeto de $\mathsf{C}$;

  \item para cada par de objetos $X,Y$ de $\mathsf{D}$, una colección de flechas
  $\mathsf{D}(X,Y)$, consistente exclusivamente en flechas de $\mathsf{C}(X,Y)$;

  \item para cada objeto $X$ de $\mathsf{D}$, la flecha identidad en $\mathsf{D}$ coincide con la identidad correspondiente en $\mathsf{C}$;

  \item la composición en $\mathsf{D}$ es la restricción de la composición de $\mathsf{C}$:
  si $f\colon X\to Y$ y $g\colon Y\to Z$ son flechas de $\mathsf{D}$, entonces 
  $g\circ f$ (calculada en $\mathsf{C}$) es una flecha de $\mathsf{D}$.
\end{enumerate}

La subcategoría $\mathsf{D}$ se dice \emph{plena} si, además, para cada par de objetos
$X,Y$ de $\mathsf{D}$, toda flecha de $\mathsf{C}$ con dominio $X$ y codominio $Y$
pertenece también a la colección $\mathsf{D}(X,Y)$.
\end{definition}



\begin{definition}\label{def: funtor}
Sean \( \mathsf{C} \) y \( \mathsf{D} \) categorías.
Un \emph{funtor covariante} es un mapeo \( F\colon\mathsf{C}\to\mathsf{D} \) que asigna:
\begin{enumerate}[\normalfont (1)]
  \item a cada objeto \( x \) de \( \mathsf{C} \), un objeto \( F(x) \) de \( \mathsf{D} \);
  \item a cada morfismo \( f\colon x\to y \) de \( \mathsf{C} \), un morfismo \( F(f)\colon F(x)\to F(y) \) de \( \mathsf{D} \);
\end{enumerate}
que preserva identidades y composición, es decir:
\[
F(\Id_{x})=\Id_{F(x)},
\qquad
F(g\circ f)=F(g)\circ F(f).
\]
Si \( F \) invierte la dirección de los morfismos, de manera que
de un morfismo \( f\colon x\to y \) en \( \mathsf{C} \)
produce un morfismo \( F(f)\colon F(y)\to F(x) \) en \( \mathsf{D} \),
entonces \( F \) se llama \emph{funtor contravariante}.
En tal caso se lo escribe \( F\colon \mathsf{C}\to\mathsf{D}^{op} \).
\end{definition}

\begin{definition}
Dada una categoría \( \mathsf{C} \), se define el \emph{funtor identidad}
\[
\mathrm{I}_{\mathsf{C}}\colon \mathsf{C}\to\mathsf{C},
\qquad
\mathrm{I}_{\mathsf{C}}(x)=x,\quad \mathrm{I}_{\mathsf{C}}(f)=f.
\]
\end{definition}


\begin{definition}\label{def:transf-nat}
Sean $F,G \colon \mathsf{C} \to \mathsf{D}$ dos funtores covariantes. 
Una \emph{transformación natural} $\theta \colon F \Rightarrow G$
es una asignación que a cada objeto $X$ de $\mathsf{C}$ le
asocia un morfismo
\[
\theta_X \colon F(X) \longrightarrow G(X)
\]
en la categoría $\mathsf{D}$, de manera tal que para cada
morfismo $f \colon X \to Y$ en $\mathsf{C}$, el siguiente diagrama conmuta:
\[
\begin{tikzcd}
F(X) \arrow[r, "\theta_X"] \arrow[d, "F(f)"'] &
G(X) \arrow[d, "G(f)"] \\
F(Y) \arrow[r, "\theta_Y"'] &
G(Y).
\end{tikzcd}
\]
Es decir, 
\[
G(f)\circ \theta_X \;=\; \theta_Y \circ F(f).
\]
\end{definition}


\begin{definition}
Sean \( \mathsf{C} \) y \( \mathsf{D} \) categorías.
Una \emph{equivalencia de categorías} entre \( \mathsf{C} \) y \( \mathsf{D} \)
está dada por un par de funtores covariantes
\[
F\colon\mathsf{C}\to\mathsf{D},
\qquad
G\colon\mathsf{D}\to\mathsf{C},
\]
y por dos transformaciones naturales
\[
\theta\colon \mathrm{I}_{\mathsf{C}}\Rightarrow G\circ F,
\qquad
\phi\colon \mathrm{I}_{\mathsf{D}}\Rightarrow F\circ G,
\]
tales que cada morfismo componente \( \theta_{C} \) y \( \phi_{D} \) es un isomorfismo.
En este caso se dice que \( \mathsf{C} \) y \( \mathsf{D} \) son \emph{equivalentes}.
\end{definition}

\begin{definition}\label{def:equivalencia-dual}
Sean \( \mathsf{C} \) y \( \mathsf{D} \) dos categorías.
Diremos que \( \mathsf{C} \) y \( \mathsf{D} \) son \emph{dualmente equivalentes}
si existe una equivalencia de categorías entre \( \mathsf{C} \) y \( \mathsf{D}^{op} \) (o entre \( \mathsf{C}^{op} \) y \( \mathsf{D} \)). En este caso se dice que \( \mathsf{C} \) y \( \mathsf{D} \) son categorías \emph{duales}.
\end{definition}

\begin{definition}\label{def:isomorfismo-categorias}
Sean \( \mathsf{C} \) y \( \mathsf{D} \) dos categorías.
Decimos que \( \mathsf{C} \) y \( \mathsf{D} \) son \emph{isomorfas}
si existen funtores
\[
F\colon \mathsf{C} \longrightarrow \mathsf{D},
\qquad
G\colon \mathsf{D} \longrightarrow \mathsf{C},
\]
tales que se cumplen las igualdades estrictas
\[
F\circ G = \mathrm{I}_{\mathsf{D}},
\qquad
G\circ F = \mathrm{I}_{\mathsf{C}}.
\]
En este caso se escribe
\[
\mathsf{C}\;\cong\;\mathsf{D}.
\]
\end{definition}




\section{Retículos}

Los retículos, y en particular los retículos distributivos acotados, constituyen el reducto algebraico de la mayoría de las estructuras que aparecerán a lo largo de esta tesis.  
En esta sección introduciremos sus propiedades básicas y algunos resultados clásicos que utilizaremos posteriormente.  
Para un estudio más detallado referimos al lector a \cite{DP, SS}.

\begin{definition}
Un \emph{retículo} es un álgebra \( \mathbb{L}=(L,\wedge,\vee) \) de tipo \((2,2)\) tal que, para cualesquiera \( a,b,c \in L \), se cumplen las siguientes igualdades:
\begin{enumerate}[\normalfont (1)]
  \item \( a \wedge (b \wedge c) = (a \wedge b) \wedge c \) y \( a \vee (b \vee c) = (a \vee b) \vee c \);
  \item \( a \wedge b = b \wedge a \) y \( a \vee b = b \vee a \);
  \item \( a \wedge a = a \) y \( a \vee a = a \);
  \item \( a \wedge (b \vee a) = a \) y \( a \vee (b \wedge a) = a \).
\end{enumerate}
\end{definition}

La noción de retículo está estrechamente vinculada a la de conjunto parcialmente ordenado.  
Si \( \mathbb{L}=(L,\wedge,\vee) \) es un retículo, se define sobre \( L \) una relación de orden mediante las operaciones del lenguaje algebraico de la siguiente manera:
\[
a \leq b
\quad\Longleftrightarrow\quad
a = a \wedge b
\quad\Longleftrightarrow\quad
b = b \vee a.
\]
Recíprocamente, si \( (L,\leq) \) es un conjunto parcialmente ordenado que admite ínfimo y supremo para cada par de elementos, la estructura \( (L,\wedge,\vee) \) dada por
\[
a\wedge b = \inf\{a,b\}, \qquad a\vee b = \sup\{a,b\},
\]
es un retículo.  
Así, un retículo puede considerarse indistintamente como un álgebra de tipo \((2,2)\) o como un conjunto parcialmente ordenado para el cual siempre existen ínfimos y supremos de subconjuntos finitos.

\begin{definition}
Un retículo \( \mathbb{L}=(L,\wedge,\vee) \) posee un \emph{primer elemento} (o elemento inferior) si existe \( 0 \in L \) tal que \( 0 \wedge a = 0 \) para todo \( a \).  
De modo dual, posee un \emph{último elemento} (o elemento superior) si existe \( 1 \in L \) tal que \( a \vee 1 = 1 \) para todo \( a \).  
Cuando existen ambos, se dice que el retículo es \emph{acotado}.
\end{definition}

\begin{definition}\label{def:reticulo_distributivo}
Un retículo \( \mathbb{L}=(L,\wedge,\vee) \) es \emph{distributivo} si para todo \( a,b,c \in L \) se cumple
\[
a \wedge (b \vee c) = (a \wedge b) \vee (a \wedge c).
\]
Esta condición es equivalente a su forma dual:
\[
a \vee (b \wedge c) = (a \vee b) \wedge (a \vee c).
\]
\end{definition}

\subsection{Filtros e ideales}

\begin{definition}\label{def:filtro}
Sea \( \mathbb{L}=(L,\wedge,\vee) \) un retículo.  
Un subconjunto no vacío \( F \subseteq L \) se llama \emph{filtro} si satisface:
\begin{enumerate}[\normalfont (1)]
  \item \( F \) es creciente con respecto al orden de \( \mathbb{L} \);
  \item si \( a,b \) pertenecen a \( F \), entonces \( a\wedge b \) también pertenece a \( F \)
        (es decir, \( F \) es cerrado bajo ínfimos).
\end{enumerate}
Un filtro \( F \) se dira \emph{propio} si $F \neq L$, es decir, está propiamente contenido en \( L \).
\end{definition}

Cuando \( \mathbb{L} \) es acotado con elemento superior \( 1 \),
la condición anterior equivale a exigir que \( 1 \) pertenezca a \( F \),
que \( F \) sea creciente y cerrado bajo ínfimos.

\begin{definition}
Un filtro \( P \) de \( \mathbb{L} \) se denomina \emph{filtro primo}
si es propio y cumple que para todos \( a,b \in L \),
si \( a\vee b \) pertenece a \( P \), entonces \( a \) o \( b \) pertenece a \( P \).
\end{definition}

Cuando \( \mathbb{L} \) es distributivo, denotaremos por \( \mathrm{Fi}(\mathbb{L}) \)
el conjunto de todos los filtros de \( \mathbb{L} \),
y por \( \mathcal{X}(\mathbb{L}) \) el conjunto de los filtros primos.

\begin{definition}
Sean \( \mathbb{L} \) un retículo distributivo acotado 
y \( X \) un subconjunto no vacío de \( L \). El \emph{filtro generado por \( X \)} es el menor filtro de \( \mathbb{L} \)
que contiene a todos los elementos de \( X \). Escribiremos \( \Fg^{\mathbb{L}}(X) \) para indicar al filtro generado por el conjunto $X$.
\end{definition}

El siguiente es un resultado conocido de la teoría de retículos distributivos acotados. Solo enunciaremos dicho resultado. Una prueba del mismo puede encontrarse en [{\cite[Lema~2.4.5]{SS}}].


\begin{lemma}\label{lem:filtro_generado}
Sea \( \mathbb{L} \) un retículo distributivo acotado.  Si \( X \) un subconjunto no vacío de \( L \) entonces: 
\[
\Fg^{\mathbb{L}}(X)
   = \{\, a \in L :
        \text{existen } x_{1},\dots,x_{n}\text{ en }X
        \text{ tales que } x_{1}\wedge\dots\wedge x_{n} \le a \,\}.
\]
\end{lemma}

Las nociones de filtro y filtro primo jugarán un rol central a lo largo de este trabajo. Sin embargo, pueden definirse sus nociones duales, las cuales también jugarán un papel muy importante.

\begin{definition}\label{def:ideal}
Sea \( \mathbb{L}=(L,\wedge,\vee) \) un retículo. Un subconjunto no vacío \( I \subseteq L \) se dirá \emph{ideal} si satisface las siguiente condiciones:
\begin{enumerate}[\normalfont (1)]
  \item \( I \) es decreciente respecto del orden de \( \mathbb{L} \),
  \item si \( a,b \) pertenecen a \( I \), entonces \( a\vee b \) pertenece a \( I \)
        (es decir, \( I \) es cerrado bajo supremos).
\end{enumerate}
Un ideal \( I \) se dirá \emph{propio} si $I \neq L$, es decir, está propiamente contenido en $L$.
\end{definition}

Si \( \mathbb{L} \) es un retículo acotado con elemento inferior \( 0 \),
entonces \( I \) es un ideal si y sólo si contiene a \( 0 \), es decreciente
y es cerrado bajo supremos.

\begin{definition}
Un ideal \(  I \) de \( \mathbb{L} \) se dirá \emph{ideal primo}
si es propio y para todo par \( a,b \in L \),
si \( a\wedge b \in I \) entonces  \( a \in I \) o \( b \in I \).
\end{definition}

Denotaremos por \( \mathrm{Ide}(\mathbb{L}) \) el conjunto de todos los ideales de \( \mathbb{L} \),
y por \( \mathcal{Y}(\mathbb{L}) \) el conjunto de todos los ideales primos.

\begin{definition} Sean \( \mathbb{L} \) un retículo distributivo acotado 
y  \( X \) un subconjunto no vacío de \( L \). El \emph{ideal generado por \( X \)} es el menor ideal de $\mathbb{L}$ que contiene a todos los elementos de \( X \) y lo indicaremos por \( \Ig^{\mathbb{L}}(X) \).
\end{definition}

\begin{lemma}\label{lem:ideal_generado}Sea \( \mathbb{L} \) un retículo distributivo acotado. Si \( X \) es un subconjunto no vacío de \( L \) entonces
\[
\Ig^{\mathbb{L}}(X)
   = \{\, a \in L :
        \text{existen } x_{1},\dots,x_{n}\text{ en }X
        \text{ tales que } a \le x_{1}\vee\dots\vee x_{n} \,\}.
\]
\end{lemma}

\begin{theorem}\label{thm:complemento-primo}
Sea $\mathbb L=(L,\wedge,\vee,0,1)$ un retículo distributivo acotado.  
Para todo $P\subseteq L$ se tiene:
\[
P\in\mathcal{X}(\mathbb{L})
\quad\Longleftrightarrow\quad
P^{c}\in \mathcal{Y}(\mathbb{L}),
\]
donde $P^{c}=L\setminus P$.
\end{theorem}

\begin{proof}[{\textbf{Demostración:}}]
Probaremos ambas implicaciones.
\begin{itemize}
\item $\Rightarrow$) Sea $P$ un filtro primo. Vamos a verificar que $P^{c}$ es un ideal primo. Comenzaremos mostrando que $P^{c}$ es decreciente. En efecto supongamos que $a\notin P$ y $b\le a$. Por absurdo, si $b\in P$ se sigue que $a\in P$, puesto que $P$ es creciente, lo que es una contradicción. Veamos ahora que que $P^{c}$ es cerrado por $\vee$. En efecto, supongamos, por absurdo, que $a,b\notin P$ y $a\vee b \in P$. Se sigue entonces que, al ser $P$ un filtro primo, $a\in P$ o $b\in P$, lo que es una contradicción. Luego, $a \vee b \in P^{c}$. Por último, veamos que $P^{c}$ es un ideal primo. Supongamos que $a\wedge b\in P^{c}$ y asumiremos que $a,b\in P$. Teniendo en cuenta que $P$ es filtro, se sigue entonces que $a\wedge b\in P$ por cierre de $P$ bajo $\wedge$, lo que resulta en una contradicción. Por tanto $a\notin P$ o $b\notin P$, es decir $P^{c}$ resulta un ideal primo

\smallskip
($\Leftarrow$)  
Si $I=P^{c}$ es un ideal primo, un razonamiento dual muestra que  
$I^{c}=P$ es creciente, cerrado bajo $\wedge$ y primo respecto de $\vee$.  
Por lo tanto $P$ es un filtro primo.
\end{itemize} 
\end{proof}

El siguiente resultado juega un papel central en la teoría de retículos distributivos acotados y, en particular, en este trabajo.

\begin{theorem}[{\cite[Teorema~2.4.13]{SS}}, Teorema del Filtro Primo] \label{thm:filtro_primo}
Sean \( \mathbb{L}=(L,\wedge,\vee,0,1) \) un retículo distributivo acotado, \( F \) un filtro e \( I \) un ideal tales que \( F\cap I=\emptyset \). Entonces existe un filtro primo \( P \) de \( \mathbb{L} \)
tal que \( F \subseteq P \) y \( P \cap I=\emptyset \).
\end{theorem}

\begin{definition}\label{def:homomorfismo_reticulos}
Sean \( \mathbb{L}_{1}=(L_{1},\wedge_{1},\vee_{1}) \) y
\( \mathbb{L}_{2}=(L_{2},\wedge_{2},\vee_{2}) \) retículos.
Una función \( f\colon L_{1}\to L_{2} \) se denomina
\emph{homomorfismo de retículos} si para todo par \( a,b\in L_{1} \) se verifican las siguiente condiciones:
\[
f(a\wedge_{1} b)=f(a)\wedge_{2} f(b), 
\qquad
f(a\vee_{1} b)=f(a)\vee_{2} f(b).
\]
\end{definition}


Es inmediato que la función identidad y la composición de homomorfismos de retículos
son nuevamente homomorfismos. Esta simple observación nos permite establecer la siguiente proposición.


\begin{proposition}
Los retículos distributivos acotados, junto con los homomorfismos de retículos, constituyen una categoría, a la que llamaremos $\BDL$.
La composición de flechas se define como la composición usual de funciones y la identidad de cada objeto está dada por el morfismo identidad \( \Id \).
\end{proposition}


\subsection{El retículo dual}

En esta subsección recordaremos la construcción del retículo dual de 
un retículo distributivo acotado, así como algunas consecuencias 
inmediatas que utilizaremos más adelante.


\begin{definition}\label{def:retículo-dual} Sea $\mathbb{L} = (L,\wedge,\vee,0,1)$ un retículo distributivo acotado. Se define su \emph{retículo dual} de la siguiente manera: 
\[
\mathbb{L}^{-1} = \left(L,\vee^{\mathbb{L}^{-1}},\wedge^{\mathbb{L}^{-1}},1^{\mathbb{L}^{-1}},0^{\mathbb{L}^{-1}}\right),
\]
donde, $a \vee^{\mathbf{L^{-1}}} b = a \we^{\mathbb{L}} b$, $a \we^{\mathbf{L^{-1}}} b = a \vee^{\mathbb{L}} b$, $0^{\mathbb{L}^{-1}} = 1^{\mathbb{L}}$ y $1^{\mathbb{L}^{-1}} = 0^{\mathbb{L}}$.
\end{definition}


\begin{remark}\label{rem:orden-dual}
A partir de la Definición~\ref{def:retículo-dual}, el orden de 
$\mathbb{L}^{-1}$ viene dado por la relación
\[
b \leq^{-1} a
\quad\Longleftrightarrow\quad
a \leq b \ \text{en}\ \mathbb{L}.
\]
Equivalente y más explícitamente:
\[
b \leq^{-1} a
\quad\Longleftrightarrow\quad
b = b \wedge^{\mathbb{L}^{-1}} a
\quad\Longleftrightarrow\quad
a = a \vee^{\mathbb{L}^{-1}} b.
\]
\end{remark}



\begin{remark} Toda afirmación válida para retículos distributivos acotados puede ser expresada dualmente en el retículo  dual \( \mathbb{L}^{-1} \).
En particular, los filtros de \( \mathbb{L} \) corresponden a los ideales de \( \mathbb{L}^{-1} \),
y viceversa.
\end{remark}

\begin{lemma}\label{lem:Li-dual}
Para todo retículo $\mathbb{L}$ y todo $i\in\{\pm1\}$ se tiene
\[
(\mathbb{L}^{-1})^{i}=\mathbb{L}^{-i}.
\]
\end{lemma}

\begin{proof}[{\textbf{Demostración:}}]
Si $i=1$, entonces claramente $(\mathbb{L}^{-1})^{1}=\mathbb{L}^{-1}$.

Si $i=-1$, al aplicar nuevamente $(-1)$ vemos que se revierten las operaciones de $\mathbb{L}^{-1}$ y recupera $\mathbb{L}$, por lo que $\left(\mathbb{L}^{-1}\right)^{-1}=\mathbb{L}$.
\end{proof}

\begin{definition}\label{def:i-ideal}
Sea $\mathbb{L}$ un retículo distributivo acotado y $i\in\{1,-1\}$.
Un subconjunto no vacío $I\subseteq L$ se llamará \emph{$i$-ideal} si cumple que:
\begin{enumerate}[\normalfont 1)]
\item Si $a\in I$ y $b\in L$ son tales que $b\le^i a$, entonces $b\in I$;
\item Si $a,b\in I$, entonces $a\vee^{\mathbb{L}^i} b\in I$.
\end{enumerate}
De este modo:
\begin{itemize}
  \item si $i=1$, los $i$-ideales coinciden con los ideales de $\mathbb{L}$;
  \item si $i=-1$, los $i$-ideales coinciden con los filtros de $\mathbb{L}$.
\end{itemize}
\end{definition}


\begin{lemma}\label{lem:primos-op-complemento}
Sea $\mathbb{L}$ un retículo distributivo acotado y $P\subseteq L$.
Entonces
\[
P \in \mathcal{X}(\mathbb{L}^{-1})
\quad\Longleftrightarrow\quad
P^c \in \mathcal{X}(\mathbb{L}).
\]
\end{lemma}

\begin{proof}[{\textbf{Demostración:}}]
\noindent(\(\Rightarrow\))  
Si \(P\in\mathcal{X}(\mathbb{L}^{-1})\), entonces \(P\) es un filtro primo en \(\mathbb{L}^{-1}\). Teniendo en cuenta la definición de retículo dual, se sigue que $P$ es un ideal primo en \(\mathbb{L}\), por tanto \(P \in \mathcal{Y}(\mathbb{L})\).
Por el Teorema~\ref{thm:complemento-primo}, el complemento de un ideal primo
es un filtro primo; en consecuencia, \(P^{c}\in\mathcal{X}(\mathbb L)\).

\medskip
\noindent(\(\Leftarrow\))  
Podemos aplicar un razonamiento análogo, recordando que $(P^c)^c=P$.
\end{proof}



\section{Dualidad de Priestley}

La dualidad de Priestley constituye el caso base sobre el cual se extenderán las dualidades consideradas en capítulos posteriores.  
Describe una correspondencia contravariante entre la categoría algebraica de los retículos distributivos acotados y una categoría de espacios topológicos parcialmente ordenados.  
Esta versión utiliza el espectro de \emph{filtros primos}, en lugar del de ideales primos introducido originalmente por H.~A.~Priestley en~\cite{P}.  
Para un desarrollo clásico y referencias adicionales remitimos a~\cite{DP,CLP,SS}.

\paragraph{Espacios de Priestley}

\begin{definition}\label{def:espacio_priestley2} Un \emph{espacio de Priestley} es una terna \( (X,\leq,\tau) \) tal que:
\begin{enumerate}[\normalfont (1)]
  \item \( (X,\tau) \) es un espacio topológico compacto;
  \item \( (X,\leq) \) es un conjunto parcialmente ordenado;
  \item para todo par \( x,y\in X \) con \( x\nleq y \), existe un conjunto \( U \) que es clopen y creciente tal que \( x\in U \) y \( y\notin U \).
\end{enumerate}
\end{definition}

\begin{remark}
La propiedad (3) de la Definicion \ref{def:espacio_priestley2} nos indica que el espacio topológico es totalmente disconexo en el orden. Dicha propiedad es conocida como axioma de separación de Priestley.
\end{remark}

\begin{definition}\label{def:familias_Priestley}
Sea \( (X,\leq,\tau) \) un espacio de Priestley. Definimos:
\begin{align*}
\mathcal{D}(X) &= \{\, U\subseteq X : U \text{ es clopen y creciente } \},\\
\Delta(X)      &= \{\, V\subseteq X : V \text{ es clopen y decreciente } \},\\
\mathrm{OpUp}(X) &= \{\, O\in\tau : O \text{ es creciente } \},\\
\mathrm{ClUp}(X) &= \{\, C\subseteq X : C \text{ es cerrado y creciente } \}.
\end{align*}
\end{definition}

El siguiente resultado (ver  [{\cite[Lema~6.5.1]{SS}}]) sobre espacios de Priestley nos permite dar una buena caracterización de abiertos y cerrados del espacio, por medio de uniones e intersecciones de elementos de la base.

\begin{lemma}\label{lem:Priestley_separacion2}
Sea \( (X,\leq,\tau) \) un espacio de Priestley. Entonces \( (X,\leq,\tau) \) es  un espacio Hausdorff. Más aún, si \( C\subseteq X \) es cerrado y creciente, y \( x\notin C \), entonces existe \( U\in \mathcal{D}(X) \)
tal que \( C\subseteq U \) y \( x\notin U \).
\end{lemma}

\begin{corollary}\label{cor:Priestley_decomp}
En todo espacio de Priestley \( (X,\leq,\tau) \) se cumple:
\begin{enumerate}[\normalfont (1)]
  \item cada cerrado creciente es intersección (posiblemente infinita) de elementos de \( \mathcal{D}(X) \);
  \item cada abierto creciente es unión (posiblemente infinita) de elementos de \( \mathcal{D}(X) \).
\end{enumerate}
\end{corollary}

\begin{definition}\label{def:morfismos_Priestley}
Sean \( (X_{1},\leq_{1},\tau_{1}) \) y \( (X_{2},\leq_{2},\tau_{2}) \) espacios de Priestley.
Una función \( f\colon X_{1}\to X_{2} \) es un \emph{morfismo de espacios de Priestley}
si se cumple que:
\begin{enumerate}
    \item $f$ es continua.
    \item $f$ es monótona.
\end{enumerate}
\end{definition}

\begin{proposition}
Los espacios de Priestley, junto con las funciones continuas y monótonas,
constituyen una categoría, a la que llamaremos \(\mathsf{PS}\).
La composición de flechas se define como la composición usual de funciones
y la identidad de cada objeto está dada por el morfismo identidad \(\Id\).
\end{proposition}



\paragraph{Función $\sigma_\mathbb{L}$ de Stone}\;

\medskip

Ahora recordaremos la noción clásica de función de Stone, que asocia a cada elemento de un retículo distributivo acotado el conjunto de filtros 
primos que lo contienen. 

\begin{definition}\label{def:Stone-func}
Sea \( \mathbb{L}=(L,\wedge,\vee,0,1) \) un retículo distributivo acotado. La \emph{función de Stone} asociada a \( \mathbb{L} \) es una aplicación $\sigma_{\mathbb{L}}\colon L\to\mathcal{P}(\mathcal{X}(\mathbb{L}))$, definida de la siguiente manera:
\[\sigma_{\mathbb{L}}(a)=\{P\in \mathcal{X}(\mathbb{L}) : a\in P\}, \]
donde, recordemos, \( \mathcal{X}(\mathbb{L}) \) es el conjunto de filtros primos de \( \mathbb{L} \).
\end{definition}

Recordemos que, para cada conjunto $X \neq \emptyset$ el álgebra $\left(\mathcal{P}(X),\cap,\cup,\emptyset,X\right)$ es un retículo distributivo acotado. En particular, si $\mathbb{L}$ es un retículo distributivo acotado, entonces
\[ \left( \mathcal{P}(\mathcal{X}(\mathbb{L})),\cap,\cup,\emptyset,\mathcal{X}(\mathbb{L}) \right), \]
es un retículo distributivo acotado y $\sigma_{\mathbb{L}} \colon \mathbb{L} \to \mathcal{P}(\mathcal{X}(\mathbb{L}))$ resulta una aplicación entre ambos retículos. El resultado que sigue mostrará que, en efecto, resulta ser un homomorfismo de retículos distributivos acotados.

\begin{lemma}[{\cite[Lema~2.5.10]{SS}}]\label{lem:Stone_prop}
Para todo \( a,b\in L \):
\[
\sigma_{\mathbb{L}}(0)=\emptyset,\quad
\sigma_{\mathbb{L}}(1)=\mathcal{X}(\mathbb{L}),\quad
\sigma_{\mathbb{L}}(a\wedge b)=\sigma_{\mathbb{L}}(a)\cap\sigma_{\mathbb{L}}(b),\quad
\sigma_{\mathbb{L}}(a\vee b)=\sigma_{\mathbb{L}}(a)\cup\sigma_{\mathbb{L}}(b).
\]
\end{lemma}

\begin{corollary} Todo retículo distributivo $\mathbb{L}$ es isomorfo a una subálgebra de $\mathcal{P}(\mathcal{X}(\mathbb{L}))$ por medio de la aplicación de Stone $\sigma_{\mathbb{L}}$.
\end{corollary}

\paragraph{El funtor \(\X\colon \mathsf{BDL}\to\mathsf{PS}\)}\;

\medskip

En este apartado introducimos el funtor que asocia a cada retículo distributivo acotado un espacio de Priestley. Su propósito es traducir información algebraica en información topológica y ordenada, de modo de preparar el terreno para las dualidades que aparecerán más adelante. 

Comenzaremos describiendo la acción de dicho funtor sobre los objetos, luego su acción sobre las flechas, y finalmente comprobaremos que las asignaciones obtenidas constituyen efectivamente un funtor contravariante.


Sea $\mathbb{L} = (L,\vee,\wedge,0,1)$ un retículo distributivo acotado. Definimos la estructura
\[
\X(\mathbb{L}) = (\mathcal{X}(\mathbb{L}),\tau_{\mathbb{L}},\subseteq_{\mathbb{L}}),
\]
donde:
\begin{itemize}
  \item \( \subseteq_{\mathbb{L}} \) es la relación de orden dada por la inclusión sobre el conjunto de los filtros primos;
  \item la topología \( \tau_{\mathbb{L}} \) está generada por la subbase
        \(
        \mathscr{S}_{\mathbb{L}} = \{ \sigma_{\mathbb{L}}(a) \colon a \in \mathbb{L} \} \cup \{ \sigma_{\mathbb{L}}(a)^{c} \colon a \in \mathbb{L} \}.
        \)
\end{itemize}
Teniendo en cuenta la definición de subbase, es sencillo comprobar que los conjuntos de la forma \( \sigma_{\mathbb{L}}(a)\setminus\sigma_{\mathbb{L}}(b) \)
con \( a,b\in L \) forman una base de abiertos para \( \tau_{\mathbb{L}} \).

\begin{theorem}\label{thm:X_bien_definido}
Para todo retículo distributivo acotado \( \mathbb{L} \),
la terna
\[
\X(\mathbb{L}) = (\mathcal{X}(\mathbb{L}),\tau_{\mathbb{L}},\subseteq_{\mathbb{L}}),
\]
es un espacio de Priestley.
\end{theorem}

\begin{proof}[{\textbf{Demostración:}}]
Sea $\mathbb{L}$ un retículo distributivo acotado. Vamos a comprobar que \(\X(\mathbb{L})\) satisface los tres requisitos de la Definición~\ref{def:espacio_priestley2}.

\begin{enumerate}[\normalfont (1)]
\item Veamos que $(\mathcal{X}(\mathbb{L}),\tau_{\mathbb{L}})$ es compacto.  
Sean $\{U_i\}_{i\in I}$ abiertos que cubren $\mathcal{X}(\mathbb{L})$ y
supongamos, hacia un absurdo, que no existe un subcubrimiento finito.

Cada $U_i$ es unión de conjuntos básicos de la forma
\[
\sigma_{\mathbb L}(a_1)\cap\cdots\cap\sigma_{\mathbb L}(a_m)
\cap\sigma_{\mathbb L}(b_1)^c\cap\cdots\cap\sigma_{\mathbb L}(b_n)^c.
\]
Si un conjunto de este tipo contiene a un filtro primo $P$, entonces
\[
a_1,\dots,a_m\in P 
\qquad\text{y}\qquad 
b_1,\dots,b_n\notin P.
\]

Como no hay subcubrimiento finito, para cada elección finita de índices
$F\subseteq I$ existe un filtro primo $P_F$ que no pertenece a ninguno de los
$U_i$ con $i\in F$.  
Para cada tal $P_F$, definamos
\[
A_F=\{\,a\in L : a\in P_F\,\}, 
\qquad 
B_F=\{\,b\in L : b\notin P_F\,\}.
\]

El punto clave es que, para todo $i\in I$, existe algún filtro primo $P_F$
que satisface simultáneamente todas las condiciones
\[
a\notin P_F\quad\text{y}\quad b\in P_F
\]
que definen los conjuntos básicos incluidos en $U_i$.  
En otras palabras, para cada $i$ podemos seleccionar un filtro $P_F$ que
evita explícitamente todos los requisitos que harían que $P_F$ cayera en $U_i$.

Consideremos ahora el conjunto
\[
F=\bigcup_{F\subseteq I\ \text{finito}} A_F.
\]
Este conjunto es no vacío y cerrado hacia arriba, y siempre que $a,b\in F$,
alguna de las elecciones $P_F$ contiene ambos, por lo que $a\wedge b\in F$.
Por tanto, $F$ genera un filtro propio de $\mathbb L$.

Por el Teorema del filtro primo,
existe un filtro primo $P$ que contiene a $F$.  
Como cada $P_F$ evita todos los $U_i$, el filtro $P$ también evita todos los
$U_i$.  
Esto contradice la suposición de que los $U_i$ cubren $\mathcal{X}(\mathbb L)$.

Por lo tanto, existe un subcubrimiento finito y el espacio es compacto.


\item Veamos que \( (\mathcal{X}(\mathbb{L}),\subseteq_{\mathbb{L}}) \) es un conjunto parcialmente ordenado. La relación $\subseteq_{\mathbb L}$ es la inclusión entre filtros primos, y por ser una inclusión, es reflexiva, antisimétrica y transitiva. Luego $(\mathcal{X}(\mathbb L),\subseteq_{\mathbb L})$ es un conjunto parcialmente ordenado.

\item Veamos que se cumple el axioma de separación de Priestley. Sea $P,Q\in\mathcal{X}(\mathbb L)$ con $P\nsubseteq Q$.  
Entonces existe $a\in P\setminus Q$. En particular:
\[
P\in\sigma_{\mathbb L}(a),\qquad Q\in\sigma_{\mathbb L}(a)^c.
\]

Observemos que:

- $\sigma_{\mathbb L}(a)$ es creciente: si $P\in\sigma_{\mathbb L}(a)$ y $P\subseteq R$, entonces $a\in R$ y por tanto $R\in\sigma_{\mathbb L}(a)$;

- $\sigma_{\mathbb L}(a)$ es clopen porque pertenece a la subbase generadora y su complemento también pertenece a ella.

Así, el conjunto clopen creciente $U=\sigma_{\mathbb L}(a)$ separa $P$ y $Q$ como exige la Definición~\ref{def:espacio_priestley2}.

\end{enumerate}

\medskip

Hemos verificado las tres condiciones, por lo que $\X(\mathbb L)$ es un espacio de Priestley.
\end{proof}
\medskip



A continuación vamos a describir la acción del funtor sobre las flechas.

Sea \( f\colon \mathbb{L}_{1}\to\mathbb{L}_{2} \) una flecha en la categoría \( \mathsf{BDL} \). Definimos la aplicación $\X(f)\colon \X(\mathbb{L}_{2}) \to \X(\mathbb{L}_{1})$ de la siguiente manera:
\[ \X(f)(P)=f^{-1}[P]. \]
En el siguiente resultado, comprobaremos que efectivamente, $\X(f)$ es una flecha en la categoría $\PS$.

\begin{lemma}\label{lem:X-cont-mon} 
Sea \( f\colon \mathbb{L}_{1}\to\mathbb{L}_{2} \) una flecha en la categoría $\mathsf{BDL}$. Entonces
la aplicación \( \X(f) \) es una flecha en $\PS$.
\end{lemma}

\begin{proof}[{\textbf{Demostración:}}] 
Veamos que $\X(f)$ es continua y preserva el orden. En efecto, para comprobar la continuidad tomemos un conjunto $U$ abierto básico de $\X(\mathbb{L}_{1})$, es decir, $U = \sigma_{\mathbb{L}_{1}}(a)$, para algún  \( a\in L_{1} \). Entonces
\[
\X(f)^{-1}[\sigma_{\mathbb{L}_{1}}(a)]
  = \{P\in X(\mathbb{L}_{2}) : a\in f^{-1}[P]\}
  = \{P : f(a)\in P\}
  = \sigma_{\mathbb{L}_{2}}(f(a)),
\]
que es un clopen creciente; luego \( \X(f) \) es continua.
La monotonía es inmediata porque la preimagen preserva inclusiones.
\end{proof}

\begin{theorem}\label{thm:X-functor}
Las asignaciones previamente descriptas \( \mathbb{L}\mapsto\X(\mathbb{L}) \) y \( f\mapsto\X(f) \)
definen un funtor contravariante
\[
\X\colon \mathsf{BDL}\to\mathsf{PS}.
\]
\end{theorem}

\begin{proof}[{\textbf{Demostración:}}]
Recordemos que, por Definición~\ref{def: funtor}, para verificar que una asignación
define un funtor contravariante debemos comprobar dos condiciones:
\begin{enumerate}
    \item \emph{Preservación de identidades.}  
Sea \( \mathbb{L} \) un objeto de \( \mathsf{BDL} \).  
Entonces
\[
\X(\Id_{\mathbb{L}})(P)=\Id_{\mathbb{L}}^{-1}[P]=P,
\]
para todo \( P\in\mathcal{X}(\mathbb{L}) \), lo cual muestra que  
\( \X(\Id_{\mathbb{L}})=\Id_{\X(\mathbb{L})} \).
    \item \emph{Preservación de la composición.}  
Sean
\[
\mathbb{L}_{1}\xrightarrow{f}\mathbb{L}_{2}\xrightarrow{g}\mathbb{L}_{3}
\]
morfismos en \( \mathsf{BDL} \). Para cada \( P\in\mathcal{X}(\mathbb{L}_{3}) \),
\[
\X(g\circ f)(P)
=(g\circ f)^{-1}[P]
=f^{-1}\bigl[g^{-1}[P]\bigr]
=\X(f)\bigl(\X(g)(P)\bigr).
\]
Por lo tanto,
\[
\X(g\circ f)=\X(f)\circ\X(g),
\]
lo cual verifica la regla de composición propia de los funtores contravariantes.
\end{enumerate}
Con ambas propiedades satisfechas, concluimos que
\(
\X\colon\mathsf{BDL}\to\mathsf{PS}
\)
es un funtor contravariante.
\end{proof}
En conclusión, el funtor $\X$ asigna a cada retículo distributivo acotado una estructura topológico-ordenada que captura de manera adecuada a su comportamiento algebraico, y envía homomorfismos de retículos a aplicaciones continuas y monótonas.


\paragraph{El funtor \(\D\colon \mathsf{PS}\to\mathsf{BDL}\)}\;

\medskip

En este apartado introducimos el funtor que asocia a cada espacio de Priestley un retículo distributivo acotado formado por sus clopens crecientes. 
Comenzaremos describiendo su acción sobre los objetos, luego su acción sobre las flechas, 
y finalmente demostraremos que dichas asignaciones definen efectivamente un funtor 
contravariante entre las categorías involucradas.

Sea \( (X,\leq,\tau) \) un espacio de Priestley.
Definimos
\[
\D(X)=(\mathcal{D}(X),\cap,\cup,\emptyset,X),
\]
donde \( \mathcal{D}(X) \) es el conjunto de los clopen crecientes de \(X\).

Veamos que esta estructura es realmente un retículo 
distributivo acotado.

\begin{lemma}\label{lem:D-obj-reticulo}
Para todo espacio de Priestley \( (X,\leq,\tau) \),
la estructura 
\[
\D(X)=(\mathcal{D}(X),\cap,\cup,\emptyset,X)
\]
es un retículo distributivo acotado.
\end{lemma}

\begin{proof}[{\textbf{Demostración:}}]
Mostraremos sucesivamente que \(\mathcal{D}(X)\) es cerrado bajo las operaciones 
propuestas, que estas definen un retículo y que dicho retículo es distributivo y acotado.

\smallskip
\noindent\emph{Cerradura.}
Sean \(U,V\in \mathcal{D}(X)\). Como ambos son clopen y crecientes, tanto su intersección como su unión finita \(U\cap V\) es abierta, cerrada y creciente.
Por lo tanto, \(U\cap V,\,U\cup V\in \mathcal{D}(X)\).

\smallskip
\noindent\emph{Estructura de retículo.}
Las operaciones \(\cap\) y \(\cup\) satisfacen las leyes de conmutatividad, 
asociatividad y absorción porque son las operaciones usuales de conjuntos. 
Por cierre, dichas operaciones permanecen dentro de \(\mathcal{D}(X)\), por lo que 
\((\mathcal{D}(X),\cap,\cup)\) es un retículo.

\smallskip
\noindent\emph{Distributividad.}
Como la intersección y la unión de conjuntos satisfacen las leyes distributivas,
el retículo resultante es distributivo.

\smallskip
\noindent\emph{Acotado.}
El conjunto vacío \(\emptyset\) y el total \(X\) son ambos clopen y crecientes;
actúan como los elementos mínimo y máximo del retículo respectivamente. Por lo tanto, \(\D(X)\) es un 
retículo distributivo acotado.
\end{proof}

Veamos ahora su acción sobre morfismos. Sea $g\colon (X_{1},\leq_{1},\tau_{1})\to(X_{2},\leq_{2},\tau_{2})$ un morfismo en la categoría \( \mathsf{PS} \). Definimos
\[
\D(g)\colon \D(X_{2})\to\D(X_{1}),
\qquad
\D(g)(U)=g^{-1}[U].
\]


\begin{lemma}\label{lem:D-morfismo-bdl}
Sea 
\( g\colon (X_{1},\leq_{1},\tau_{1})\to(X_{2},\leq_{2},\tau_{2}) \)
un morfismo en \( \mathsf{PS} \).
Entonces la aplicación
\[
\D(g)\colon \D(X_{2})\to\D(X_{1}),
\qquad
\D(g)(U)=g^{-1}[U].
\]
es un homomorfismo de retículos distributivos acotados.
\end{lemma}

\begin{proof}[{\textbf{Demostración:}}]
Debemos probar dos cosas: (i) que la preimagen de un clopen creciente sigue 
siendo un clopen creciente, y (ii) que la preimagen preserva las operaciones del retículo.

\smallskip
\noindent\emph{(i) La preimagen de un clopen creciente es clopen creciente.}
Sea \(U\in \mathcal{D}(X_{2})\). Entonces \(U\) es clopen y creciente.

\begin{itemize}
    \item Como \(g\) es continua, \(g^{-1}[U]\) es abierto; y como \(U\) es cerrado, 
          también \(g^{-1}[U]\) lo es.
    \item Sea \(x,y\in X_{1}\) tales que \(x\leq_{1} y\) y \(x\in g^{-1}[U]\).  
          Entonces \(g(x)\leq_{2}g(y)\) (monotonía de \(g\)) y \(g(x)\in U\).  
          Como \(U\) es creciente, \(g(y)\in U\), por lo que \(y\in g^{-1}[U]\).
\end{itemize}
Así, \(g^{-1}[U]\in \mathcal{D}(X_{1})\).

\smallskip
\noindent\emph{(ii) Preservación de las operaciones.}
Para cualesquiera \(U,V\in \mathcal{D}(X_{2})\),
\[
g^{-1}[U\cap V]=g^{-1}[U]\cap g^{-1}[V],
\qquad
g^{-1}[U\cup V]=g^{-1}[U]\cup g^{-1}[V].
\]
Además,
\[
g^{-1}[X_{2}]=X_{1},
\qquad
g^{-1}[\emptyset]=\emptyset.
\]
Todas estas igualdades se siguen directamente de las propiedades básicas de la 
preimagen. Por lo tanto, \(\D(g)\) es un homomorfismo de retículos distributivos acotados.
\end{proof}
Veamos por último que estas asignaciones definen un funtor:
\begin{theorem}
Las asignaciones $X\longmapsto\D(X)$ y $g\longmapsto\D(g)$ definen un funtor contravariante
\[
\D\colon\mathsf{PS}\to\mathsf{BDL}.
\]
\end{theorem}

\begin{proof}[{\textbf{Demostración:}}]
Sea \(X\) un objeto de \(\mathsf{PS}\).  
La preimagen respecto de la identidad satisface
\[
\D(\Id_{X})(U)=\Id_{X}^{-1}[U]=U,
\]
es decir \(\D(\Id_{X})=\Id_{\D(X)}\).

Si 
\(
X_{1}\xrightarrow{g}X_{2}\xrightarrow{h}X_{3}
\),
entonces para todo \(U\in\D(X_{3})\),
\[
\D(h\circ g)(U)
=(h\circ g)^{-1}[U]
=g^{-1}[h^{-1}[U]]
=\D(g)(\D(h)(U)),
\]
por lo que \(\D(h\circ g)=\D(g)\circ \D(h)\).  
Esto muestra la compatibilidad con la composición.
\end{proof}

\smallskip

Con esto queda completamente especificado el funtor contravariante 
\(\D\colon\mathsf{PS}\to\mathsf{BDL}\), que asigna a cada espacio de Priestley el 
retículo distributivo acotado de sus clopens crecientes y a cada flecha 
monótona y continua su correspondiente preimagen.



\paragraph{Las transformaciones naturales y la dualidad}\;

\medskip

En esta parte presentamos las aplicaciones que permiten comparar, para cada retículo distributivo acotado, la estructura original con la reconstruida mediante la doble aplicación de los funtores construidos en las secciones anteriores. Del mismo modo, para cada espacio de Priestley introduciremos una aplicación que permite recuperar el punto original a partir de los clopens crecientes que lo contienen. Estas construcciones serán esenciales para demostrar la dualidad completa entre ambas categorías.

\medskip

Recordemos primero que, para cada retículo distributivo acotado $\mathbb{L}$, la aplicación de Stone
\[
\sigma_{\mathbb{L}}\colon L\longrightarrow \mathcal{P}(\mathcal{X}(\mathbb{L})),
\qquad 
\sigma_{\mathbb{L}}(a)=\{\,P\in\mathcal{X}(\mathbb{L}) : a\in P\,\},
\]
es un homomorfismo de retículos. El siguiente resultado afirma que esta aplicación describe completamente el retículo cuando se la interpreta dentro de la estructura $\D(\X(\mathbb{L}))$.

\begin{theorem}\label{thm:sigma_iso}
Para cada retículo distributivo acotado $\mathbb{L}$,
la aplicación de Stone 
\[
\sigma_{\mathbb{L}}\colon \mathbb{L}\longrightarrow \D(\X(\mathbb{L}))
\]
es un isomorfismo de retículos distributivos acotados.
\end{theorem}

\begin{proof}[{\textbf{Demostración:}}]
Demostraremos que $\sigma_{\mathbb{L}}$ es un homomorfismo inyectivo y sobreyectivo.

\smallskip
\noindent\emph{(i) $\sigma_{\mathbb{L}}$ es un homomorfismo.}  
El Lema~\ref{lem:Stone_prop} establece que
\[
\sigma_{\mathbb{L}}(a\wedge b)=\sigma_{\mathbb{L}}(a)\cap\sigma_{\mathbb{L}}(b),
\qquad
\sigma_{\mathbb{L}}(a\vee b)=\sigma_{\mathbb{L}}(a)\cup\sigma_{\mathbb{L}}(b),
\]
y que $\sigma_{\mathbb{L}}(0)=\emptyset$, $\sigma_{\mathbb{L}}(1)=\mathcal{X}(\mathbb{L})$.
Esto prueba que $\sigma_{\mathbb{L}}$ preserva todas las operaciones del retículo.

\smallskip
\noindent\emph{(ii) Inyectividad.}  
Sean $a,b\in L$ tales que $\sigma_{\mathbb{L}}(a)=\sigma_{\mathbb{L}}(b)$.
Suponiendo $a\nleq b$, el Teorema del Filtro Primo
(Teorema~\ref{thm:filtro_primo}) nos garantiza la existencia de
$P\in\mathcal{X}(\mathbb{L})$ con $a\in P$ y $b\notin P$,
lo cual implica
\[
P\in\sigma_{\mathbb{L}}(a)\setminus\sigma_{\mathbb{L}}(b),
\]
contradicción.
Por lo tanto $a\leq b$. Un argumento simétrico prueba $b\leq a$, con lo cual $a=b$.

\smallskip

\noindent\emph{(iii) Sobreyectividad.}  
Sea $U\in\mathcal{D}(\mathcal{X}(\mathbb{L}))$.
Como $U$ y $U^{c}$ son cerrados de un espacio compacto, ambos son compactos.

Para cada $P\in U$ y cada $Q\in U^{c}$, dado que $U$ es creciente tenemos
$P\nsubseteq Q$. Por el Teorema del Filtro Primo, existe $a_{PQ}\in L$ tal que
\[
a_{PQ}\in P
\qquad\text{y}\qquad
a_{PQ}\notin Q.
\]
Es decir,
\[
P\in\sigma_{\mathbb L}(a_{PQ}),
\qquad
Q\in\sigma_{\mathbb L}(a_{PQ})^{c}.
\]

Fijado $P\in U$, la familia
$\{\sigma_{\mathbb L}(a_{PQ})^{c}:Q\in U^{c}\}$
cubre $U^{c}$.  
Como $U^{c}$ es compacto, existen $Q_{1},\dots,Q_{n}\in U^{c}$ tales que
\[
U^{c}\subseteq \sigma_{\mathbb L}(a_{PQ_{1}})^{c}\cup\cdots\cup
               \sigma_{\mathbb L}(a_{PQ_{n}})^{c}.
\]
Equivalentemente,
\[
U^{c}\subseteq 
\sigma_{\mathbb L}(a_{P})^{c},
\qquad
\text{donde } 
a_{P}=a_{PQ_{1}}\wedge\cdots\wedge a_{PQ_{n}}.
\]
Por tanto,
\[
\sigma_{\mathbb L}(a_{P})\subseteq U,
\qquad\text{y además } P\in\sigma_{\mathbb L}(a_{P}).
\]
Los conjuntos $\{\sigma_{\mathbb L}(a_{P}):P\in U\}$ cubren $U$.  
Como $U$ es compacto, existen $P_{1},\dots,P_{m}\in U$ tales que
\[
U=\sigma_{\mathbb L}(a_{P_{1}})\cup\cdots\cup\sigma_{\mathbb L}(a_{P_{m}}).
\]
Sea ahora
\[
a=a_{P_{1}}\vee\cdots\vee a_{P_{m}}.
\]
Entonces
\[
\sigma_{\mathbb L}(a)
=\sigma_{\mathbb L}(a_{P_{1}})\cup\cdots\cup\sigma_{\mathbb L}(a_{P_{m}})
=U.
\]
Esto prueba que todo $U\in\mathcal{D}(\mathcal{X}(\mathbb{L}))$ es de la forma $\sigma_{\mathbb L}(a)$ para algún $a\in L$, es decir, que 
$\sigma_{\mathbb L}$ es sobreyectiva.


\smallskip

Concluimos que $\sigma_{\mathbb{L}}$ es un isomorfismo de retículos distributivos acotados.
\end{proof}

\medskip

Pasemos ahora a la construcción dual.  
Sea $(X,\leq,\tau)$ un espacio de Priestley. Queremos reconstruir cada punto del espacio
a partir del conjunto de todos los clopens crecientes que lo contienen. Esto se logra mediante la siguiente aplicación.

\[
\epsilon_{X}\colon X\longrightarrow \X(\D(X)),
\qquad
\epsilon_{X}(x)=\{\,U\in \mathcal{D}(X): x\in U\,\}.
\]

\begin{theorem}\label{thm:epsilon_iso}
Para todo espacio de Priestley $(X,\leq,\tau)$,  
la aplicación $\epsilon_{X}\colon X\to \X(\D(X))$ es un isomorfismo en la categoría $\mathsf{PS}$.
\end{theorem}

\begin{proof}[{\textbf{Demostración:}}]
Notemos que como $(X,\leq,\tau)$ un espacio de Priestley, por el Lema \ref{lem:Priestley_separacion2}, este espacio es Hausdorff, por lo que basta con probar que es biyectiva y continua para probar que es un homeomorfismo. Además, probaremos monotonía para terminar de ver que es un morfismo en $\PS$.

\smallskip
\noindent\emph{(i) Monotonía.}  
Sean $x\leq y$.  
Si $U$ es un clopen creciente que contiene a $x$, entonces, por definición de conjunto creciente, también contiene a $y$.  
Esto implica $\epsilon_{X}(x)\subseteq\epsilon_{X}(y)$.

\smallskip
\noindent\emph{(ii) Continuidad.}  
Sea $U\in \mathcal{D}(X)$.  
Los abiertos básicos de $\X(\D(X))$ son los de la forma $\sigma_{\D(X)}(U)$.  
Calculemos su preimagen:
\[
\epsilon_{X}^{-1}[\sigma_{\D(X)}(U)]
= \{x\in X : U\in\epsilon_{X}(x)\}
= \{x\in X : x\in U\}
= U,
\]
que es un clopen creciente de $X$.  
Luego $\epsilon_{X}$ es continua.

\smallskip
\noindent\emph{(iii) Inyectividad.}  
Sean $x\neq y$.  
Por el axioma de separación de Priestley, existe un clopen creciente $U$ con $x\in U$ y $y\notin U$.  
Entonces $U\in\epsilon_{X}(x)$ pero $U\notin\epsilon_{X}(y)$, así que $\epsilon_{X}(x)\neq\epsilon_{X}(y)$.

\smallskip

\noindent\emph{(iv) Sobreyectividad.}  
Queremos ver que $\epsilon_{X}$ es sobreyectiva.
Como $\epsilon_{X}$ es continua (por (\emph{(ii)})) y $X$ es compacto, el conjunto
$\epsilon_{X}(X)$ es compacto en $\X(\D(X))$, y en particular es cerrado.

Supongamos, hacia un absurdo, que $\epsilon_{X}$ no es sobreyectiva.  
Entonces existe $Q\in\mathcal{X}(\mathcal{D}(X))$ tal que $Q\notin\epsilon_{X}(X)$.
Puesto que $\epsilon_{X}(X)$ es cerrado y $Q\notin\epsilon_{X}(X)$, existe,
por el axioma de separación de Priestley aplicado en $\X(\D(X))$, un clopen
creciente $V$ tal que
\[
Q\in V
\qquad\text{y}\qquad
V\cap\epsilon_{X}(X)=\varnothing.
\]

Podemos suponer sin pérdida de generalidad que $V$ es un elemento básico,
es decir, que existen $U,W\in\mathcal{D}(X)$ tales que
\[
V=\sigma_{\D(X)}(U)\cap\sigma_{\D(X)}(W)^{c}.
\]
En tal caso,
\[
\epsilon_{X}^{-1}[V]
=\epsilon_{X}^{-1}[\sigma_{\D(X)}(U)]
  \cap \epsilon_{X}^{-1}[\sigma_{\D(X)}(W)^{c}]
=U\cap W^{c}.
\]
Pero la condición $V\cap\epsilon_{X}(X)=\varnothing$ implica
\[
\epsilon_{X}^{-1}[V]=\varnothing,
\qquad\text{es decir,}\qquad
U\cap W^{c}=\varnothing.
\]
Luego $U\subseteq W$ y por tanto
\[
\sigma_{\D(X)}(U)\subseteq\sigma_{\D(X)}(W).
\]

Como $Q\in\sigma_{\D(X)}(U)$, se tiene $U\in Q$.
Y como $\sigma_{\D(X)}(U)\subseteq\sigma_{\D(X)}(W)$, también
$W\in Q$.  
Pero $Q\in V$ implica $Q\notin\sigma_{\D(X)}(W)$, es decir,
$W\notin Q$.
Esto es una contradicción. Por lo tanto $\epsilon_{X}$ es sobreyectiva.


\noindent Con esto se prueba que $\epsilon_{X}$ es un isomorfismo en $\mathsf{PS}$.
\end{proof}

\medskip

\begin{theorem}[Condición de naturalidad]\label{thm:naturality} Sean 
\( f\colon \mathbb{L}_{1}\to\mathbb{L}_{2} \) un morfismo en \( \mathsf{BDL} \)
y 
\( g\colon X_{1}\to X_{2} \) un morfismo en \( \mathsf{PS} \).
Entonces los siguientes diagramas conmutan:
\end{theorem}

\begin{proof}[{\textbf{Demostración:}}]
Demostraremos la conmutatividad de cada diagrama por separado.

\begin{enumerate}
    \item Naturalidad de \(\sigma\).
    
Sea \( a\in L_{1} \) y sea \( P\in\mathcal{X}(\mathbb{L}_{2}) \).
La rama izquierda del cuadrado envía \( a \) a \( f(a) \), y luego:
\[
\sigma_{\mathbb{L}_{2}}(f(a))
=\{\,P\in\mathcal{X}(\mathbb{L}_{2}) : f(a)\in P\,\}.
\]

La rama superior–derecha envía primero
\[
a\mapsto \sigma_{\mathbb{L}_{1}}(a)=\{\,Q\in\mathcal{X}(\mathbb{L}_{1}) : a\in Q\,\},
\]
y luego aplica \( \D(\X(f)) \), cuya acción es la preimagen:
\[
\D(\X(f))(\sigma_{\mathbb{L}_{1}}(a))
= \{\,P\in\mathcal{X}(\mathbb{L}_{2}) : f^{-1}[P]\in\sigma_{\mathbb{L}_{1}}(a)\,\}.
\]

Pero
\[
f^{-1}[P]\in\sigma_{\mathbb{L}_{1}}(a)
\;\Longleftrightarrow\;
a\in f^{-1}[P]
\;\Longleftrightarrow\;
f(a)\in P.
\]

Por lo tanto:
\[
\D(\X(f))(\sigma_{\mathbb{L}_{1}}(a))
= \{\,P : f(a)\in P \,\}
= \sigma_{\mathbb{L}_{2}}(f(a)).
\]

Ambas composiciones coinciden punto por punto; el primer diagrama conmuta.

\medskip
\item Naturalidad de \(\epsilon\).

Sea \( x\in X_{1} \).  
La rama izquierda envía \( x\mapsto g(x) \), y luego:
\[
\epsilon_{X_{2}}(g(x))
= \{\,U\in \mathcal{D}(X_{2}) : g(x)\in U\,\}.
\]

Por otra parte:
\[
\epsilon_{X_{1}}(x)=\{\,V\in \mathcal{D}(X_{1}) : x\in V\,\}.
\]
Aplicando \( \X(\D(g)) \), que actúa como preimagen, obtenemos:
\[
\X(\D(g))(\epsilon_{X_{1}}(x))
= \{\,U\in \mathcal{D}(X_{2}) : g^{-1}[U]\in\epsilon_{X_{1}}(x)\,\}.
\]

Pero
\[
g^{-1}[U]\in\epsilon_{X_{1}}(x)
\;\Longleftrightarrow\;
x\in g^{-1}[U]
\;\Longleftrightarrow\;
g(x)\in U.
\]

Luego:
\[
\X(\D(g))(\epsilon_{X_{1}}(x))
= \{\,U\in \mathcal{D}(X_{2}) : g(x)\in U\,\}
= \epsilon_{X_{2}}(g(x)).
\]

Así, el segundo diagrama también conmuta.
\end{enumerate}
De esta manera, las  familias
\[
\sigma=\{\sigma_{\mathbb{L}}\}_{\mathbb{L}}\qquad y \qquad \epsilon=\{\epsilon_{X}\}_{X}
\]
constituyen las transformaciones naturales
\[
\sigma\colon \mathrm{Id}_{\mathsf{BDL}} \Rightarrow \D\circ\X. \qquad y \qquad \epsilon\colon \mathrm{Id}_{\mathsf{PS}} \Rightarrow \X\circ\D.
\]
\end{proof}

\medskip

Con la verificación de la naturalidad de las aplicaciones
\(
\sigma=\{\sigma_{\mathbb{L}}\}_{\mathbb{L}}
\)
y
\(
\epsilon=\{\epsilon_{X}\}_{X},
\)
ya contamos con todos los ingredientes necesarios para relacionar
la estructura original con su reconstrucción a través de los funtores
\(\X\) y \(\D\). Aquí se reúnen estos resultados y establece de manera precisa que ambas construcciones forman una dualidad completa entre
retículos distributivos acotados y espacios de Priestley.

\begin{theorem}
Los funtores 
\[
\X\colon \mathsf{BDL}\to\mathsf{PS},
\qquad
\D\colon \mathsf{PS}\to\mathsf{BDL},
\]
junto con las transformaciones naturales
\(
\sigma\colon \mathrm{Id}_{\mathsf{BDL}} \Rightarrow \D\circ\X
\)
y
\(
\epsilon\colon \mathrm{Id}_{\mathsf{PS}} \Rightarrow \X\circ\D,
\)
establecen una dualidad entre las categorías
\(\mathsf{BDL}\) y \(\mathsf{PS}\).
\end{theorem}
Con esto se completa la construcción de ambos funtores, la verificación de su acción sobre objetos y morfismos, y la comprobación de las transformaciones naturales que los vinculan. La conmutatividad de los diagramas anteriores asegura que cada estructura puede recuperarse íntegramente a partir de su dual, consolidando así la dualidad de Priestley como correspondencia precisa entre lo algebraico y lo topológico–ordenado.



\chapter{Dualidad tipo Priestley para $(i,j)$-retículos}\label{Ch 3}



En el presente capítulo, desarrollaremos una dualidad de tipo Priestley para la categoría cuyos objetos son retículos distributivos acotados dotados de ciertos operadores unarios y cuyas flechas serán los homomorfismos de retículos que preservan los operadores en cuestión. Para esto, introduciremos una categoría cuyos objetos son espacios de Priestley dotados de ciertas relaciones binarias y las flechas serán homomorfismos entre dichos espacios que satisfacen ciertas condiciones adicionales. El objetivo de este capítulo será desarrollar una teoría general que abarque todos los operadores unarios normales que pueden ser definidos sobre retículos distributivos acotados.

Comenzaremos introduciendo las definiciones elementales, las categorías involucradas y algunos resultados básicos que serán utilizados a lo largo del presente capítulo. Luego, estudiaremos una equivalencia dual entre las categorías algebraicas de $(i,j)$-retículos con $(-i,-j)$-retículos. Por último, estableceremos la equivalencia dual entre la categoría algebraica de $(i,j)$-retículos con la categoría cuyos objetos son $(i,j)$-espacios de Priestley.  

\section{Definiciones elementales}

El objetivo de esta sección es presentar las nociones básicas que serán utilizadas a través de este capítulo. Comenzaremos introduciendo aquellas necesarias para definir la categoría cuyos objetos son $(i,j)$-retículos. Luego, haremos lo propio para la categoría cuyos objetos son espacios de Priestley relacionales, los cuales, como se verá posteriormente, serán los objetos duales a los $(i,j)$-retículos.
 

\begin{definition} \label{def tipos} Sea $\mathbb{L}$ un retículo distributivo acotado y $m:\mathbb{L}\to \mathbb{L}$ un operador unario. Diremos que
\begin{itemize}
\item[1.] $m$ es de tipo $(1,1)$ si satisface las siguientes condiciones:
\begin{itemize}
\item[a.] $m(0) = 0$
\item[b.] $m(a \vee b) = m(a) \vee m(b)$.
\end{itemize}
\item[2.] $m$ es de tipo $(1,-1)$ si satisface las siguientes condiciones:
\begin{itemize}
\item[c.] $m(0) = 1$
\item[d.] $m(a \vee b) = m(a) \wedge m(b)$.
\end{itemize}
\item[3.] $m$ es de tipo $(-1,1)$ si satisface las siguientes condiciones:
\begin{itemize}
\item[e.] $m(1) = 0$
\item[f.] $m(a \wedge b) = m(a) \vee m(b)$.
\end{itemize}
\item[4.] $m$ es de tipo $(-1,-1)$ si satisface las siguientes condiciones:
\begin{itemize}
\item[g.] $m(1) = 1$
\item[h.] $m(a \wedge b) = m(a) \wedge m(b)$.
\end{itemize}
\end{itemize}
\end{definition}

En \cite[Seccion 1]{CLP} los autores introducen la siguiente definición.

\begin{definition} \label{def: join homo} Sean $\mathbb{L}$ y $\mathbb{M}$ dos retículos distributivos acotados y $j \colon \mathbb{L} \to \mathbb{M}$ una aplicación. Diremos que $j$ es un \emph{join homomorfismo} si satisface las siguientes condiciones:
\begin{itemize}
\item[1.] $j(0^{\mathbb{L}}) = 0^{\mathbb{M}}$
\item[2.] $j(a \vee^{\mathbb{L}} b) = j(a) \vee^{\mathbb{M}} j(b)$
\end{itemize}
\end{definition}

\begin{remark} Sean $\mathbb{L}$ un retículo distributivo acotado y $(i,j) \in\{-1,1\}^{2}$. Teniendo en cuenta la Definición \ref{def tipos} y  Definición \ref{def: join homo}, 
Decimos que $m$ es un operador normal de tipo $(i, j)$ si $m$ es un 
\begin{itemize}
    \item[1.] $m\left(0^{\mathbb{L}^{i}}\right)=0^{\mathbb{L}^{j}}$
    \item[2.] $\forall a,b\in L: m\left(a \vee^{\mathbb{L}^{i}} b\right)= m(a) \vee^{\mathbb{L}^{j}} m(b)$,
\end{itemize}
donde 
\[ 0^{\mathbb{L}^{i}} := \left\{ \begin{array}{cc} 0^{\mathbb{L}} & \text{ si } i = 1 \\ 1^{\mathbb{L}} & \text{ si } i = -1 \end{array} \right. \hspace{0.5cm} \text{ y } \hspace{0.5cm}  a \vee^{\mathbb{L}^{i}} b := \left\{ \begin{array}{cc} a \vee^{\mathbb{L}} b & \text{ si } i = 1 \\ a \wedge^{\mathbb{L}} b & \text{ si } i = -1 \end{array} \right.  \]
Llamaremos $(i,j)$-retículo al par $(\mathbb{L},m)$.
\end{remark}

\begin{remark}\label{rem:interpretacion-ij}
En \cite{CLP} se estudian el concepto de $\vee$-homomorfismo sobre retículos distributivos acotados. Recordemos que un \emph{join-homomorfismo} entre dos retículos distributivos acotados $\mathbb{L}_1$ y $\mathbb{L}_2$ es una función $h\colon \mathbb{L}_1\to\mathbb{L}_2$ que satisface $h(0^{\mathbb{L}_1})=0^{\mathbb{L}_2}$ y $h(a\vee b)=h(a)\vee h(b)$ para todo $a,b\in L_1$. En particular, si $\mathbb{L}$ es un retículo distributivo acotado y $m$ es un $\vee$-homomorfismo sobre $\mathbb{L}$, entonces $(\mathbb{L},m)$ es un $(1,1)$-retículo. Es decir que la noción de $(i,j)$-retículo generaliza la de $\vee$-homomorfismo, pues si $(\mathbb{L},m)$ es un $(i,j)$-retículo entonces $m\colon\mathbb{L}^i\to\mathbb{L}^j$ es un $\vee$-homomorfismo entre los retículos $\mathbb{L}^i$ y $\mathbb{L}^j$. Utilizaremos en el próximo capítulo esta interpretación de los operadores normales de tipo $(i, j)$.
\end{remark}


El estudio de operadores unarios sobre retículos distributivos acotados resulta de gran interés, no sólo desde el punto de vista de la lógica modal, sino también por los diferentes ejemplos que pueden encontrarse, ya sea en el álgebra como en la topología. A continuación brindaremos algunos ejemplos de $(i,j)$-retículos.

\begin{example} \label{Ex 1} Consideremos el par \((\mathbb{L}, \neg)\), donde \(\mathbb{L} = \{0,\, a,\, b,\, 1\}\) es el retículo distributivo acotado de cuatro elementos y el operador unario \( \neg: \mathbb{L} \to \mathbb{L}\) definido de la siguiente manera:

\begin{center}
\begin{tabular}{c|c}
$x$  & \( \neg x \)  \tabularnewline 
\hline
$0$  & $0$     \tabularnewline
$a$  & $b$     \tabularnewline
$b$  & $a$     \tabularnewline
$1$  & $1$     \tabularnewline
\end{tabular}
\end{center}

Algunos cálculos simples muestran que el operador $\neg$ puede ser considerado como un operador normal de tipo $(1,1)$ y un operador normal de tipo $(-1,-1)$ simultáneamente.

\centering
\scalebox{1}{%
\begin{tikzpicture}[node distance=1cm, auto, every node/.style={scale=0.9}]
    % --- Retículo original (superior) ---
    \node (T1) at (0,2) {$1$};
    \node (Ta) at (-1,1) {$a$};
    \node (Ta') at (1,1) {$b$};
    \node (T0) at (0,0) {$0$};
    % Aristas del retículo superior
    \draw (T0) -- (Ta) -- (T1);
    \draw (T0) -- (Ta') -- (T1);
    
    % --- Retículo imagen por negación (inferior), alineado horizontalmente ---
    \begin{scope}[shift={(4,2)}]
        % Recordemos:  \(\neg(0)=1\), \(\neg(a)=a'\), \(\neg(a')=a\) y \(\neg(1)=0\).
        \node (Btop) at (0,0) {$1$};
        \node (Bleft) at (-1,-1) {$b$};
        \node (Bright) at (1,-1) {$a$};
        \node (Bbot) at (0,-2) {$0$};
        % Aristas del retículo inferior
        \draw (Bbot) -- (Bleft) -- (Btop);
        \draw (Bbot) -- (Bright) -- (Btop);
    \end{scope}
    
    % --- Flechas que representan la aplicación de \(\neg\) ---
    \draw[->, dashed] (Ta') -- (Bleft) node[midway, right]{};
\end{tikzpicture}}
\end{example}

\begin{example} \label{Ex 2} Consideremos ahora el par $\left(\mathbb{L},\sim\right)$, donde $\mathbb{L}$ es el retículo de cuatro elementos del ejemplo anterior y $\sim$ el operador unario definido por 

\begin{center}
\begin{tabular}{c|c}
$x$  & \( \sim x \)  \tabularnewline 
\hline
$0$  & $1$     \tabularnewline
$a$  & $b$     \tabularnewline
$b$  & $a$     \tabularnewline
$1$  & $0$     \tabularnewline
\end{tabular}
\end{center}
En este caso, puede observarse que $(\mathbb{L},\sim)$ puede ser considerado un $(1,-1)$ operador y un $(-1,1)$ operador simultáneamente.
\end{example}

\begin{example}\label{ex:clausura-interior}
Sea $(X,\tau)$ un espacio topológico. 
Recordemos de los preliminares las aplicaciones
\[
\Cl_{\tau}(U)=\bigcap\{F\subseteq X : F \text{ es cerrado en } \tau \text{ y } U\subseteq F\},
\qquad
\Int_{\tau}(U)=\bigcup\{O\in\tau : O\subseteq U\}.
\]
Estas funciones pueden considerarse como operadores unarios sobre el retículo $(\mathcal{P}(X),\subseteq)$.

El operador de clausura $\Cl_{\tau}$ satisface 
\[
\Cl_{\tau}(\emptyset)=\emptyset
\quad\text{y}\quad
\Cl_{\tau}(U\cup V)=\Cl_{\tau}(U)\cup \Cl_{\tau}(V),
\]
por lo que es un operador normal de tipo $(1,1)$ sobre $\mathcal{P}(X)$.
De manera dual, el operador de interior $\Int_{\tau}$ cumple
\[
\Int_{\tau}(X)=X
\quad\text{y}\quad
\Int_{\tau}(U\cap V)=\Int_{\tau}(U)\cap \Int_{\tau}(V),
\]
siendo así un operador normal de tipo $(-1,-1)$ sobre el mismo retículo.
\end{example}


\medskip
En el contexto de $(i,j)$-retículos, queremos no sólo preservar la estructura de retículo sino también la acción del operador unario.  
Por tanto, definiremos a continuación qué significa que una función entre dos tales estructuras respete simultáneamente el orden y el operador. Esta noción generaliza la de homomorfismo de retículos al incorporar el comportamiento del operador normal de tipo $(i,j)$.
\medskip


\begin{definition} Sea un par $(i,j) \in \{-1,1\}^{2}$ fijo. Si $(\mathbb{L}_{1},m_{1})$ y $(\mathbb{L}_{2},m_{2})$ son dos $(i,j)$-retículos y $f \colon \mathbb{L}_{1} \to \mathbb{L}_{2}$ es una función. Diremos que $f$ es un homomorfismo de $(i,j)$-retículos si:
\begin{itemize}
\item[1.] $f$ es homomorfismo de retículos.
\item[2.] Para cada $a \in L_{1}$ se tiene que $f(m_1(a)) = m_2(f(a))$. Es decir que el siguiente diagrama conmuta

\[
\begin{tikzcd}[row sep=large, column sep=large]
  \mathbb{L}_1 
    \arrow[r, "f" above] 
    \arrow[d, "m_1" left] 
  & \mathbb{L}_2 
    \arrow[d, "m_2" right] \\
  \mathbb{L}_1 
    \arrow[r, "f" below] 
  & \mathbb{L}_2
\end{tikzcd}
\]
\end{itemize}
\end{definition} 





\begin{proposition} \label{Categoría {$BDL(i,j)$}} Sea $(i,j) \in \{-1,1\}^{2}$. Entonces la clase de los $(i, j)$-retículos, con los homomorfismos de $(i,j)$-retículos forman una categoría, cuya composición de flechas esta dada por la composición usual de funciones y la flecha identidad esta dada por la función identidad. 
\end{proposition}

\begin{proof}[{\textbf{Demostración:}}]
Teniendo en cuenta que la composición de flechas está definida como la composición usual de funciones, bastará con comprobar que la composición usual de homomorfismos de $(i,j)$-retículos es nuevamente un homomorfismo de $(i,j)$-retículos y además que la identidad es un homomorfismo de $(i,j)$-retículos. En efecto:
\begin{itemize}
\item[1.] Composición de morfismos:  

Sean $f\colon (\mathbb{L}_1, m_1)\rightarrow(\mathbb{L}_2, m_2)$ y $g\colon (\mathbb{L}_2, m_2)\rightarrow(\mathbb{L}_3, m_3)$ homomorfismos de \((i,j)\)-retículos.
\begin{itemize}
  \item Por definición, \(f\) y \(g\) son homomorfismos de retículos, luego su composición también es un homomorfismo de retículos.
  \item Además, para todo \(a \in L_1\):
  \[
    (g \circ f)\bigl(m_1(a)\bigr)
      = g\bigl(f(m_1(a))\bigr)
      = g\bigl(m_2(f(a))\bigr)
      = m_3\bigl(g(f(a))\bigr)
      = m_3\bigl((g \circ f)(a)\bigr).
  \]
  Por tanto, \(g\circ f\) conmuta con los operadores.
\end{itemize}
Así \(g\circ f\) es un homomorfismo de \((i,j)\)-retículos. 


\item[2.] Identidad:

Sea $(\mathbb{L},m)$ un $(i,j)$-retículo y definamos la función
\begin{align*}
\Id_{L}&\colon L\to L\\
\Id_{L}&(a) = a
\end{align*}

\medskip
\noindent\emph{(a) Homomorfismo de retículos.}  
Para todo $a,b\in L$,
\begin{align*}
\Id_{L}(a\vee b)&=a\vee b=\Id_{L}(a)\vee\Id_{L}(b),\\
\Id_{L}(a\wedge b)&=a\wedge b=\Id_{L}(a)\wedge\Id_{L}(b),\\
\Id_{L}(0)&=0,\quad\Id_{L}(1)=1.
\end{align*}

\medskip
\noindent\emph{(b) Conmutatividad con el operador.}  
Para todo $a\in L$,
\begin{align*}
\Id_{L}\bigl(m(a)\bigr)&=m(a)=m\bigl(\Id_{L}(a)\bigr)
\end{align*}

\medskip
\noindent\emph{(c) Ley de identidad.}  
Sea $f\colon(\mathbb{L}_1,m_1)\to(\mathbb{L}_2,m_2)$ un homomorfismo de $(i,j)$-retículos. Entonces
\begin{align*}
f\circ\Id_{L}&=f\\
\Id_{L}\circ f&=f
\end{align*}
\medskip
De (a), (b) y (c) se concluye que para cada objeto $(\mathbb{L},m)$ la flecha definida por 
\[
\Id_{(\mathbb{L},m)}:=\Id_{L}
\]
es la flecha identidad.

\end{itemize}
Así, de 1. y 2. se puede concluír que los $(i,j)$-retículos con los homomorfismos de $(i, j)$-retículos constituyen
una categoría.
\end{proof}

Sea $(i,j) \in \{-1,1\}^{2}$. Escribiremos  $\BDL(i,j)$ para indicar la categoría cuyos objetos son $(i,j)$-retículos y cuyas flechas son los homomorfismos de $(i,j)$-retículos.


\vspace{0.2cm}

A continuación, presentaremos las definiciones de las estructuras topológicas relacionales, y sus correspondientes morfismos, para establecer la categoría que resultará dual a $\BDL(i,j)$. Antes de definir estas estructuras, fijaremos una notación que nos será útil.

\begin{definition}\label{def:orden-i}
Sean \((X,\le)\) un conjunto ordenado e \(i\in\{1,-1\}\). Definimos:
\[
x \;\le^i\; y
\;:=\;
\begin{cases}
x\le y,& i=1,\\
y\le x,& i=-1,
\end{cases}
\]
y para cualquier \(U\subseteq X\),
\[
U^i
\;:=\;
\begin{cases}
U,& i=1,\\
X\setminus U,& i=-1.
\end{cases}
\]
\end{definition}


\begin{definition} \label{{$(i, j)$}-Relación} Sean $(X,\tau, \leq)$ un espacio de Priestley, $(i,j) \in\{-1,1\}^{2}$, y $R$ una relación binaria definida sobre $X$. Diremos que $R$ es una {$(i, j)$}-relación si satisface las siguientes condiciones:
\begin{itemize}
    \item[1.] $\forall U\in \mathcal{D}(X): [R^{-1}(U^i)]^j\in \mathcal{D}(X)$
    \item[2.] Si $(x,y)\notin R\Rightarrow \exists U\in \mathcal{D}(X): y\in U^i\text{ y } R(x)\cap U^i=\emptyset$
\end{itemize}
Llamaremos \emph{$(i,j)$-espacio de Priestley} a la estructura $(X,\tau, \leq, R)$, donde $(X,\leq,\tau)$ es un espacio de Priestley y $R$ es una $(i,j)$-relación.
\end{definition}

\begin{definition}\label{def:morfismo-(i,j)espacio}
Sean $(i,j) \in \{-1,1\}^{2}$, $(X, \tau_X, \leq_X, R_X)$ e $(Y, \tau_Y, \leq_Y, R_Y)$ dos $(i,j)$-espacios de Priestley y $f: X \to Y$ una función. Diremos que $f$ es un morfismo de $(i,j)$-espacios de Priestley si satisface las siguientes condiciones:
\begin{enumerate}
    \item $f$ es un morfismo en espacios de Priestley, es decir, $f$ es continua y monótona.
    \item Para todo par \((x, y) \in R_X\), se tiene \((f(x), f(y)) \in R_Y\).
    \item Para todo \(x \in X\) y \(z \in Y\) tales que \((f(x), z) \in R_Y\), existe \(y \in X\) tal que \((x, y) \in R_X\) y \(z \leq^i_Y f(y) \).
\end{enumerate}
\end{definition}


\begin{proposition} La clase de los $(i,j)$-espacios de Priestley con los morfismos de $(i,j)$-espacios de Priestley constituyen una categoría, cuya composición de flechas es la composición usual de funciones y la flecha identidad es la función identidad.  
\end{proposition}


\begin{proof}[\textbf{Demostración}]
Al igual que en la proposición previa, bastará con comprobar que tanto la composición usual de morfismos de $(i,j)$-espacios de Priestley como la identidad son morfismos de $(i,j)$-espacios de Priestley.
\begin{itemize}
\item[1.] Composición de morfismos

Sean $f: (X, \tau_X, \leq_X, R_X) \to (Y, \tau_Y, \leq_Y, R_Y)$ y $g: (Y, \tau_Y, \leq_Y, R_Y) \to (Z, \tau_Z, \leq_Z, R_Z)$ dos morfismos de $(i,j)$-espacios de Priestley. Veamos que la composición satisface las condiciones (1),(2) y (3) de la definicion \ref{def:morfismo-(i,j)espacio}:
\begin{itemize}
    \item Como $f$ y $g$ son morfismos de espacios de Priestley y por lo tanto son continuos y monótonos, su composición $g \circ f: X \to Z$ es continua y monótona, y por lo tanto es un morfismo de espacios de Priestley.
    \item Para todo par \((x, y) \in R_X\), tenemos que $(f(x), f(y)) \in R_Y$, ya que $f$ es un morfismo de $(i,j)$-espacios de Priestley. Puesto que $g$ tambien lo es:
    \[
    (g(f(x)), g(f(y))) \in R_Z.
    \]
    Por lo tanto, $((g \circ f)(x), (g \circ f)(y)) \in R_Z$, y \(g \circ f\) preserva las relaciones.
    \item Sean $x \in X$, $z \in Z$ tal que $(g(f(x)),z) \in R_{Z}$. Dado que $g \colon Y \to Z$ es una flecha, existe $y \in Y$ tal que $(f(x),y) \in R_{Y}$ y $z \leq_{Z}^{i} g(f(x))$. Dado que $f \colon X \to Y$ es una flecha, existe $x^{\prime} \in X$ tal que $(x,x^{\prime}) \in R_{X}$ y $y \leq_{Y}^{i} f(x^{\prime})$. Por la monotonía de $g$, se tiene que $g(y) \leq_{Z}^{i} g(f(x^{\prime}))$. Por lo tanto, $(x,x^{\prime}) \in R_{X}$, con $z \leq_{Z}^{i} g(f(x^{\prime}))$. Por lo tanto, $g \circ f$ satisface la propiedad (3)
\end{itemize}
Por lo tanto, \(g \circ f\) es un morfismo de $(i,j)$-espacios de Priestley.

\item[2.] Identidad

Sea $(X,\tau_X,\le_X,R_X)$ un $(i,j)$-espacio de Priestley. Definimos
\begin{align*}
\Id_X\colon X &\to X\\
x &\mapsto x
\end{align*}
Veamos que también cumple la definición \ref{def:morfismo-(i,j)espacio}.
\begin{itemize}
\item Claramente $\Id_X$ es continua y monótona, luego es un morfismo de espacios de Priestley.

\item Para todo $(x,y)\in R_X$,
\[
(\Id_X(x),\Id_X(y))=(x,y)\in R_X.
\]

\item   Si $(\Id_X(x),z)\in R_X$, es decir $(x,z)\in R_X$, basta tomar $y=z$, de modo que $(x,y)\in R_X$ y  
\[
\Id_X(y)=y\;\le_X^i\;z.
\]
\end{itemize}

De esta manera se concluye que para cada objeto $(X,\tau_X,\le_X,R_X)$ la flecha definida por 
\[
\Id_{(X,\tau_X,\le_X,R_X)} := \Id_X.
\]
es la flecha identidad.

\end{itemize}
Así, por lo anterior, podemos concluir que los $(i,j)$-espacios de Priestley con los $(i,j)$-morfismos de Priestley forman una categoría.
\end{proof}
Sea $(i,j) \in \{-1,1\}^{2}$. Escribiremos $\PS(i,j)$ para denotar la categoría cuyos objetos son los $(i,j)$-espacios de Priestley y cuyas flechas son los morfismos de $(i,j)$-espacios.


\medskip
En síntesis, en esta sección se han introducido las categorías
$\BDL(i,j)$ y $\mathsf{PS}(i,j)$,
que extienden respectivamente la categoría de los retículos distributivos acotados
y la de los espacios de Priestley. En las secciones siguientes construiremos los funtores contravariantes
que establecen la equivalencia dual entre estas categorías, mostrando así que la dualidad de Priestley se extiende naturalmente al contexto de los $(i,j)$-retículos.







\section{La conexión entre estructuras de tipo $(i,j)$ con estructuras $(-i,-j)$}

El objetivo de la presente sección será mostrar que existe una fuerte conexión entre estructuras de tipo $(i,j)$ con las estructuras de tipo $(-i,-j)$, esto es, existe un isomorfismo de categorías entre las categorías $\BDL(i,j)$ y $\BDL(-i,-j)$ y lo propio sucede con las categorías $\PS(i,j)$ y $\PS(-i,-j)$. Estos isomorfismos, los cuales
resultarán de gran ayuda para el objetivo general de este capítulo, capturan el principio de dualidad presente en los retículos y en los espacios de Priestley.

\subsection{Isomorfismo entre $\BDL(i, j)$ y $\BDL(-i,-j)$} \label{S3.2}

En esta subsección estableceremos el isomorfismo entre las categorías $\BDL(i,j)$ y $\BDL(-i,-j)$. Para ello, comenzaremos con una observación estructural que permite reinterpretar cualquier $(i,j)$-retículo como un $(-i,-j)$-retículo al pasar al retículo dual.

Sea $(\mathbb{L},m)$ un $(i,j)$-retículo. Notemos que operador normal de tipo \((i,j)\) puede ser considerado como un operador de tipo \((-i,-j)\) actuando sobre el retículo dual \(\mathbb{L}^{-1}\). Esto se sigue teniendo en cuenta la definición \ref{def:retículo-dual}, ya que en \(\mathbb{L}^{-1}\) se intercambian las operaciones de ínfimo y supremo (y también los elementos \(0\) y \(1\)). 


\vspace{0.2cm}

 
A continuación, haremos uso de la observación previa para definir una asignacion $(-)^{*} \colon \BDL(i,j) \to \BDL(-i,-j)$, que actúa de la siguiente manera:

\vspace{0.2cm}

\noindent \textbf{Accion sobre Objetos:} Sea $(\mathbb{L},m)$ es un $(i,j)$-retículo. Definimos
\[ \left( \mathbb{L},m \right)^{*} = \left(\mathbb{L}^{-1},m^{*}\right), \]
donde $m^{*} = m$, pero actúa sobre $\mathbb{L}^{-1}$.

\vspace{0.2cm}

\noindent \textbf{Sobre Flechas:} Sea $f \colon (\mathbb{L}_{1},m_{1}) \to (\mathbb{L}_{2},m_{2})$ un homomorfismo de $(i,j)$-retículos. Definimos la aplicación $f^{*} \colon (\mathbb{L}_{1},m_{1})^{*} \to (\mathbb{L}_{2},m_{2})^{*}$, donde $f^{*} = f$.

\vspace{0.2cm}

Teniendo en cuenta las asignaciones $(\mathbb{L},m) \mapsto (\mathbb{L},m)^{*}$ y $f \mapsto f^{*}$ podemos establecer el siguiente resultado:

\begin{lemma}\label{Funtor (i,j) a (-i,-j)} Las asignaciones $(\mathbb{L},m) \mapsto (\mathbb{L},m)^{*}$ y $f \mapsto f^{*}$ definen un funtor covariante  $(-)^* \colon \BDL(i,j) \to \BDL(-i,-j)$.
\end{lemma}
\begin{proof}[\textbf{Demostración:}]
Para mostrar que $(-)^*$ es un funtor covariante, veamos que cumple las siguientes propiedades de la Definición \ref{def: funtor}:

\begin{enumerate}
    \item Buena definición sobre objetos:

Sea $(\mathbb{L},m)$ un $(i,j)$-retículo, es decir, se tiene
\begin{align*}
  m\left(0^{\mathbb{L}^i}\right) & = 0^{\mathbb{L}^j},\\[1mm]
  m\left(a\vee^{\mathbb{L}^i} b\right) & = m(a) \vee^{\mathbb{L}^j} m(b) \quad \text{para todo } a,b\in L.
\end{align*}
Teniendo en cuenta que el lema \ref{lem:Li-dual}, se tiene que:
\[ 0^{{(\mathbb{L}^{-1})}^{-i}} = 0^{\mathbb{L}^{i}} \hspace{0.5cm} \text{ y } \hspace{0.5cm} a \vee^{(\mathbb{L}^{-1})^{-i}}  b = a \vee^{\mathbb{L}^{i}} b \]
Asimismo, 
\begin{align*} m^{*}\Bigl(0^{{(\mathbb{L}^{-1})}^{-i}}\Bigr) & = m \left( 0^{\mathbb{L}^{i}} \right) \\
& = 0^{\mathbb{L}^{j}} \\
& = 0^{{(\mathbb{L}^{-1})}^{-j}}
\end{align*}
Por otro lado, notemos que
\begin{align*}
  m^{*}\Bigl(a \vee^{(\mathbb{L}^{-1})^{-i}} b \Bigr) & = m^{*}\Bigl(a \vee^{\mathbb{L}^{i}} b\Bigr) \\
& = m\Bigl(a \vee^{\mathbb{L}^{i}} b\Bigr) \\
& = m(a)\vee^{\mathbb{L}^{j}} m(b) \\
& = m^{*}(a) \vee^{(\mathbb{L}^{-1})^{-j}} m^{*}(b)
\end{align*}
Por lo tanto, $(\mathbb{L},m)^*=\bigl(\mathbb{L}^{-1}, m\bigr)$ es un $(-i,-j)$-retículo, y el la asignación de $(-)^*$ sobre los objetos está bien definida.

\medskip

\item Buena definición sobre flechas: 

Sean $(\mathbb{L}_1, m_1)$, $(\mathbb{L}_2, m_2)$ dos $(i,j)$-retículos y $ f: (\mathbb{L}_1, m_1) \to (\mathbb{L}_2, m_2)$ un homomorfismo de $(i,j)$-retículos. Es sencillo comprobar que $f^* \colon \mathbb{L}_1^{-1} \to \mathbb{L}_2^{-1}$ es un homomorfismo de retículos. Además, se sigue que
\begin{align*}
f^{*}(m^{*}(a)) & = f^{*}(m(a)) \\
                & = f(m(a)) \\
                & = m(f(a)) \\
                & = m^{*}(f^{*}(a))
\end{align*}
Entonces  $f^*$ es un homomorfismo de $(-i,-j)$-retículos. Por lo que la asignación $(-)^{*}$ sobre las flechas está bien definida.

\item Preservación de la identidad:  

Sea $\Id_{(\mathbb{L}, m)}\colon(\mathbb{L}, m)\to(\mathbb{L}, m)$ la flecha identidad en $\BDL(i,j)$. Por la definición de $(-)^*$ sobre flechas:
\begin{align*}
(\Id_{(\mathbb{L}, m)})^* = \Id_{(\mathbb{L}, m)}= \Id_{L}= \Id_{(\mathbb{L}^{-1},m)} = \Id_{(\mathbb{L},m)^*}
\end{align*}
ya que la identidad no altera el orden. De modo que $(-)^*$ envía identidades en identidades.

\medskip
\item Preservación de la composición:

Sean $f\colon(\mathbb{L}_1,m_1)\to(\mathbb{L}_2,m_2)$ y 
$g\colon(\mathbb{L}_2,m_2)\to(\mathbb{L}_3,m_3)$ dos homomorfismos de $\BDL(i,j)$. Entonces:
\begin{align*}
(g\circ f)^* &= g\circ f = g^*\circ f^*
\end{align*}
puesto que por definición $f^* = f$ y $g^* = g$.
\end{enumerate}
Se sigue entonces que el par de asignaciones $(-)^{*} \colon \BDL(i,j) \to \BDL(-i,-j)$ es un funtor covariante. 

\end{proof}
\smallskip

\noindent\textbf{Observación.} Es posible definir el funtor
\[
(-)_{*}\colon \BDL(-i,-j)\longrightarrow \BDL(i,j)
\]
de la misma forma que el funtor $(-)^{*}$, es decir

\noindent \textbf{Sobre Objetos:} Sea $(\mathbb{L},m)$ un $(-i,-j)$-retículo. Definimos
\[ \left( \mathbb{L},m \right)_{*} = \left(\mathbb{L}^{-1},m_{*}\right), \]
donde $m_{*} = m$.

\noindent \textbf{Sobre Flechas:} Sea $f \colon (\mathbb{L}_{1},m_{1}) \to (\mathbb{L}_{2},m_{2})$ un homomorfismo de $(-i,-j)$-retículos. Definimos la aplicación $f_{*} \colon (\mathbb{L}_{1},m_{1})_{*} \to (\mathbb{L}_{2},m_{2})_{*}$, donde $f_{*} = f$.

Se puede mostrar, de la misma manera que en el Lema \ref{Funtor (i,j) a (-i,-j)} que se trata de un funtor covariante.


A continuación, mostraremos que los funtores $(-)^{*}$ y $(-)_{*}$ son inversos y establecen un isomorfismo entre las categorías $\BDL(i,j)$ y $\BDL(-i,-j)$.
\begin{theorem}\label{Equivalencia de $BDL(i,j)$} La categoría $\BDL(i,j)$ es isomorfa a la categoría $\BDL(-i,-j)$
\end{theorem}

\begin{proof}[\textbf{Demostración}]

Veamos que $(-)^{*}$ y $(-)_{*}$ son inversos estrictos y por tanto exhiben un isomorfismo de categorías. Basta verificar que ambas composiciones actúan como la identidad sobre objetos y sobre morfismos.

\medskip
\noindent\textbf{Composición $(-)_{*} \circ (-)^{*}$.}
Sobre objetos, para todo $(i,j)$-retículo  $(\mathbb{L},m)$:
\begin{align*}
\bigl((\mathbb{L},m)^{*}\bigr)_{*}
&= (\mathbb{L}^{-1},m)_{*} \\
&= ((\mathbb{L}^{-1})^{-1},m) \\
&= (\mathbb{L},m)
\end{align*}
Sobre morfismos, sea $f\colon(\mathbb{L}_1,m_1)\to(\mathbb{L}_2,m_2)$ entonces
\begin{align*}
\bigl(f^{*}\bigr)_{*}
&= f_{*}\\
&= f
\end{align*}
De modo que $(-)_{*}\circ(-)^{*}=\Id_{\BDL(i,j)}$

\medskip
\noindent\textbf{Composición $(-)^{*} \circ (-)_{*}$.}
Análogamente, para todo objeto y morfismo en $\BDL(-i,-j)$ se verifica
\begin{align*}
\bigl((\mathbb{L},m)_{*}\bigr)^{*}&=(\mathbb{L},m)\\
\bigl(f_{*}\bigr)^{*}&=f
\end{align*}
de modo que $(-)^{*}\circ(-)_{*}=\Id_{\BDL(-i,-j)}$.

\medskip
Por la definición \ref{def:isomorfismo-categorias}, estos hechos prueban el isomorfismo de categorías:
\begin{align*}
\BDL(i,j)\cong\BDL(-i,-j)
\end{align*}
\end{proof}



\subsection{Isomorfismo entre $\PS(i, j)$ y $\PS(-i,-j)$} \label{S3.3}

Del mismo modo en que hemos conectado la categoría $\BDL(i,j)$ con la categoría $\BDL(-i,-j)$, en la presente sección estableceremos la equivalencia entre las categorías $\PS(i,j)$ y $\PS(-i,-j)$. Esto nos permite trasladar propiedades de  $(i,j)$-espacios de Priestley a sus correspondientes interpretaciones en el $(-i,-j)$ espacio de Priestley isomorfo.

A continuación, se establecerá una asignación  $(-)^{+} \colon \PS(i,j) \to \PS(-i,-j)$, que actúa sobre objetos y sobre morfismos de la siguiente manera: 

\noindent \textbf{Sobre objetos:} Sea $(X,\tau,\le,R)$ un $(i,j)$-espacio de Priestley, entonces
\[ (X,\tau,\leq,R)^{+} = \left(X,\tau,\leq^{-1},R\right) \]
donde recordemos que $\le^{-1}$ representa el orden dual a $\le$. Teniendo en cuenta que $\mathcal{D}(X)$ denota la familia de clopens crecientes para el espacio $(X,\tau,\leq,R)$, escribiremos 
\begin{align*}
\mathcal{D}^{+}(X) &= \{ V \subseteq X : V \text{ es clopen y creciente de } (X,\tau,\leq,R)^{+} \}
\end{align*}
para denotar a los clopens crecientes sobre el orden dual. Nótese que son los clopens decrecientes sobre $(X,\tau,\leq,R)$.

\noindent \textbf{Sobre flechas:} Si $(X_{1},\tau_{1},\leq_{1},R_{1})$ y $(X_{2},\tau_{2},\leq_{2},R_{2})$ son dos $(i,j)$ - espacios de Priestley y $ f \colon X_{1} \to X_{2}$ es un morfismo de $(i,j)$-espacios de Priestley, entonces $f^{+} = f$.  


\begin{lemma} El par de asignaciones $(X,\leq,\tau,R) \mapsto (X,\leq,\tau,R)^{+}$ y $f \mapsto f^{+}$, establecen un funtor $(-)^{+} \colon \mathsf{PS}(i,j) \to \mathsf{PS}(-i,-j)$ covariante.
\end{lemma}
\begin{proof}[\textbf{Demostración}]
Nuevamente, para mostrar que $(-)^{+}$ es un funtor covariante, veamos que cumple las propiedades de la Definición \ref{def: funtor}:

\begin{enumerate}
    \item Buena definición sobre objetos
    
Vamos a comenzar comprobando que la asignación de $(-)^{+} \colon \PS(i,j) \to \PS(-i,-j)$ sobre objetos está bien definida. En efecto, sea $(X,\tau,\leq,R)$ un $(i,j)$-espacio de Priestley. Notemos que la estructura $(X,\tau,\leq^{-1})$ también es un espacio de Priestley, ya que  $(X,\leq^{-1})$ es un conjunto parcialmente ordenado y $(X,\tau)$ es un espacio compacto. Consideremos ahora $x,y\in X$ con $x \nleq^{-1} y$. Dado que $(X,\tau,\leq)$ es un espacio de Priestley, existe un conjunto $U \in \mathcal{D}(X)$ tal que  
\begin{align*}
    y \in U \quad \text{y} \quad x\notin U.
\end{align*}
Consideremos el conjunto $V=U^c$. Es claro que $V \in \mathcal{D}^{+}(X)$. Además, $x\in V$, y dado que $y\in U$, se tiene $y\notin V$. Por lo tanto, $(X,\tau,\leq^{-1})$ verifica la condición de separación de Priestley y en consecuencia $(X,\tau,\leq^{-1})$ es un espacio de Priestley.  \\
\vspace{0.2cm}

Veamos ahora que $R$ es una \((-i,-j)\)-relación, i.e, vamos a comprobar que:
\begin{enumerate}
\item $\forall\, V\in \mathcal{D}^{+}(X), \quad \Bigl[R^{-1}(V^{-i})\Bigr]^{-j} \in \mathcal{D}^{+}(X)$

Para todo $V\in \mathcal{D}^{+}(X)$ se tiene que 
\begin{align*}
V^c \in \mathcal{D}(X)
\end{align*}
y sabemos, puesto que $R$ es una $(i,j)$-relación que
\begin{align*}
\Bigl[ R^{-1}((V^{c})^{i}) \Bigr]^{j} \in \mathcal{D}(X)
\end{align*}
Es fácil ver que $(V^{c})^i=V^{-i}$ para todo $i\in \{-1,1\}$. Por lo que
\begin{align*}
\Bigl[R^{-1}(V^{-i})\Bigr]^{-j}= 
\Bigl[\Bigl(R^{-1}((V^{c})^i)\Bigr)^j\Bigr]^{c}\in \mathcal{D}^{+}(X)
\end{align*}

\item Veamos ahora que si $(x,y)\notin R $ entonces existe $  V\in D^{+}(X) $ tal que $ y\in V^{-i}$ y $R(x)\cap V^{-i}=\emptyset$. En efecto, comencemos asumiendo que  $(x,y)\notin R$. Dado que $R$ es una $(i,j)$-relación, tenemos que $\exists\, U\in \mathcal{D}(X)$ tal que  
$y\in U^i$ y $R(x)\cap U^i=\emptyset$. Consideremos $V=U^c$. Es claro que $V\in \mathcal{D}^{+}(X)$, $V^{-i}=U^{i}$  y se cumple que 
\begin{align*}
    y\in U^i=V^{-i} \text{ y } R(x)\cap U^{i}=R(x)\cap V^{-i}=\emptyset
\end{align*}
\end{enumerate}
Se concluye que $R$ es una \((-i,-j)\)-relación en el espacio $(X,\tau,\leq^{-1})$. Por lo tanto,
\begin{align*}
(X,\tau,\le,R)^{+}=(X,\tau,\leq^{-1},R)
\end{align*}
es un $(-i,-j)$-espacio de Priestley.

\medskip

\item Buena definición sobre flechas

Veamos ahora que la asignación $f \mapsto f^{+}$ sobre las flechas está bien definida. Sean $(X_{1},\tau_{1},\leq_{1},R_{1})$ y $(X_{2},\tau_{2},\leq_{2},R_{2})$ dos $(i,j)$-espacios espacios de Priestley y $f \colon X_{1} \to X_{2}$ un morfismo de $(i,j)$-espacios. Teniendo en cuenta que la aplicación $f$ es continua y creciente, se sigue que la aplicación $f^{+} \colon \left(X_{1},\tau_{1},\leq_{1},R_{1}\right)^{+} \to \left(X_{2},\tau_{2},\leq_{2},R_{2}\right)^{+}$ es también continua y creciente. Asimismo, dado que tanto el espacio $\left(X_{1},\tau_{1},\leq_{1},R_{1}\right)^{+}$ como el espacio $\left(X_{2},\tau_{2},\leq_{2},R_{2}\right)^{+}$ no alteran la relación $R_{1}$ y $R_{2}$, se sigue que $f^{+}$ es una flecha en la categoría $\PS(-i-j)$. 

\item Preservación de la identidad:

Sea $\Id_{(X, \tau_X, \leq_X, R_X)}\colon (X,\tau,\le,R)\to (X,\tau,\le,R)$ la flecha identidad en $PS(i,j)$. Por la definición de $(-)^{+}$ sobre flechas:
\begin{align*}
(\Id_{(X, \tau_X, \leq_X, R_X)})^{+} &= \Id_{(X, \tau_X, \leq_X, R_X)} = \Id_{X} = \Id_{(X,\tau_X,\le_X^{-1},R_X)} = \Id_{(X,\tau_X,\le_X,R_X)^{+}}
\end{align*}
ya que la identidad no altera el orden. De modo que $(-)^{+}$ envía identidades en identidades.

\medskip
\item Preservación de la composición:

Sean $f\colon (X_1,\tau_1,\le_1,R_1)\to (X_2,\tau_2,\le_2,R_2)$ y 
$g\colon (X_2,\tau_2,\le_2,R_2)\to (X_3,\tau_3,\le_3,R_3)$ dos morfismos en $PS(i,j)$. Entonces:
\begin{align*}
(g\circ f)^{+} &= g\circ f\\
&= g^{+} \circ f^{+},
\end{align*}
puesto que por definición $f^{+} = f$ y $g^{+} = g$.
\end{enumerate}

Se sigue entonces que $(-)^{+} \colon \PS(i,j) \to \PS(-i,-j)$ es un funtor covariante. 
\end{proof}

\smallskip


\noindent\textbf{Observación.} Podemos definir el funtor 
\[
(-)_{+}\;\colon\;\PS(-i,-j)\;\longrightarrow\;\PS(i,j)
\]
mediante la misma “dualización” de orden que en $(-)^{+}$:

\begin{itemize}
  \item Si \((X,\tau,\le,R)\) es un objeto de \(\PS(-i,-j)\), entonces
  \[
    (X,\tau,\le,R)_{+} := (X,\tau,\le^{-1},R).
  \]
  \item Si 
  \[
    f\colon (X_1,\tau_1,\le_1,R_1)\to (X_2,\tau_2,\le_2,R_2)
  \]
  es una flecha en \(\PS(-i,-j)\), entonces
  \[
    f_{+} := f.
  \]
\end{itemize}

En el siguiente teorema mostraremos que las composiciones
\((-)_{+}\circ(-)^{+}\) y \((-)^{+}\circ(-)_{+}\) son las identidad correspondientes.



\bigskip

\begin{theorem}\label{Equivalencia de $PS(i,j)$} La categoría $\PS(i,j)$ es isomorfa a $\PS(-i,-j)$.
\end{theorem}
\begin{proof}[\textbf{Demostración}]
Habiendo definido el funtor
\begin{align*}
(-)^{+} : \PS(i,j) &\to \PS(-i,-j)
\end{align*}
mediante la dualización del orden y manteniendo la relación $R$, vemos que esta construcción solo exige que $(i,j)\in\{-1,1\}^2$, por lo que podemos reutilizar idéntica la misma regla para $(-i,-j)$ (de la misma manera que en la demostración del teorema  \ref{Equivalencia de $BDL(i,j)$}), obteniendo
\begin{align*}
(-)_{+} : \PS(-i,-j) &\to \PS(-(-i),-(-j)) = \PS(i,j)
\end{align*}
puesto que $-(-i)=i$ y $-(-j)=j$.

Para concluir que $(-)^{+}$ y $(-)_{+}$ son inversos estrictos y exhiben un isomorfismo de categorías, basta verificar que ambas composiciones son la identidad en objetos y en morfismos.

\medskip
\noindent\textbf{Composición $(-)_{+}\circ(-)^{+}$.}  
Sobre objetos, para todo $(i,j)$-espacio de Priestley $(X,\tau,\le,R)$:
\begin{align*}
\bigl((X,\tau,\le,R)^{+}\bigr)_{+}
&= (X,\tau,\le^{-1},R)_{+}\\
&= (X,\tau,(\le^{-1})^{-1},R)\\
&= (X,\tau,\le,R)
\end{align*}
Sobre morfismos, sea $f\colon (X_1,\tau_1,\le_1,R_1)\to (X_2,\tau_2,\le_2,R_2)$, puesto que en ambos funtores la acción sobre flechas es la identidad, tenemos
\begin{align*}
\bigl(f^{+}\bigr)_{+}
&= f_{+}\\
&= f
\end{align*}
Por lo tanto $(-)_{+}\circ(-)^{+}=\Id_{\PS(i,j)}$.

\medskip
\noindent\textbf{Composición $(-)^{+}\circ(-)_{+}$.}  
Análogamente, para todo objeto y morfismo en $\PS(-i,-j)$ se verifica
\begin{align*}
\bigl((X,\tau,\le,R)_{+}\bigr)^{+}&=(X,\tau,\le,R)\\
\bigl(f_{+}\bigr)^{+}&=f
\end{align*}
de modo que $(-)^{+}\circ(-)_{+}=\Id_{\PS(-i,-j)}$.

\medskip
Por la definición \ref{def:isomorfismo-categorias}, hemos probado el isomorfismo de categorías:
\begin{align*}
\PS(i,j)\cong\PS(-i,-j)
\end{align*}
\end{proof}


\bigskip

Los resultados presentados en las secciones \ref{S3.2} y \ref{S3.3} nos permiten afirmar que las categorías $\BDL(i,j)$ y $\PS(i,j)$ son isomorfas a las categorías $\BDL(-i,-j)$ y $\PS(-i,-j)$, respectivamente. Esta equivalencia será de gran utilidad
en el desarrollo del presente capítulo, ya que nos permitirá simplificar considerablemente las demostraciones de la sección siguiente, evitando el tratamiento redundante de múltiples casos.



\section{Dualidad de tipo Priestley para $(i,j)$-retículos}

El objetivo de la presente sección será establecer una dualidad topológica entre la categoría de $(i,j)$-retículos y la categoría de $(i,j)$-espacios de Priestley. Tal como lo hemos mencionado sobre el final de la sección anterior, los isomorfismos establecidos nos permiten conectar los índices $(i,j)$ y $(-i,-j)$. Por este motivo, las pruebas que haremos en esta sección se harán para los casos $(1,1)$ y $(-1,1)$. \\

\vspace{0.2cm}

Comenzaremos esta sección estableciendo un lema que será de gran utilidad a lo largo del trabajo.

\begin{lemma} \label{Ideales} Sean $m:\mathbb{L}\to \mathbb{L}$, $m$ un operador normal de tipo $(i, j)$ y sea $P$ un filtro primo de $\mathbb{L}$. Entonces $\left(m^{-1}[P]^c\right)^j$ es un $i$-ideal según la definición \ref{def:i-ideal}.
\end{lemma}
\begin{proof}[\textbf{Demostración}:]

Como hemos dicho, probaremos en los casos $(1,1)$ y $(-1,1) $.

\begin{itemize}

\item Caso $(1,1)$

Sea $m : \mathbb{L} \to \mathbb{L}$ un operador normal de tipo $(1,1)$ y $P$ un filtro primo de $\mathbb{L}$. Nuestro objetivo es demostrar que
\begin{align*}
\left(m^{-1}[P]^c\right)^1 = m^{-1}[P]^c = \{ a \in L : m(a) \notin P \}
\end{align*}
es un $1$-ideal, es decir, que es un ideal en $\mathbb{L}$. En efecto:
\begin{enumerate}
  \item \textit{$0^{\mathbb{L}}\in m^{-1}[P]^c$}:
  
  $m(0^{\mathbb{L}}) = 0^{\mathbb{L}} \notin P$, ya que $P$ es propio. Luego $0^{\mathbb{L}}\in m^{-1}[P]^c$.
\item \emph{$m^{-1}[P]^c$ es decreciente} 

Si $a\in m^{-1}[P]^c$ y $b\le a$, entonces
  \begin{align*}
    b\le a \;\Longrightarrow\; m(b)\le m(a).
  \end{align*}
  Si $m(b)\in P$, como $P$ es filtro, $m(a)\in P$, lo cual es un absurdo. Luego $m(b)\notin P$, es decir $b\in m^{-1}[P]^c$.

  \item \emph{Cerradura bajo $\vee$:}
  
Para $a,b\in m^{-1}[P]^c$,
  \begin{align*}
    m(a\vee^{\mathbb{L}} b)=m(a)\vee^{\mathbb{L}} m(b).
  \end{align*}
  Si suponemos $m(a)\vee^{\mathbb{L}} m(b)\in P$, por primalidad de $P$ uno de $m(a)$ o $m(b)$ estaría en $P$, contradicción. Por tanto 
  \begin{align*}
    m(a\vee^{\mathbb{L}} b)\notin P
    \;\Longrightarrow\;
    a\vee^{\mathbb{L}} b\in m^{-1}[P]^c.
  \end{align*}
\end{enumerate}

\medskip

\item Caso $(-1, 1)$:

Sea $m : \mathbb{L} \to \mathbb{L}$ un operador normal de tipo $(-1,1)$ y sea $P$ un filtro primo de $\mathbb{L}$. Queremos demostrar que 
\begin{align*}
\left(m^{-1}[P]^c\right)^1 = m^{-1}[P]^c = \{ a \in L : m(a) \notin P \}
\end{align*}
es un $(-1)$-ideal, es decir, un filtro en $\mathbb{L}$:

\begin{enumerate}
  \item \textit{$1\in m^{-1}[P]^c$}:
  
$m(1^{\mathbb{L}})=m(0^{\mathbb{L}^{-1}})=0^{\mathbb{L}}\notin P$ ya que $P$ es propio. Luego $1^{\mathbb{L}}\in m^{-1}[P]^c$.



  \item \emph{$m^{-1}[P]^c$ es creciente:} 
  

  Sea $a\in m^{-1}[P]^c$, es decir, $m(a)\notin P$, y sea $b\in L$ tal que
$b\leq^{-1} a$. Por la definición del orden dual, esto equivale a $a\leq b$ en
$\mathbb{L}$, y en el retículo dual se cumple
\[
a\vee^{\mathbb{L}^{-1}} b \;=\; a.
\]

Supongamos, hacia una contradicción, que $b\notin m^{-1}[P]^c$, esto es,
$m(b)\in P$. Como $m$ es un operador normal de tipo $(-1,1)$, preserva la
operación $\vee^{\mathbb{L}^{-1}}$, de modo que
\[
m(a)
\;=\;
m\bigl(a\vee^{\mathbb{L}^{-1}} b\bigr)
\;=\;
m(a)\vee^{\mathbb{L}} m(b).
\]
Si $m(b)\in P$, entonces $m(a)\vee^{\mathbb{L}} m(b)\in P$, y por lo tanto
$m(a)\in P$, lo cual contradice $a\in m^{-1}[P]^c$. Luego necesariamente
$m(b)\notin P$, es decir, $b\in m^{-1}[P]^c$. Por lo tanto, $m^{-1}[P]^c$ es
decreciente respecto de $\leq^{-1}$, o lo que es lo mismo, creciente respecto del orden de $\mathbb{L}$.

\item \emph{Cerradura bajo $\wedge$:}

Debemos probar que si $a,b\in m^{-1}[P]^c$, entonces
$a\vee^{\mathbb{L}^{-1}} b\in m^{-1}[P]^c$.

Sean $a,b\in m^{-1}[P]^c$, esto es, $m(a)\notin P$ y $m(b)\notin P$. Como $m$
es de tipo $(-1,1)$, preserva la operación $\vee^{\mathbb{L}^{-1}}$, de manera
que
\[
m\bigl(a\vee^{\mathbb{L}^{-1}} b\bigr)
\;=\;
m(a)\vee^{\mathbb{L}} m(b).
\]
Supongamos que $m\bigl(a\vee^{\mathbb{L}^{-1}} b\bigr)\in P$. Entonces
$m(a)\vee^{\mathbb{L}} m(b)\in P$, y como $P$ es un filtro primo, se sigue que
$m(a)\in P$ o $m(b)\in P$, en contradicción con $a,b\in m^{-1}[P]^c$. Por lo tanto,
\[
m\bigl(a\vee^{\mathbb{L}^{-1}} b\bigr)\notin P,
\]
es decir,
\[
a\vee^{\mathbb{L}^{-1}} b\in m^{-1}[P]^c.
\]
Así, $m^{-1}[P]^c$ es cerrado bajo $\vee^{\mathbb{L}^{-1}}$, es decir bajo $\wedge$ en $\mathbb{L}$.
\end{enumerate}
\end{itemize}
Así hemos probado que para ambos casos que $\left(m^{-1}[P]^c\right)^j$ es un $i$-ideal.
\end{proof}

\subsection{El $(i,j)$-espacio de Priestley asociado a un $(i,j)$-retículo distributivo}

En esta parte del capítulo nos adentraremos en la construcción que asocia a cada retículo distributivo acotado con un operador normal de tipo $(i, j)$, un $(i, j)$ espacio de Priestley. Lejos de ser un mero ejercicio de definición, esta representación sitúa la estructura algebraica en un universo topológico donde la compacidad y la separación de clopens se transforman en herramientas que nos permiten estudiar las propiedades de los operadores normales de tipo $(i, j)$. La atención aquí se centra en cómo la dualidad clásica se extiende al contexto de los operadores en cuestión.\\


Sea $(\mathbb{L},m)$ un $(i,j)$-retículo. Consideremos la estructura 
\[  \mathfrak{X}(\mathbb{L},m) := \left({\mathcal{X}}(\mathbb{L}),\tau,\subseteq,R_{m} \right), \]
donde $\left({\mathcal{X}}(\mathbb{L}),\tau,\subseteq\right)$ es el espacio de Priestley asociado a $\mathbb{L}$ y $R_{m} \subseteq \mathcal{X}(\mathbb{L}) \times \mathcal{X}(\mathbb{L})$ es una relación binaria definida de la siguiente manera:
\[ (P,Q) \in R_{m} \text{ si y sólo si } Q^{i} \subseteq m^{-1}[P]^{j}. \]


\begin{lemma}\label{Lema de Representacion de BDL(i,j)}
Si $(\mathbb{L}, m)$ un $(i,j)$-retículo entonces $ \mathfrak{X}(\mathbb{L},m)$ es un $(i,j)$-espacio de Priestley. 
\end{lemma}




\begin{proof}[{\textbf{Demostración:}}]Teniendo en cuenta la dualidad de Priestley para retículos distributivos acotados, sabemos que $(X,\tau,\leq)$ es un espacio de Priestley. Por lo tanto, será suficiente comprobar que $R_m$ es una $(i,j)$-relación. Tal como lo hemos mencionado previamente, haremos las demostraciones de los casos $(1,1)$ y $(-1,1)$, respectivamente.\\

\vspace{0.2cm}

\begin{itemize}
    \item Caso $(1,1)$


Veamos que $R_m$ es una $(1,1)$-relación. Es decir, probaremos que
\begin{enumerate}
\item $\forall U\in \mathcal{D}(\mathcal{X}(\mathbb{L})): [R_m^{-1}(U^i)]^j = R_m^{-1}(U)\in \mathcal{D}(\mathcal{X}(\mathbb{L}))$

\vspace{0.2cm}

Consideremos $U \in \mathcal{D}(\mathcal{X}(\mathbb{L}))$. Sabemos, por la dualidad de Priesley del Capítulo 2, que $U$ es de la forma $U = \sigma_{\mathbb{L}}(a)$ para algún $a\in L$. Vamos a comprobar que
\begin{align*}
    R_m^{-1}[\sigma_{\mathbb{L}}(a)] = \sigma_{\mathbb{L}}(m(a)).
\end{align*}
Sea $P\in R_m^{-1}(\sigma_{\mathbb{L}}(a))$. Entonces $R_m(P)\cap \sigma_{\mathbb{L}}(a) \neq \emptyset$, por lo tanto existe $Q\in \mathcal{X}(\mathbb{L})$ tal que $(P,Q)\in R_m$ y $a\in Q$. Como $(P,Q)\in R_m$ implica $Q\subseteq m^{-1}[P]$, tenemos que  $a\in Q \subseteq m^{-1}[P]$, es decir, $m(a)\in P$. Por lo tanto, $P\in \sigma_{\mathbb{L}}(m(a))$.\\  

\vspace{0.2cm}

Consideremos ahora $P\in \sigma_{\mathbb{L}}(m(a))$, es decir que $m(a)\in P$ i.e., $a\in m^{-1}[P]$. Teniendo en cuenta el Lema \ref{Ideales}, se sigue que $m^{-1}[P]^c$ es un $1$-ideal y además notemos que 
\begin{align*}
      [a)\cap m^{-1}[P]^c = \emptyset.
\end{align*}
Por el Teorema del filtro primo (Teorema \ref{thm:filtro_primo}), existe $Q\in\mathcal{X}(\mathbb{L})$ tal que $a\in Q \quad \text{y} \quad Q\cap m^{-1}[P]^c = \emptyset$. Esto último es equivalente a decir que
$Q\subseteq m^{-1}[P]$. En consecuencia $(P,Q)\in R_m$ y $a\in Q$; por lo tanto, $Q\in \sigma_{\mathbb{L}}(a)$ y se sigue que $P\in R_m^{-1}(\sigma_{\mathbb{L}}(a))$. \\

\vspace{0.2cm}

\item Si $(P,Q)\notin R_m\Rightarrow \exists U\in \mathcal{D}(\mathcal{X}(\mathbb{L})): Q\in U^i = U\text{ y } R_m(P)\cap U=\emptyset$

\vspace{0.2cm}

Si $(P,Q)\notin R_m$, entonces $Q\nsubseteq m^{-1}[P]$; es decir, existe $a\in Q$ con $a\notin m^{-1}[P]$. Sea
\begin{align*}
        U=\sigma_{\mathbb{L}}(a).
\end{align*}
Claramente, $Q\in \sigma_{\mathbb{L}}(a)=U$. Supongamos, buscando una contradicción, que existe $Z\in R_m(P)\cap U$. Entonces, $Z\in \sigma_{\mathbb{L}}(a)$ implica que $a\in Z$ y, dado que $Z\subseteq m^{-1}[P]$, se obtendría $a\in m^{-1}[P]$, lo cual contradice nuestra suposición. Así, se concluye que $R_m(P)\cap U = \emptyset$.
\end{enumerate}
En consecuencia, se demuestra que $R_m$ es una $(1,1)$-relación y 
\begin{align*}
(\mathcal{X}(\mathbb{L}), \tau_L, \subseteq, R_m)
\end{align*}
es un $(1,1)$-espacio de Priestley.\\

\vspace{0.2cm}

\item Caso $(-1,1)$


Analizaremos ahora el caso $(-1,1)$. Veamos que $R_m$ es una $(-1,1)$-relación. Es decir, probaremos que
\begin{enumerate}
\item $\forall U\in \mathcal{D}(\mathcal{X}(\mathbb{L})): [R_m^{-1}(U^i)]^j = R_m^{-1}(U^c)\in \mathcal{D}(\mathcal{X}(\mathbb{L}))$

\vspace{0.2cm}


Consideremos \(U \in \mathcal{D}(\mathcal{X}(\mathbb{L}))\). Es decir, $U=\sigma_{\mathbb{L}}(a)$ para algún $a\in L$. De igual manera que en el caso anterior, bastará demostrar que
\begin{align*}
R_m^{-1}\bigl(\sigma_{\mathbb{L}}(a)^c\bigr)=\sigma_{\mathbb{L}}\bigl(m(a)\bigr)
\end{align*}

Sea \(P\in R_m^{-1}\bigl(\sigma_{\mathbb{L}}(a)^c\bigr)\). Entonces, $R_m(P)\cap \sigma_{\mathbb{L}}(a)^c\neq\emptyset$, por lo tanto, existe \(Q\in \mathcal{X}(\mathbb{L})\) tal que
\begin{align*}
    (P,Q)\in R_m \quad \text{y} \quad Q\in \sigma_{\mathbb{L}}(a)^c,
\end{align*}
es decir, \(a\notin Q\) y \((P,Q)\in R_m\). Por definición de $R_{m}$, se sigue que $Q^c\subseteq m^{-1}[P]$. De esta manera $m(a)\in P$ i.e., \(P\in \sigma_{\mathbb{L}}\bigl(m(a)\bigr)\).\\

\vspace{0.2cm}
    
Consideremos ahora \(P \in \sigma_{\mathbb{L}}\bigl(m(a)\bigr)\), es decir, \(m(a)\in P\). Esto implica que $a\in m^{-1}[P]$. Teniendo en cuenta el Lema \ref{Ideales}, se sigue que \(m^{-1}[P]^c\) es un $i$-ideal, i.e., es un filtro. Puesto que $a\notin m^{-1}[P]^c$, el decreciente generado por \(a\), denotado \((a]\), no interseca a \(m^{-1}[P]^c\); en otras palabras,
\begin{align*}
    (a]\cap m^{-1}[P]^c=\emptyset.
\end{align*}
Por el teorema del filtro primo, existe un filtro primo \(Q\in \mathcal{X}(\mathbb{L})\) tal que
\begin{align*}
a\notin Q \quad \text{y} \quad m^{-1}[P]^c\subseteq Q
\end{align*}
Esto último implica que $Q^c\subseteq m^{-1}[P]$. Se sigue entonces que \((P,Q)\in R_m\) y, además, \(a\notin Q\); en otras palabras, \(Q\in \sigma_{\mathbb{L}}(a)^c\). Por lo tanto,
\begin{align*}
P\in R_m^{-1}\bigl(\sigma_{\mathbb{L}}(a)^c\bigr).
\end{align*}


\item Si $(P,Q)\notin R_m\Rightarrow \exists U\in \mathcal{D}(\mathcal{X}(\mathbb{L})): Q\in U^i = U^c\text{ y } R_m(P)\cap U^c=\emptyset$

\vspace{0.2cm}

Supongamos ahora que \((P,Q)\notin R_m\). Se sigue entonces que \(Q^c\nsubseteq m^{-1}[P]\); es decir, existe \(a\in Q^c\) con \(a\notin m^{-1}[P]\). Sea $U=\sigma_{\mathbb{L}}(a)$. Notemos que $Q\in \sigma_{\mathbb{L}}(a)^c=U^c$. Supongamos, buscando una contradicción, que existe \(Z\in R_m(P)\cap U^c\). Estos datos implican que
\begin{align*}
Z^c\subseteq m^{-1}[P] \quad \text{y} \quad a\notin Z 
\end{align*}
y por lo tanto se obtendría que $a\in m^{-1}[P]$, lo cual contradice la suposición que teniamos de $a\notin m^{-1}[P]$. Así, se concluye que
\begin{align*}
R_m(P)\cap U^c=\emptyset.
\end{align*}
\end{enumerate}

Por lo tanto, \(R_m\) es una \((-1,1)\)-relación y, en consecuencia,
\begin{align*}
(\mathcal{X}(\mathbb{L}),\tau_L,\subseteq, R_m)
\end{align*}
es un \((-1,1)\)-espacio de Priestley.
\end{itemize}
\end{proof}


Consideremos ahora $(\mathbb{L}_1, m_1)$ y $(\mathbb{L}_2, m_2)$ dos objetos de la categoría $\BDL(i,j)$ y sea 
$f: (\mathbb{L}_1, m_1) \to (\mathbb{L}_2, m_2)$ una flecha en $\BDL(i,j)$. Por la dualidad de Priestley para retículos distributivos acotados, sabemos que la aplicación $\mathfrak{X}(f) \colon \mathcal{X}(\mathbb{L}_2) \to \mathcal{X}(\mathbb{L}_1)$ definida por $\mathfrak{X}(f)(P) = f^{-1}[P]$ es una función continua y monótona entre los espacios de Priestley duales a los retículos. A continuación, vamos a comprobar que es una flecha en la categoría $\PS(i,j)$.


\begin{lemma} \label{Lema de Functorialidad de BDL(i,j)} 
Sean $(\mathbb{L}_1, m_1)$ y $(\mathbb{L}_2, m_2)$ dos $(i,j)$-retículos y  
\begin{align*}
f: (\mathbb{L}_1, m_1) \to (\mathbb{L}_2, m_2)
\end{align*}
un homomorfismo de $(i,j)$-retículos. Entonces la aplicación 
\begin{align*}
\mathfrak{X}(f)&: (\mathcal{X}(\mathbb{L}_2), \tau_{L_2}, \subseteq, R_{m_2}) \to (\mathcal{X}(\mathbb{L}_1), \tau_{L_1}, \subseteq, R_{m_1})
\end{align*}
definida por $\mathfrak{X}(f)(P) := f^{-1}[P]$, es un morfismo de $(i,j)$-espacio de Priestley.
\end{lemma}

\begin{proof}[{\textbf{Demostración:}}]Probaremos entonces que $\mathfrak{X}(f)$ es un morfismo de $(i,j)$-espacios de Priestley. Teniendo en cuenta los resultados ya probados en la dualidad de Priestley, ya sabemos que $\mathfrak{X}(f)$ es continuo y monótono. Nuevamente dividiremos la demostración en los casos $(1,1)$ y $(-1,1)$ y veremos el resto de las propiedades.

\begin{itemize}
    \item Caso $(1,1)$

Nos falta comprobar las siguientes condiciones:
\begin{enumerate}
    \item Para cada $P,Q \in \mathcal{X}(\mathbb{L}_{2})$, si $(P,Q) \in R_{m_2}$ entonces $(f^{-1}[P], f^{-1}[Q]) \in R_{m_1}$

\vspace{0.2cm}

Sean $P,Q \in \mathcal{X}(\mathbb{L}_{2})$ tales que $Q \subseteq m_{2}^{-1}[P]$, vamos a comprobar que $\mathfrak{X}(f)(Q) \subseteq m_{1}^{-1}\left[ f^{-1}[P] \right]$ i.e., $f^{-1}[Q] \subseteq m_{1}^{-1}\left(f^{-1}[P]\right)$. Consideremos $a \in f^{-1}[Q]$, se sigue que $f(a) \in Q$. Por hipótesis, $Q \subseteq m_{2}^{-1}[P]$, de donde se sigue que $f(a) \in m_{2}^{-1}[P]$ i.e., $m_{2}(f(a)) \in P$. Dado que $f$ es una flecha en $\BDL(i,j)$, se sigue que $m_{2}(f(a)) = f(m_{1}(a))$, de donde se sigue que $m_{1}(a) \in f^{-1}[P]$ y por lo tanto, $a \in m_{1}^{-1}[f^{-1}[P]]$, como queríamos demostrar.

\vspace{0.2cm}

    \item Si $P \in \mathcal{X}(\mathbb{L}_{2})$, $Z \in \mathcal{X}(\mathbb{L}_{1})$ y $(f^{-1}[P], Z) \in R_{m_1} $ entonces existe $Q \in \mathcal{X}(L_2)$ tal que $ (P,Q) \in R_{m_2}$ y $Z \subseteq f^{-1}[Q]$

\vspace{0.2cm}

Consideremos ahora $P \in \mathcal{X}(\mathbb{L}_{2})$, $Z \in \mathcal{X}(\mathbb{L}_{1})$ tales que $(f^{-1}[P], Z) \in R_{m_1}$
se sigue entonces que $Z \subseteq m_1^{-1}[f^{-1}[P]]$. De esto último, tenemos que  $f(Z) \subseteq m_2^{-1}[P]$. Consideremos ahora el filtro generado por $f[Z]$,  $F = \Fg^{\mathbb{L}_{2}}\left( f[Z] \right)$. Vamos a comprobar que 
\[  F \cap m_2^{-1}[P]^c = \emptyset \]
En efecto, si $a \in F \cap m_{2}^{-1}[P]^{c}$ entonces existe $z \in Z$ tal que $f(z) \leq a$ y $m_{2}(a) \notin P$. Dado que $m_{2}$ es de tipo $(1,1)$, $m_{2}(f(z)) \leq m_{2}(a)$ y por lo tanto, $m_{2}(f(z)) \notin P$. Dado que $f$ es flecha en $\BDL(i,j)$, se sigue que $z \notin m_{1}^{-1}[f^{-1}[P]]$, lo que resulta en una contradicción, ya que $z \in Z \subseteq m_{1}^{-1}[f^{-1}[P]]$. Luego, por el teorema del filtro primo, existe $Q \in \mathcal{X}(L_2)$ tal que $F \subseteq Q$ y $Q \subseteq m_2^{-1}[P]$, lo que implica que 
\[  Z \subseteq f^{-1}[Q] \hspace{0.2cm} \text{ y } \hspace{0.2cm} Q \in R_{m_2}(P) \]
\end{enumerate}
Por lo que $\mathfrak{X}(f)$ es una flecha de $\PS(1,1)$.

\item Caso $(-1,1)$

Ahora la condición $Q^i\subseteq m_2^{-1}[P]$ se traduce en $m_2^{-1}[P]\supseteq Q^{-1}=Q^c$, y la segunda inclusión debe invertirse:
\begin{enumerate}
    \item $\forall(P,Q) \in R_{m_2}\Rightarrow (f^{-1}[P], f^{-1}[Q]) \in R_{m_1}$
    
\vspace{0.2cm}
    
    La propiedad de preservación de la relación se prueba de la misma manera que para el caso $(1,1)$, ya que no depende de $i$. 
    
\vspace{0.2cm}
    
    \item $\forall(f^{-1}[P], Z) \in R_{m_1}\Rightarrow\exists Q \in \mathcal{X}(L_2): Q \in R_{m_2}(P)$ y $Z \supseteq f^{-1}[Q]$
    
\vspace{0.2cm}

Por definición de $R_{m_1}$, la condición $(f^{-1}[P],Z)\in R_{m_1}$ equivale a $Z^c\subseteq m_1^{-1}\bigl(f^{-1}[P]\bigr)$. Al aplicar la flecha $f$ obtenemos $f(Z^c)\subseteq m_2^{-1}[P]$. 
Sea ahora el ideal generado por $f[Z^c]$, $I = \Ig^{\mathbb{L}_{2}}\left( f[Z^c] \right)$ Supongamos, para llegar a contradicción, que existe $a\in I\cap m_2^{-1}[P]^c$. Entonces hay $z\in Z^c$ con $f(z)\le a$ y además $m_2(a)\notin P$. Como $m_2$ es monótono, se cumple $m_2\bigl(f(z)\bigr)\le m_2(a)\notin P$. Pero la conmutatividad $f\circ m_1 = m_2\circ f$ implica $f\bigl(m_1(z)\bigr)=m_2\bigl(f(z)\bigr)\notin P$,  
lo que significa $m_1(z)\notin f^{-1}[P]$, en contradicción con $z\in Z^c\subseteq m_1^{-1}[f^{-1}[P]]$. De ello se sigue  
\[I\cap m_2^{-1}[P]^c=\emptyset\] 
Como $P$ es primo, existe entonces un filtro primo $Q\subseteq L_2$ que contiene a $m_2^{-1}[P]^c$ y es disjunto de $I$. Puesto que $f(Z^c)\subseteq I\subseteq Q$, tenemos $Z^c\subseteq f^{-1}[Q]$ y por lo tanto $f^{-1}[Q]\supseteq Z$. Finalmente, $Q^c\subseteq m_2^{-1}[P]$ da $(P,Q)\in R_{m_2}$, como queríamos.
\end{enumerate}
\end{itemize}
\end{proof}
Así hemos comprobado que a todo $(i,j)$‑retículo corresponde un $(i,j)$‑espacio de Priestley de manera natural, utilizando ni mas ni menos que una extensión del funtor $\mathfrak{X}$ de la dualidad de Priestley original.


\subsection{El $(i,j)$-retículo distributivo asociado a un $(i,j)$-espacio de Priestley}

Pasamos ahora al sentido inverso: asociar un $(i,j)$-retículo distributivo a un $(i,j)$-espacio de Priestley. Esta construcción extrae del espacio su retículo de clopens crecientes y define sobre él un operador normal de tipo $(i, j)$ que refleja la acción de la relación. De este modo, hacemos visible la correspondencia que sustentará la dualidad.


Sea $(X, \tau, \leq, R)$ un $(i,j)$-espacio de Priestley. Consideremos la estructura
\begin{align*}
\D&(X) := (\mathcal{D}(X), \cap, \cup, m_{R} ,\emptyset, X),
\end{align*}
donde
\begin{itemize}
\item $(\mathcal{D}(X), \cap, \cup ,\emptyset, X)$ es el retículo distributivo acotado dado por los clopens crecientes de $(X,\tau,\leq)$.
\item $m_R$ es un operador unario sobre $\mathcal{D}(X)$ definido por 
\begin{align*} m_{R}(U) & := R^{-1}(U^i)^j 
\end{align*}    
\end{itemize}


\begin{lemma}\label{Lema de Reconstrucción} Si $(X, \tau, \leq, R)$ es un $(i,j)$-espacio de Priestley entonces $\D(X)$ es un $(i,j)$-retículo.
\end{lemma}
\begin{proof}[{\textbf{Demostración:}}]Mostraremos, nuevamente en dos casos, que $\D(X)$ es un $(i,j)$-retículo. Es claro que $\left(\mathcal{D}(X),\cap,\cup,\emptyset,X\right)$ es un retículo distributivo acotado, por lo que debemos ver que $m_{R}(U)$ es de tipo $(i,j)$. Es decir que para todo $(i,j)\in \{-1,1\}^2$, lo que se busca es que nos quede un operador 
    \[m_R:\D(X)^i\to \D(X)^j\]
que sea un $\vee$-homomorfismo, de acuerdo a la Observación \ref{rem:interpretacion-ij}. Es decir, debemos probar que satisface lo siguiente:
    \begin{enumerate}
     \item $m_R\left(0^{\D(X)^{i}}\right)=0^{\D(X)^j}$
     \item $m_R(A\vee^{\D(X)^{i}}B)=m_R(A)\vee^{\D(X)^{j}}m_R(B)$ 
    \end{enumerate} 
\begin{itemize}
    \item Caso $(1,1)$

Debemos probar que $m_R$ es de tipo $(1,1)$. Recordemos que $0^{\D(X)}=\emptyset$ y $A\vee^{\D(X)}B=A\cup B$. Es decir que las condiciones se traducen en:
\begin{enumerate}
  \item $m_R(\emptyset) = \emptyset$
    \begin{align*}
    m_R(\emptyset) &= \{x : R(x) \cap \emptyset \neq \emptyset\} = \emptyset
    \end{align*}
  \item $\forall\, A,B \in \mathcal{D}(X): m_R(A \cup B) = m_R(A) \cup m_R(B)$
    \begin{align*}
    m_R(A \cup B) &= \{x : R(x) \cap (A \cup B) \neq \emptyset\} \\
    &= \{x : R(x) \cap A \neq \emptyset\} \cup \{x : R(x) \cap B \neq \emptyset\} \\
    &= m_R(A) \cup m_R(B).
    \end{align*}
\end{enumerate}
Por lo tanto, $(\mathcal{D}(X), \cap, \cup, \emptyset, X, m_R)$ es un $(1,1)$-retículo.

\item Caso $(-1,1)$

Veamos que $m_R$ es de tipo $(-1,1)$. Recordemos que $0^{\D(X)^{-1}}=X$ y $A\vee^{\D(X)^{-1}}B=A\cap B$. Es decir que las condiciones se traducen en:

\begin{enumerate}
    \item $m_R(X) = \emptyset$
    \begin{align*}
    m_R(X) &= \{x : R(x) \cap X^c \neq \emptyset\}\\
    & = \{x : R(x) \cap \emptyset \neq \emptyset\}\\
    & = \emptyset
    \end{align*}
    \item $\forall\, A,B \in \mathcal{D}(X): m_R(A \cap B) = m_R(A) \cup m_R(B)$
    \begin{align*}
    m_R(A \cap B) &= \{x : R(x) \cap (A \cap B)^c \neq \emptyset\} \\
    &= \{x : R(x) \cap (A^c \cup B^c) \neq \emptyset\} \\
    &= \{x : R(x) \cap A^c \neq \emptyset\} \cup \{x : R(x) \cap B^c \neq \emptyset\} \\
    &= m_R(A) \cup m_R(B).
    \end{align*}
\end{enumerate}
Luego $(\mathcal{D}(X), \cap, \cup, \emptyset, X, m_R)$ es un $(-1,1)$-retículo.


\end{itemize}
\end{proof}

Sea $f: (X_{1}, \tau_{1}, \leq_{1}, R_{1}) \to (X_{2}, \tau_{2}, \leq_{2}, R_{2})$ un morfismo de $(i,j)$-espacios de Priestley. Definimos la aplicación 
\[ \D(f): (\mathcal{D}(X_2
), \cap, \cup, \emptyset, X_2, m_{R_{2}}) \to (\mathcal{D}(X_{1}), \cap, \cup, \emptyset, X_1, m_{R_{1}}),\] 
donde $\D(f)(U) := f^{-1}[U]$ para cada $U \in \mathcal{D}(X_{2})$.


\begin{lemma}\label{Lema de Functorialidad para D} Si $f: (X_{1}, \tau_{1}, \leq_{1}, R_{1}) \to (X_{2}, \tau_{2}, \leq_{2}, R_{2})$ es una flecha en $\PS(i,j)$ entonces la aplicación $\D(f) \colon {\cal{D}}(X_{2})\to {\cal{D}}(X_{1})$ es una flecha en $\BDL(i,j)$.
\end{lemma}
\begin{proof}[{\textbf{Demostración:}}]Otra vez, veamos para ambos casos que $\D(f)$ es un morfismo de $(i, j)$-retículos. Puesto que se encuentra definido igual que en la dualidad de Priestley, $\D(f)$ es un morfismo de retículos distributivos acotados. Veamos que conmuta con los operadores $m_R$. Sea $U\in \mathcal{D}(Y)$. Debemos probar que
\begin{align*}
\D(g)(m_S(U)) = m_R(\D(g)(U)),
\end{align*}
es decir
\begin{align*}
g^{-1}[m_S(U)] = m_R(g^{-1}[U]).
\end{align*}
Procedemos a demostrar la doble contención.

\medskip
\begin{itemize}
    \item Caso $(1,1)$: 

\begin{itemize}
  \item[$\subseteq$)] Si \(x\in g^{-1}[m_{S}(U)]\), entonces \(g(x)\in m_{S}(U)\), esto es \(S(g(x))\cap U\neq\emptyset\). Luego existe \(z\in U\) con \((g(x),z)\in S\). Como \(g\) es morfismo de espacios, existe \(y\in X\) tal que \((x,y)\in R\) y \(z\le g(y)\). Pero \(z\in U\) y \(U\) es creciente, así \(g(y)\in U\) y por tanto \(y\in R(x)\cap g^{-1}[U]\neq\emptyset\). Hemos probado que 
  \[
    x\in m_{R}\bigl(g^{-1}[U]\bigr).
  \]

  \item[$\supseteq$)] Si \(x\in m_{R}(g^{-1}[U])\), entonces existe \(y\in X\) con \((x,y)\in R\) y \(g(y)\in U\). Por homomorfismo \((g(x),g(y))\in S\), así \(S(g(x))\cap U\neq\emptyset\) y \(g(x)\in m_{S}(U)\). Luego \(x\in g^{-1}[m_{S}(U)]\).
\end{itemize}

\medskip

\item Caso $(-1,1)$: 

Recordemos que en este caso
\[
m_{S}(U)=\{\,x: S(x)\cap U^{c}\neq\emptyset\},
\quad
m_{R}(V)=\{\,x: R(x)\cap V^{c}\neq\emptyset\}.
\]

\begin{enumerate}
  \item[$\subseteq$)] Si \(x\in g^{-1}[m_{S}(U)]\), entonces \(S(g(x))\cap U^{c}\neq\emptyset\), y existe \(z\in U^{c}\) con \((g(x),z)\in S\). Por ser $g$ morfismo, existe \(y\in X\) con \((x,y)\in R\) y \(g(y)\le z\). Como \(z\in U^c\) y \(U^c\) es decreciente, \(g(y)\in U^c\), luego \(y\in R(x)\cap g^{-1}[U]^{c}\neq\emptyset\) y \(x\in m_{R}(g^{-1}[U])\).

  \item[$\supseteq$)] Si \(x\in m_{R}(g^{-1}[U])\), entonces existe \(y\in X\) con \((x,y)\in R\) y \(g(y)\notin U\). Por homomorfismo \((g(x),g(y))\in S\), así \(S(g(x))\cap U^{c}\neq\emptyset\) y \(g(x)\in m_{S}(U)\). Luego \(x\in g^{-1}[m_{S}(U)]\).
\end{enumerate}
\end{itemize}

En ambos casos se cumple \(g^{-1}[m_{S}(U)]=m_{R}(g^{-1}[U])\), por lo que \(\D(g)\) es homomorfismo de \((i,j)\)-retículos.
\end{proof}
Al terminar vemos que tenemos a disposición un procedimiento sistemático para obtener $(i,j)$‑retículos a partir de $(i,j)$-espacios de Priestley y lo mismo con sus respectivos morfismos.


\subsection{La dualidad}

Con ambos funtores —de $(i,j)$-retículos a $(i,j)$-espacios de Priestley y viceversa— correctamente definidos en objetos y en flechas, estamos listos para articular la dualidad. Aquí presentamos las transformaciones naturales que enlazan los funtores compuestos con la identidad, un paso crucial cuya validez garantiza la equivalencia completa entre las categorías algebraica y topológica. Esta sección no repite las construcciones previas, sino que las culmina mostrando su compatibilidad global.\\


Recordemos que, según la Definición~\ref{def:transf-nat},  
una transformación natural \(\theta\colon F\Rightarrow G\) entre dos funtores
\(F,G\colon\mathsf C\to\mathsf D\)
es una asignación que a cada objeto \(X\) de \(\mathsf C\) le asocia un morfismo
\[
\theta_X\colon F(X)\longrightarrow G(X),
\]
y que satisface que para todo morfismo \(f\colon X\to Y\) en \(\mathsf C\) se cumple
\[
G(f)\circ \theta_X \;=\; \theta_Y \circ F(f).
\]


En particular, recordemos que $\eta\colon F\Rightarrow G$ forma parte de un isomorfismo natural si para cada objeto $A$ de $\mathsf{A}$, la flecha $\eta_{A}$ es un isomorfismo en $\mathsf B$. Así, dado $(\mathbb{L},m)$ un objeto de la categoría $\BDL(i,j)$, vamos a comprobar que la aplicación de Stone
\[ \sigma_{\mathbb{L}} \colon \mathbb{L} \rightarrow (\D \circ \mathfrak{X})(\mathbb{L}), \]
definida por $\sigma_{\mathbb{L}}(a) = \{ P \in \mathcal{X}(\mathbb{L}) \colon a \in P \}$ es un isomorfismo natural entre los funtores covariantes $\Id_{\BDL(i,j)}$ y  $\D \circ \mathfrak{X}$, de la misma manera que en la dualidad de Priestley original.

\begin{lemma}\label{Lema de Equivalencia Natural para BDL} Para todo $(i,j)$-retículo $(\mathbb{L}, m)$, la aplicación
$\sigma_{\mathbb{L}} \colon \mathbb{L} \rightarrow (\D \circ \mathfrak{X})(\mathbb{L})$ es un isomorfismo natural en $\BDL(i,j)$.
\end{lemma}
\begin{proof}[\textbf{Demostración}]
Para todo $(i,j)$-retículo $(\mathbb{L}, m)$, aplicando el funtor $\mathfrak{X}$ obtenemos
\begin{align*}
    \mathfrak{X}(\mathbb{L}, m)=(\mathcal{X}(\mathbb{L}), \tau_L, \subseteq, R_m)
\end{align*}
El Lema \ref{Lema de Representacion de BDL(i,j)} nos garantiza que $\mathfrak{X}(\mathbb{L}, m)$ es un $(i,j)$-espacio de Priestley, y el Lema \ref{Lema de Functorialidad de BDL(i,j)} que el funtor $\mathcal{X}$ respeta morfismos. Aplicando el funtor $\D$ obtenemos
\begin{align*}
    \D(\mathfrak{X}(\mathbb{L}, m))=(\mathcal{D}(\mathcal{X}(\mathbb{L})),\cap,\cup,\emptyset,\mathcal{X}(\mathbb{L}),m_{R_m})
\end{align*}
y los lemas \ref{Lema de Reconstrucción} y \ref{Lema de Functorialidad para D} nos aseguran que este último es un $(i,j)$-retículo y que el funtor $\D$ respeta morfismos.

Por la dualidad de Priestley, sabemos que para cada retículo distributivo acotado $\mathbb{L}$, la función de Stone $\sigma_{\mathbb{L}}$ es un isomorfismo de retículos. Para mostrar que $\sigma_{\mathbb{L}}$ es un isomorfismo en $\BDL(i,j)$ bastará con probar que conmuta con los operadores, es decir,
\[ \sigma_{\mathbb{L}}(m(a))=m_{R_m}(\sigma_{\mathbb{L}}(a)). \]
Sin embargo, en el Lema \ref{Lema de Representacion de BDL(i,j)} hemos comprobado que para cada $U = \sigma_{\mathbb{L}}(a)$, se da la igualdad 
\[ \sigma_{\mathbb{L}}(m(a)) = R_{m}^{-1}\left( \sigma_{\mathbb{L}}(a)^{i} \right)^{j} = m_{R_{m}}(\sigma_{\mathbb{L}}(a)) \]
Por lo tanto, $\sigma_{\mathbb{L}}$ resulta un isomorfismo en $\BDL(i,j)$. La naturalidad se sigue de la probada en la dualidad de Priestley original. El siguiente diagrama muestra la compatibilidad natural de $\sigma_{\mathbb{L}}$ con los funtores:

\[
\begin{tikzcd}[row sep=large, column sep=large]
  (L, \vee, \wedge, 0, 1, m)
    \arrow[r, "\mathfrak{X}" above]
    \arrow[dr, "\sigma_{\mathbb{L}}" left]
  & (\mathcal{X}(\mathbb{L}), \tau_L, \subseteq, R_m)
    \arrow[d, "\D" right] \\
  & (\D(\mathcal{X}(\mathbb{L})), \cup, \cap, \emptyset, \mathcal{X}(\mathbb{L}), m_{R_m})
\end{tikzcd}
\]

\end{proof}

Recordemos que, para la dualidad de tipo Priestley para retículos distributivos acotados, existe un isomorfismo natural $\epsilon \colon \Id_{\PS(i,j)} \Rightarrow \mathfrak{X} \circ \D$ definido de la siguiente manera: para cada $(X,\tau,\leq)$ objeto en $\PS$, la aplicación $\epsilon_{X} \colon X \to (\mathfrak{X}\circ \D)(X)$, definida por
\[ \epsilon_{X}(x) = \{ U \in \mathcal{D}(X) \colon x \in U \} \]
es un isomorfismo en $\PS$. Nuestro objetivo será comprobar que, para cada $(i,j)$-espacio, la aplicación $\epsilon_{X}$ es un isomorfismo en $\PS(i,j)$ y, en consecuencia, $\epsilon \colon \Id_{\PS(i,j)} \Rightarrow (\mathfrak{X}\circ \D)$ es un isomorfismo natural.


\begin{lemma}\label{Lema de Equivalencia Natural para PS} Para todo $(X, \tau, \leq, R)$ objeto de $\PS(i,j)$, la aplicación
$\epsilon_{X}$ es un isomorfismo natural en $\PS(i,j)$.
\end{lemma}

\begin{proof}[{\textbf{Demostración:}}]Para todo $(i,j)$-espacio de Priestley $(X,\tau,\leq,R)$, si aplicamos el funtor $\D$ obtenemos la estructura
\[ \D(X)=(\mathcal{D}(X),\cap,\cup,\emptyset,X,m_R). \]
Los lemas \ref{Lema de Reconstrucción} y \ref{Lema de Functorialidad para D} nos aseguran que este último es un $(i,j)$-retículo y que el funtor $\D$ respeta morfismos. Luego, aplicando el funtor $\mathfrak{X}$ obtenemos
\[ \mathfrak{X}(\D(X))=(\mathcal{X}(\mathcal{D}(X)),\tau_{\D(X)},\subseteq,R_{m_R}) \]
El lema \ref{Lema de Representacion de BDL(i,j)} nos garantiza que se trata de un $(i,j)$-espacio de Priestley, y el lema \ref{Lema de Functorialidad de BDL(i,j)} que el funtor $\mathfrak{X}$ respeta morfismos.
Teniendo en cuenta que la aplicación
\begin{align*}
\varepsilon_X&: X \to \mathcal{X}(\mathcal{D}(X))\\
\varepsilon_X&(x)=\{U\in \mathcal{D}(X): x\in U\} \end{align*}
es un isomorfismo natural de espacios de Priestley, comprobaremos además que es una flecha en $\PS(i,j)$, o lo que es equivalente, preserva la estructura de relación $R$, es decir que 
\begin{align*}
    (x,y) \in R\Leftrightarrow(\varepsilon_X(x), \varepsilon_X(y))\in R_{m_R}
\end{align*}
Probaremos  esta doble implicación para los casos $(1,1)$ y $(-1,1)$.

\begin{itemize}
    \item Caso $(1,1)$
\begin{itemize}
\item[$\Rightarrow$)] Sea $(x,y)\in R$, queremos ver que $(\varepsilon_X(x), \varepsilon_X(y))\in R_{m_R}$, que es lo mismo que 
\begin{align*}
    \varepsilon_X(y)^i\subseteq m_R^{-1}[\varepsilon_X(x)]^j
\end{align*}
Sea $U\in \varepsilon_X(y)$ arbitrario. Entonces $y\in U$, y puesto que $(x,y)\in R$
\begin{align*}
    R(x)\cap U\neq\emptyset &\implies x\in R^{-1}(U)\\
    &\implies R^{-1}(U)\in \varepsilon_X(x)\\
    &\implies m_R(U)\in \varepsilon_X(x)
\end{align*}
Por lo que $U\in m_R^{-1}[\varepsilon_X(x)]$, lo que nos dice que $\varepsilon_X(y)\subseteq m_R^{-1}[\varepsilon_X(x)]$ y por lo tanto
\begin{align*}
    (\varepsilon_X(x), \varepsilon_X(y))\in R_{m_R}
\end{align*}

\item[$\Leftarrow$)] Sea $(\varepsilon_X(x), \varepsilon_X(y))\in R_{m_R}$, queremos ver que $(x,y)\in R$. Probaremos la contrarrecíproca. Supongamos que \((x,y)\notin R\). Por la condición de separación de la definición de $(i,j)$‑relación, existe un $U\in \mathcal{D}(X)$ tal que $y\in U$ y $R(x)\cap U =\emptyset$. Esto implica que 
\begin{align*}
    U\in \varepsilon_X(y)\quad \text{ y }\quad x\notin R^{-1}(U)=m_R(U)
\end{align*}
Así $m_R(U)\notin \varepsilon_X(x)$ y por lo tanto $U\notin m_R^{-1}[\varepsilon_X(x)]$, pero sí $U\in \varepsilon_X(y)$. De esta manera $\varepsilon_X(y)\nsubseteq m_R^{-1}[\varepsilon_X(x)]$ por lo que 
\begin{align*}
    (\varepsilon_X(x), \varepsilon_X(y))\notin R_{m_R}
\end{align*}
\end{itemize}
\item Caso $(-1,1)$

\item[$\Rightarrow$)] Supongamos $(x,y)\in R$. 
queremos ver que $(\varepsilon_X(x), \varepsilon_X(y))\in R_{m_R}$, que es lo mismo que 
\begin{align*}
    \varepsilon_X(y)^c\subseteq m_R^{-1}[\varepsilon_X(x)]
\end{align*}
Sea $U\in \varepsilon_X(y)^c$ arbitrario. Entonces $y\notin U$, y puesto que $(x,y)\in R$
\begin{align*}
    R(x)\cap U^c\neq\emptyset &\implies x\in R^{-1}(U^c)\\
    &\implies R^{-1}(U^c)\in \varepsilon_X(x)\\
    &\implies m_R(U)\in \varepsilon_X(x)
\end{align*}
Por lo que $U\in m_R^{-1}[\varepsilon_X(x)]$, lo que nos dice que $\varepsilon_X(y)^c\subseteq m_R^{-1}[\varepsilon_X(x)]$ y por lo tanto
\begin{align*}
    (\varepsilon_X(x), \varepsilon_X(y))\in R_{m_R}
\end{align*}

\item[$\Leftarrow$)] Supongamos $(\varepsilon_X(x),\varepsilon_X(y))\in R_{m_R}$, es decir
\begin{align*}
\varepsilon_X(y)^{c}\subseteq m_R^{-1}[\varepsilon_X(x)]
\end{align*}
Queremos ver que $(x,y)\in R$. Probaremos la contrarrecíproca. Supongamos que $(x,y)\notin R$. Entonces existe $U\in \mathcal{D}(X)$ tal que
\begin{align*}
y\in U^{c} \quad\text{y}\quad R(x)\cap U^{c}=\emptyset.
\end{align*}
Así tenemos que $y\notin U$, es decir, $U\in \varepsilon_X(y)^{c}$. Por otro lado, $R(x)\cap U^c=\emptyset$ implica que $x\notin R^{-1}(U^c)$, es decir, $U\notin m_R^{-1}[\varepsilon_X(x)]$. Por lo que
\begin{align*}
\varepsilon_X(y)^{c}\nsubseteq m_R^{-1}[\varepsilon_X(x)]
\end{align*}
\end{itemize}
Se concluye que
\begin{align*}
    (x,y)\in R\Leftrightarrow(\varepsilon_X(x), \varepsilon_X(y))\in R_{m_R}
\end{align*}
y $\varepsilon_X$ es un isomorfismo de $(i,j)$-espacios de Priestley. La naturalidad se sigue de la probada en la dualidad de Priestley original. En el siguiente diagrama podemos ver la compatibilidad natural de $\varepsilon_X$ con los funtores:
\[
\begin{tikzcd}[row sep=large, column sep=huge]
  (X,\tau,\le,R)
    \arrow[r, "\D" above]
    \arrow[dr, "\varepsilon_X" below right]
  & (\mathcal{D}(X),\cup,\cap,\emptyset,X,m_R)
    \arrow[d, "\mathfrak{X}" right]
  \\
  & (\mathcal{X}(\mathcal{D}(X)),\tau_{\D(X)},\subseteq,R_{m_R})
\end{tikzcd}
\]

\end{proof}

\begin{theorem} \label{Dualidad_BDL_PS}
La categoría $\BDL(i,j)$ es dual a $\PS(i,j)$. Es decir, existe una equivalencia:
\[ \BDL(i,j)^{op} \cong \PS(i,j) \]
\end{theorem}

\begin{proof}[\textbf{Demostración}]
Por los lemas \ref{Lema de Equivalencia Natural para BDL} y \ref{Lema de Equivalencia Natural para PS} las transformaciones naturales
\[ \sigma_{\mathbb{L}}: \Id_{\BDL(i,j)}\Rightarrow (\D\circ \mathfrak{X}) \quad \text{y} \quad \varepsilon: \Id_{\PS(i,j)}\Rightarrow (\mathfrak{X}\circ \D) \]
son isomorfismos. Luego por la Definición \ref{def:equivalencia-dual} los funtores 
\begin{align*}
\mathfrak{X}:\BDL(i,j)\to \PS(i,j)^{op} \quad \text{ y }\quad \D:\PS(i,j)\to \BDL(i,j)^{op}
\end{align*}
son una equivalencia dual de las categorías de $(i,j)$-retículos y $(i,j)$-espacios de Priestley.
\end{proof}
De esta manera hemos probado la dualidad para los casos:

\[
\begin{tikzcd}
  \BDL(1,1)
    \arrow[r, shift left=0.7ex, "\mathfrak{X}"]
  & \PS(1,1)^{op}
    \arrow[l, shift left=0.7ex, "\D"]
\end{tikzcd}
\]
y 
\[
\begin{tikzcd}
  \BDL(-1,1)
    \arrow[r, shift left=0.7ex, "\mathfrak{X}"]
  & \PS(-1,1)^{op}
    \arrow[l, shift left=0.7ex, "\D"]
\end{tikzcd}
\]
Sin embargo, utilizando las equivalencias probadas en la sección anterior tenemos las cadenas de equivalencias
\[
\begin{tikzcd}
  \BDL(-1,-1)
    \arrow[r, shift left=0.7ex, "{(-)^*}"]
  & \BDL(1,1)
    \arrow[l, shift left=0.7ex, "{(-)_*}"]
    \arrow[r, shift left=0.7ex, "\mathfrak{X}"]
  & \PS(1,1)^{op}
    \arrow[l, shift left=0.7ex, "\D"]
    \arrow[r, shift left=0.7ex, "{(-)^{+}}"]
  & \PS(-1,-1)^{op}
    \arrow[l, shift left=0.7ex, "{(-)_{+}}"]
\end{tikzcd}
\]
y
\[
\begin{tikzcd}
  \PS(1,-1)^{op}
    \arrow[r, shift left=0.7ex, "{(-)^{+}}"]
  & \PS(-1,1)^{op}
    \arrow[l, shift left=0.7ex, "{(-)_{+}}"]
    \arrow[r, shift left=0.7ex, "\D"]
  & \BDL(-1,1)
    \arrow[l, shift left=0.7ex, "\mathfrak{X}"]
    \arrow[r, shift left=0.7ex, "{(-)^*}"]
  & \BDL(1,-1)
    \arrow[l, shift left=0.7ex, "{(-)_*}"]
\end{tikzcd}
\]
Con lo cuál hemos probado lo siguiente:
\begin{corollary}
Para todo $(i,j)\in \{1,-1\}^{2}$ las categorías $BDL(i,j)$ y
$PS(i,j)$ son dualmente equivalentes, es decir
\[ \BDL(i,j)^{op} \cong \PS(i,j). \]
\end{corollary}













\chapter{Extensión canónica} \label{Ch 4}


La teoría de extensiones canónicas surgió a finales del siglo XX en el estudio de dualidades algebraico‑topológicas, cuando investigadores como Jónsson, Tarski, Gehrke y Harding observaron que el paso directo del álgebra —por ejemplo, retículos distributivos con operadores— a un contexto topológico implicaba a menudo construcciones ad hoc para relacionar operaciones algebraicas con relaciones en el espacio dual. En particular, para operadores modales o normales, la relación que hemos utilizado en la dualidad de tipo Priestley desarrollada en el Capítulo \ref{Ch 3} de este trabajo para la interpretación de dichos operadores, se construyen de manera aparentemente diferente en cada caso, sin seguir una construcción unificada desde la estructura algebraica original.
\medskip


Teniendo en cuenta este problema, surge naturalmente la siguiente pregunta: ¿Existe alguna manera estándar de establecer este tipo de relaciones binarias?. La \emph{extensión canónica} de un retículo distributivo acotado, introducida originalmente por Jónsson y Tarski en el contexto de álgebras booleanas con operadores y posteriormente generalizada a retículos por Gehrke y Harding \cite{GehrkeHarding2001}, ofrece una respuesta elegante a este problema. 
\medskip

En este capítulo, estudiaremos la construcción formal de la $\sigma$‑extensión de un retículo distributivo acotado y sus propiedades universales. Asimismo, estudiaremos la manera de extender operadores unarios mediante dichas extensiones. En particular, buscaremos establecer conexiones entre las extensiones canónicas para operadores unarios y la dualidad tipo Priestley ya determinada en este trabajo. 




\section{Definiciones elementales} \label{S 4.1}

En la presente sección, introduciremos la notación y las nociones elementales necesarias para la construcción de la extensión canónica para retículos distributivos acotados. El objetivo de dicha construcción algebraica es inyectar un retículo distributivo acotado dado dentro de un retículo completo. En particular, dicha construcción puede especializarse hacia álgebras con un lenguaje más complejo, especializando de manera adecuada dicha construcción.  



Comenzaremos esta sección definiendo lo que, en este trabajo, entenderemos como extensión algebraica, noción que es central en diversas áreas de la matemática, ya que permite ampliar la estructura de un objeto algebraico manteniendo sus propiedades esenciales .\\


\begin{definition} Sea $\mathbf{A} = (A,\mathcal{F})$ un álgebra de tipo $\mathcal{F}$. Una \emph{extensión algebraica} de $\mathbf{A}$ es un par $(\mathbf{B}, \iota)$, donde $\mathbf{B}$ es un álgebra de tipo $\mathcal{F}$  y
$$
\iota: A \to B
$$
es un homomorfismo inyectivo. 
\end{definition}

Teniendo en cuenta esta noción de extensión algebraica, buscaremos establecer extensiones algebraicas para retículos distributivos acotados. En particular, dado $\mathbb{L} = (L,\wedge,\vee,0,1)$ un retículo distributivo acotado, buscaremos establecer un par $\left( \mathbb{L}^{\delta},\iota \right)$ \footnote{Si no hay ambigüedad, escribiremos $\mathbb{L}^{\delta}$ en lugar de $\left( \mathbb{L}^{\delta},\iota \right)$}, donde $\iota$ es un homomorfismo inyectivo y además $\mathbb{L}^{\delta}$ es un retículo que satisface ciertas propiedades que mencionaremos a continuación.\\



\begin{definition} \label{reticulo completo} Sea $\mathbb{L}$ un retículo. Diremos que $\mathbb{L}$ es \emph{completo} si para todo subconjunto $S \subseteq L$, las operaciones de supremo arbitrario $\bigvee S$ e ínfimo arbitrario $\bigwedge S$ están definidas en $\mathbb{L}$.
\end{definition}

Es claro que todo retículo finito resulta un retículo completo. Sin embargo, no todo retículo resulta completo como se muestra en el siguiente ejemplo.

\begin{example} Dado $(X,\tau)$ un espacio topológico. El álgebra $\left( \mathcal{O}(X),\wedge,\vee,\emptyset,X \right)$, cuyo universo son los conjuntos abiertos de $(X,\tau)$ y las operaciones $\wedge$ y $\vee$ están determinadas por la intersección y la unión de conjuntos respectivamente, es un retículo distributivo acotado. Sin embargo, no es completo. Ya que la intersección arbitraria de conjuntos abiertos, no necesariamente es un abierto.
\end{example}

\begin{definition}\label{def:compactos-cocompactos}
Sea $\mathbb{L}$ un retículo completo. Un elemento $c \in L$ se dice \emph{compacto algebraico} si para toda familia $S \subseteq L$ que verifique 
\[
c = \bigvee S,
\]
existe un subconjunto finito $S_0 \subseteq S$ tal que
\[
c = \bigvee S_0.
\]
Denotamos por $\mathcal{K}(\mathbb{L})$ al conjunto de elementos compactos algebraicos de $\mathbb{L}$.

\noindent Análogamente, un elemento $d\in L$ se dice \emph{co-compacto algebraico} (o \emph{compacto algebraico en el orden dual}) si para toda familia $T\subseteq L$ que verifique
\[
d = \bigwedge T,
\]
existe un subconjunto finito $T_0\subseteq T$ tal que
\[
d = \bigwedge T_0.
\]
Denotamos por $\mathcal{K}(\mathbb{L}^{-1})$ al conjunto de elementos co-compactos algebraicos de $\mathbb{L}$ (es decir, los compactos algebraicos de $\mathbb{L}^{-1}$).
\end{definition}
En otras palabras: $\mathcal{K}(\mathbb L)$ son los elementos que son finitamente generados por \(\vee\) y
$\mathcal{K}(\mathbb L^{-1})$ son los elementos que son finitamente generados por \(\wedge\).\\

\begin{definition}\label{def:densidades-join-meet}
Sea $\mathbb{L}$ un retículo completo y $S\subseteq L$ un subconjunto.
Diremos que $S$ es \emph{$\vee$-denso} en $\mathbb{L}$ si para todo $u\in L$ se cumple:
\[
u \;=\; \bigvee\{\, s\in S : s\le u \,\}.
\]
Es decir, cada elemento de $\mathbb{L}$ es el supremo de todos los elementos de $S$ que están por debajo de él. Dualmente, diremos que $S$ es \emph{$\wedge$-denso} en $\mathbb{L}$ si para todo $u\in L$ se cumple:
    \[
    u \;=\; \bigwedge\{\, s\in S : u\le s \,\}.
    \]
    Es decir, cada elemento de $\mathbb{L}$ es el ínfimo de todos los elementos de $S$ que están por encima de él.
\end{definition}

\begin{definition}
Un retículo $\mathbb{L}$ se dice \emph{algebraico} si es completo y para todo elemento $x \in L$ se tiene
\[
x \;=\; \bigvee \{\, c \in \mathcal{K}(\mathbb{L}) : c \le x \,\},
\]
esto es, el conjunto $\mathcal{K}(\mathbb{L})$ de elementos compactos algebraicos es $\vee$-denso en $\mathbb{L}$.
\end{definition}

\begin{definition} Dualmente, diremos que $\mathbb{L}$ es \emph{dualmente algebraico} si $\mathbb{L}^{-1}$ es algebraico, es decir que para todo \(x\in L\):
\begin{align*}
    x&=\bigvee\nolimits_{\mathbb L^{-1}}\{\,d\in \mathcal{K}(\mathbb L^{-1}): d\le_{\mathbb L^{-1}} x\,\}\\
    &= \bigwedge\nolimits_{\mathbb L}\{\,d\in \mathcal{K}(\mathbb L^{-1}): d\ge_{\mathbb L} x\,\}
\end{align*}
En otras palabras, los \emph{co-compactos algebraicos} $\mathcal{K}(\mathbb{L}^{-1})$ son $\wedge$-densos en $\mathbb{L}$.
\end{definition}

\begin{definition}
Diremos que $\mathbb{L}$ es \emph{doblemente algebraico} si tanto $\mathbb{L}$ como su retículo dual $\mathbb{L}^{-1}$ son algebraicos. 
\end{definition}


\begin{definition} Sea $\mathbb{L}$ un retículo completo. Se dice que un elemento $u\in L$, con $u\neq 0$, es \emph{completamente $\vee$-irreducible} si para toda familia $S\subseteq L$ se cumple que
\begin{align*}
u = \bigvee S \quad \Longrightarrow \exists s\in S: s=u
\end{align*}
Denotamos por $\mathcal{J}^{\infty}(\mathbb{L})$ al conjunto de todos los elementos completamente $\vee$-irreducibles de $\mathbb{L}$.
\end{definition}

\begin{lemma}\label{thm:J-en-K} Sea $\mathbb{L}$ un retículo completo. Entonces,   
\begin{align*}
\mathcal{J}^{\infty}(\mathbb{L}) \subseteq \mathcal{K}(\mathbb{L}).
\end{align*}
\end{lemma}

\begin{proof}[\textbf{Demostración:}]
Sea $u\in\mathcal{J}^{\infty}(\mathbb{L})$. Supongamos que existe una familia
$S\subseteq L$ tal que
\[
u \;=\; \bigvee S.
\]
Por la definición de completamente $\vee$-irreducible, de la igualdad anterior
se deduce que existe $s\in S$ con $s=u$. Tomando el subconjunto finito $S_0=\{u\}$
obtenemos
\[
u \;=\; \bigvee S_0,
\]
por lo que
$u\in \mathcal{K}(\mathbb{L})$. Dado que la elección de $u\in\mathcal{J}^\infty(\mathbb{L})$
fue arbitraria, concluimos que
$\mathcal{J}^{\infty}(\mathbb{L}) \subseteq \mathcal{K}(\mathbb{L})$.
\end{proof}



\begin{definition} Sea $\mathbb{L}$ un retículo completo. Se dice que un elemento $v\in L$, con $v\neq 1$, es \emph{completamente $\wedge$-irreducible} si para toda familia $S\subseteq L$ se cumple que
\begin{align*}
v = \bigwedge S \quad \Longrightarrow \exists s\in S: s=v
\end{align*}
Denotamos por $\mathcal{M}^{\infty}(\mathbb{L})$ al conjunto de todos los elementos completamente $\wedge$-irreducibles de $\mathbb{L}$.
\end{definition}

\begin{lemma}\label{lem:meet-irr-implica-cocompact}
Sea $\mathbb L$ un retículo completo. Entonces
\[
\mathcal M^\infty(\mathbb L)\;\subseteq\; \mathcal{K}(\mathbb L^{-1}).
\]
Es decir, todo elemento completamente $\wedge$-irreducible es co-compacto algebraico.
\end{lemma}

\begin{proof}[\textbf{Demostración:}]
Sea $v\in\mathcal M^\infty(\mathbb L)$ y supongamos que existe una familia
$S\subseteq L$ tal que
\[
v=\bigwedge S
\]
Por la definición de completamente $\wedge$-irreducible, la igualdad anterior implica que existe algún $s\in S$ con $s=v$. Entonces el subconjunto finito $S_0=\{v\}$ satisface
\[v=\bigwedge S_0\]
y así $v$ es compacto algebraico en $\mathbb L^{-1}$, lo que demuestra la inclusión.
\end{proof}


\begin{definition} Sea $e\, \colon \mathbb{L}_1 \to \mathbb{L}_2$ un homomorfismo inyectivo de retículos, donde $\mathbb{L}_2$ es completo. Diremos que $e[L_1]$ es \emph{separante} en $\mathbb{L}_2$ si para todo $u\in \mathcal{J}^{\infty}(\mathbb{L}_2)$ y para todo $v\in \mathcal{M}^{\infty}(\mathbb{L}_2)$ con $v\le u$, existe un elemento $a\in L_1$ tal que
\begin{align*}
v \le e(a) \le u.
\end{align*}
\end{definition}


\begin{definition}
\label{def:compacto-en}
Sea $e\, \colon \mathbb{L}_1 \to \mathbb{L}_2$ un homomorfismo inyectivo de retículos, donde $\mathbb{L}_2$ es completo. Decimos que $e[L_1]$ es \emph{compacto} en $\mathbb{L}_2$ si
para cualesquiera $X,Y\subseteq L_1$ que verifiquen
\[
\bigwedge\nolimits_{\mathbb{L}_2} e[X] \;\leq\; \bigvee\nolimits_{\mathbb{L}_2} e[Y],
\]
existen subconjuntos finitos $F\subseteq X$ y $G\subseteq Y$ tales que
\[
\bigwedge\nolimits_{\mathbb{L}_2} e[F]
\;=\;
e\!\bigl(\bigwedge\nolimits_{\mathbb{L}_1} F\bigr)
\;\leq\;
e\!\bigl(\bigvee\nolimits_{\mathbb{L}_1} G\bigr)
\;=\;
\bigvee\nolimits_{\mathbb{L}_2} e[G].
\]
\end{definition}


Con las nociones introducidas a lo largo de este capitulo vamos a definir lo que, a lo largo de este trabajo, vamos a entender como extensión canónica.


\begin{definition}\label{Def:Ext canonica} Sea $\mathbb{L}$ un retículo distributivo acotado. Una extensión $(\mathbb{L}^\delta, e)$ de $\mathbb{L}$ se dirá  \emph{extensión canónica} si verifica las siguientes condiciones:
\begin{enumerate}
    \item $\mathbb{L}^\delta$ es completo y doblemente algebraico.
    \item $e[L]$ es separante en $\mathbb{L}^\delta$.
    \item $e[L]$ es compacto en $\mathbb{L}^\delta$.
\end{enumerate}
\end{definition}

El siguiente resultado sobre extensiones canónicas es de gran interés en este trabajo. En particular, afirma que una extensión canónica es esencialmente única, es decir, es única salvo isomorfismos. Esto nos permite entonces hablar esencialmente de una única extensión canónica para un retículo distributivo acotado dado (ver \cite[Proposición 2.7]{GehrkeHarding2001}).

\begin{theorem}\label{thm:unicidad-extension-canonica}
Sea $\mathbb{L}$ un retículo distributivo acotado y supongamos que
\[
  e_1\colon \mathbb{L}\;\hookrightarrow\;\mathbb{L}_1^\delta,
  \qquad
  e_2\colon \mathbb{L}\;\hookrightarrow\;\mathbb{L}_2^\delta
\]
son dos extensiones canónicas de $\mathbb{L}$. Entonces existe un único isomorfismo de retículos completos
\[
  \varphi\colon \mathbb{L}_1^\delta \;\xrightarrow{\;\cong\;}\; \mathbb{L}_2^\delta
\]
tal que $\varphi\circ e_1 = e_2$.
\end{theorem}

\begin{definition}\label{def:abiertos}
Sea $\mathbb{L}$ un retículo distributivo acotado y $(\mathbb{L}^\delta,e)$ su extensión canónica. Diremos que un elemento $x\in\mathbb{L}^\delta$ es un \emph{abierto algebraico} si existe un conjunto $S\subseteq L$ tal que
\[
x \;=\; \bigvee\nolimits_{\mathbb{L}^\delta}\{\,e(a):a\in S\,\}.
\]
Denotamos por $\mathcal{O}(\mathbb{L}^\delta)$ al conjunto de todos los elementos abiertos algebraicos de $\mathbb{L}^\delta$.
\end{definition}


\begin{definition}\label{def:cerrados} Análogamente, diremos que un elemento $x\in\mathbb{L}^\delta$ es un \emph{cerrado algebraico} si existe un conjunto $T\subseteq L$ tal que
\[x \;=\; \bigwedge\nolimits_{\mathbb{L}^\delta}\{\,e(b):b\in T\,\}\]
Denotamos por $\mathcal{C}(\mathbb{L}^\delta)$ al conjunto de todos los elementos cerrados algebraicos de $\mathbb{L}^\delta$.
\end{definition}

\begin{lemma}\label{lem:cerrados-op}
Sea $\mathbb{L}$ un ret\'{\i}culo distributivo acotado y $(\mathbb{L}^\delta,e)$ su extensi\'on can\'onica. 
Entonces, todo cerrado algebraico de $\mathbb{L}^\delta$ es un abierto algebraico de $(\mathbb{L}^\delta)^{-1}$. En consecuencia,
\[
\mathcal{C}(\mathbb{L}^\delta)
\;=\;
\mathcal{O}\!\left((\mathbb{L}^\delta)^{-1}\right).
\]
\end{lemma}
\begin{proof}[\textbf{Demostración:}]
Sea $C\in\mathcal{C}(\mathbb{L}^\delta)$. Por definici\'on, existe un conjunto $T\subseteq L$ tal que
\[
C \;=\; \bigwedge_{\mathbb{L}^\delta}\{\,e(b):b\in T\,\}.
\]
Al considerar $C$ como elemento de $(\mathbb{L}^\delta)^{-1}$, la operaci\'on $\bigwedge_{\mathbb{L}^\delta}$ se interpreta como 
$\bigvee_{(\mathbb{L}^\delta)^{-1}}$. Por lo tanto,
\[
C
\;=\;
\bigvee_{(\mathbb{L}^\delta)^{-1}}\{\,e(b):b\in T\,\},
\]
y por definici\'on esto muestra que $C$ es un abierto algebraico de $(\mathbb{L}^\delta)^{-1}$. 
La igualdad entre ambas familias de conjuntos se sigue inmediatamente.
\end{proof}


\begin{proposition}\label{prop:clopen-is-image}
Sean $\mathbb{L}$ un retículo distributivo acotado y $(\mathbb{L}^\delta,e)$ su extensión canónica.
Entonces
\[
e[L] \;=\; \mathcal{O}(\mathbb{L}^\delta)\ \cap\ \mathcal{C}(\mathbb{L}^\delta).
\]
Es decir, un elemento de $\mathbb{L}^\delta$ proviene de $\mathbb{L}$ si y sólo si es simultáneamente abierto algebraico y cerrado algebraico.
\end{proposition}

\begin{proof}[\textbf{Demostración:}]

Probaremos la doble inclusión.

\begin{itemize}
    \item $(\subseteq)$ Si $a\in L$ entonces
\[
e(a)=\bigvee\nolimits_{\mathbb{L}^\delta}\{e(a)\},
\qquad
e(a)=\bigwedge\nolimits_{\mathbb{L}^\delta}\{e(a)\},
\]
así que $e(a)$ es a la vez abierto y cerrado algebraico.

\item $(\supseteq)$ Sea $x\in\mathcal{O}(\mathbb{L}^\delta)\cap\mathcal{C}(\mathbb{L}^\delta)$.
Entonces existen $A,B\subseteq L$ tales que
\[
x \;=\; \bigvee\nolimits_{\mathbb{L}^\delta}\{e(a):a\in A\}
\qquad\text{y}\qquad
x \;=\; \bigwedge\nolimits_{\mathbb{L}^\delta}\{e(b):b\in B\}.
\]
De estas igualdades se obtiene la desigualdad
\[
\bigwedge\nolimits_{\mathbb{L}^\delta} e[B]
\;\le\;
\bigvee\nolimits_{\mathbb{L}^\delta} e[A].
\]

Como $(\mathbb{L}^\delta,e)$ es una extensión canónica, la imagen $e[L]$ es \emph{compacta} en $\mathbb{L}^\delta$ y existen subconjuntos finitos $F\subseteq B$ y $G\subseteq A$ tales que
\[
\bigwedge\nolimits_{\mathbb{L}^\delta}\{e(b):b\in F\}
\;\le\;
\bigvee\nolimits_{\mathbb{L}^\delta}\{e(a):a\in G\}.
\]
Como $e$ es homomorfismo de retículos, las operaciones finitas conmutan con $e$, luego
\[
e\!\bigl(\bigwedge_{b\in F} b\bigr)
\;\le\;
e\!\bigl(\bigvee_{a\in G} a\bigr).
\]
La inyectividad de $e$ implica en $\mathbb{L}$ que
\[
\bigwedge_{b\in F} b \;\le\; \bigvee_{a\in G} a.
\tag{$(1)$}
\]

Por la elección de $F,G$ y las expresiones de $x$ como ínfimo y supremo:
\[
x \;\le\; e\!\bigl(\bigwedge_{b\in F} b\bigr)
\quad\text{y}\quad
e\!\bigl(\bigvee_{a\in G} a\bigr)\le x.
\]
Combinando estas desigualdades con la desigualdad $(1)$ se obtiene
\[
x \;\le\; e\!\bigl(\bigwedge_{b\in F} b\bigr)
\;\le\;
e\!\bigl(\bigvee_{a\in G} a\bigr)
\;\le\; x,
\]
de donde sigue la igualdad
\[
x \;=\; e\!\bigl(\bigvee_{a\in G} a\bigr).
\]
Por tanto $x\in e[L]$, y hemos mostrado
$\mathcal{O}(\mathbb{L}^\delta)\cap\mathcal{C}(\mathbb{L}^\delta)\subseteq e[L]$.
\end{itemize}
Así, la afirmación se sigue. 
\end{proof}



En adelante, teniendo en cuenta el Teorema \ref{thm:unicidad-extension-canonica}, hablaremos de \emph{la} extensión canónica entendiendo que es esencialmente única. Además,  escribiremos \(\mathbb L^\delta\) para indicar la extensión canónica en lugar de \((\mathbb L^\delta,e)\).


Quedan ahora sí dadas las herramientas algebraicas necesarias para definir una extensión canónica particular, la \(\sigma\)-extensión de \(\mathbb{L}\), cuya construcción y propiedades demostraremos en la siguiente sección.





\section{La extensión canónica para retículos distributivos acotados}\label{S4.2}


Tal como hemos mencionado en la sección anterior, para cada retículo $\mathbb{L}$, la extensión canónica es esencialmente única i.e., todas las construcciones algebraicas que extiendan canónicamente a $\mathbb{L}$, en el sentido de la Definición \ref{Def:Ext canonica}, son isomorfas. En particular, en la presente sección, vamos a construir la extensión canónica en la cual nos centraremos en este trabajo, entendiendo que cualquier otra construcción será isomorfa a la que presentada. La misma, está basada en la aplicación de Stone:
\begin{align*}
  \sigma_{\mathbb{L}}&\colon \mathbb{L}\rightarrow\bigl(\mathcal{D}(\mathcal{X}(\mathbb{L})),\,\cup,\,\cap,\,\emptyset,\,\mathcal{X}(\mathbb{L})\bigr)
\end{align*}
Recordemos que, por los resultados obtenidos en la dualidad de Priestley, si $\mathbb{L}$ es un retículo distributivo acotado entonces la aplicación de Stone es un isomorfismo de retículos distributivos acotados. En particular, puede verse como un homomorfismo inyectivo de retículos sobre el retículo
\begin{align*}
  \Up(\mathcal{X}(\mathbb{L}))=\bigl(\mathrm{Up}(\mathcal{X}(\mathbb{L})),\,\cup,\,\cap,\,\emptyset,\,\mathcal{X}(\mathbb{L})\bigr)
\end{align*}
ya que todo clopen creciente es claramente creciente. Más aún, $\Up(\mathcal{X}(\mathbb{L}))$ es un retículo completo y por lo tanto, es un candidato natural a ser una extensión canónica.

El objetivo de este apartado será comprobar que el par $\left(\Up(\mathcal{X}(\mathbb{L})),\sigma_{\mathbb{L}} \right)$ es, en efecto, una extensión que extiende canónicamente a $\mathbb{L}$ y, por lo tanto, la entenderemos como la extensión canónica para $\mathbb{L}$. Previo a este resultado, probaremos una serie de lemas que nos serán útiles para caracterizan a los elementos completamente $\vee$-irreducibles y completamente $\wedge$-irreducibles de $\Up(\mathcal{X}(\mathbb{L}))$.
 


\begin{lemma}\label{lem:irreducibles-principales}
En el retículo completo $\Up(\mathcal{X}(\mathbb L))$,
los elementos completamente $\vee$-irreducibles son exactamente los crecientes principales:
\[
\mathcal J^\infty\bigl(\Up(\mathcal{X}(\mathbb L))\bigr)
\;=\;
\{\,\up{P}: P\in\mathcal{X}(\mathbb L)\,\}.
\]
\end{lemma}

\begin{proof}[\textbf{Demostración:}]
Probaremos la doble contención:
\begin{itemize}
    \item (\(\subseteq\)) Sea $U\in\mathcal J^\infty(\Up(\mathcal{X}(\mathbb L)))$.
Como $U$ es creciente,
\[
U \;=\; \bigcup_{Q\in U} \up{Q},
\]
puesto que, para todo $Q \in U$ se tiene que $\up{Q}\subseteq U$. Por la completa $\vee$-irreducibilidad de $U$,
esta unión sólo es posible si existe $Q\in U$ tal que $U=\up{Q}$.

\item (\(\supseteq\)) Sea $P\in\mathcal{X}(\mathbb L)$. Supongamos
\[
\up{P} \;=\; \bigcup_{i\in I} U_i
\quad\text{con } U_i\in \Up(\mathcal{X}(\mathbb L)).
\]
Como $P\in\up{P}$, existe $i_0$ con $P\in U_{i_0}$. Al ser $U_{i_0}$ creciente y contener a $P$,
se tiene $\up{P}\subseteq U_{i_0}$. Por la igualdad de uniones, también $U_{i_0}\subseteq \up{P}$,
luego $U_{i_0}=\up{P}$. Esto prueba que $\up{P}$ es completamente $\vee$-irreducible.
\end{itemize}
\end{proof}

\begin{lemma}\label{lem:meet-irreducibles-complementos}
En el retículo completo $\Up(\mathcal{X}(\mathbb L))$, los elementos completamente $\wedge$-irreducibles son exactamente los complementos de decrecientes principales. Es decir,
\[
\mathcal M^\infty\bigl(\Up(\mathcal{X}(\mathbb L))\bigr)
\;=\;
\bigl\{\,\downc{P}\;:\;P\in\mathcal{X}(\mathbb L)\,\bigr\}.
\]
\end{lemma}

\begin{proof}[\textbf{Demostración:}]
Probaremos la doble contención:

\begin{itemize}
\item (\(\subseteq\)) 
Sea $V\in\mathcal M^\infty(\Up(\mathcal{X}(\mathbb L)))$.  
Como $V$ es creciente, para cada $Q\notin V$ el conjunto decreciente $\down{Q}$ es disjunto de $V$, y se cumple la representación
\[
V \;=\; \bigcap_{Q\notin V}\downc{Q}.
\]
Por la completa $\wedge$-irreducibilidad de $V$, existe $P\notin V$ tal que $V=\downc{P}$.

\noindent\textbf{(\(\supseteq\))}  
Sea \(P\in\mathcal X(\mathbb L)\) y definamos \(V := (P]^{c}=\downc{P}\).  
Supongamos que
\[
V=\bigcap_{i\in I} U_i,
\qquad U_i\in\Up(\mathcal X(\mathbb L)).
\]
Como \(P\notin V\), se tiene
\[
P\in V^{c}
= \Bigl(\bigcap_{i\in I}U_i\Bigr)^{c}
= \bigcup_{i\in I} U_i^{c},
\]
por lo que existe \(i_0\in I\) tal que  
\(P\in U_{i_0}^{c}\), es decir, \(P\notin U_{i_0}\). Vamos a comprobar que \(V = U_{i_0}\), o lo que es lo mismo, que \((P]^{c}=U_{i_0}\).

\begin{itemize}
\item \(U_{i_0}\subseteq (P]^{c}\)

Sea \(Q\in U_{i_0}\).  
Si \(Q\in (P]\), entonces \(Q\subseteq P\).  
Dado que \(U_{i_0}\) es creciente, de \(Q\in U_{i_0}\) y \(Q\subseteq P\) se seguiría \(P\in U_{i_0}\), contradicción con \(P\notin U_{i_0}\).  
Por lo tanto, \(Q\notin (P]\), es decir,
\[
Q\in (P]^{c}.
\]
Así,
\[
U_{i_0}\subseteq (P]^{c}.
\]

\item \((P]^{c}\subseteq U_{i_0}\)

Sea \(Q\in (P]^{c}\).  
Entonces \(Q\notin (P]\); en particular, \(Q\in V=\bigcap_{i\in I}U_i\).  
Por tanto,
\[
Q\in U_{i_0}.
\]
Así,
\[
(P]^{c}\subseteq U_{i_0}.
\]
\end{itemize}

De (1) y (2) obtenemos la igualdad
\[
(P]^{c}=U_{i_0}.
\]

Por lo tanto, \(\downc{P}\) es completamente $\wedge$-irreducible.
\end{itemize}
Por lo que queda probada la doble inclusión
\end{proof}


\begin{remark}\label{rem:join-meet-densos}
De las demostraciones de los lemas~\ref{lem:irreducibles-principales} y~\ref{lem:meet-irreducibles-complementos}
obtenemos las descomposiciones canónicas siguientes para todo
$U\in\Up(\mathcal{X}(\mathbb{L}))$:
\[
U \;=\; \bigvee\{\, [P) : P\in U \,\},
\qquad
U \;=\; \bigwedge\{\, (Q]^c : Q\notin U \,\}.
\]
Se sigue entonces que el conjunto 
\(\mathcal J^\infty\bigl(\Up(\mathcal{X}(\mathbb L))\bigr)=\{[P):P\in\mathcal{X}(\mathbb L)\}\)
es $\vee$-denso. Dualmente, $\mathcal M^\infty\bigl(\Up(\mathcal{X}(\mathbb L))\bigr)
\;=\;
\bigl\{\,\downc{P}\;:\;P\in\mathcal{X}(\mathbb L)\,\bigr\}$ es $\wedge$-denso.
\end{remark}

\begin{lemma}\label{lem:densidad-heredada}
Sea $\mathbb L$ un retículo completo y $S\subseteq T\subseteq L$. 
Si $S$ es $\vee$-denso en $\mathbb L$, entonces también $T$ es $\vee$-denso en $\mathbb L$. 
De manera dual, si $S$ es $\wedge$-denso en $\mathbb L$, entonces también lo es $T$. 
\end{lemma}

\begin{proof}[\textbf{Demostración:}]
Supongamos que $S$ es $\vee$-denso. Entonces, para todo $u\in L$,
\[
u \;=\; \bigvee\{\,s\in S : s\le u\,\}.
\]
Pero como $\{\,s\in S : s\le u\,\}\subseteq \{\,t\in T : t\le u\,\}$, por
monotonía del operador $\bigvee$ se tiene
\[
\bigvee\{\,t\in T : t\le u\,\}
\;\;\ge\;\;
\bigvee\{\,s\in S : s\le u\,\}
\;=\; u.
\]
Por otra parte, todos los elementos de $\{t\in T : t\le u\}$ son menores o iguales a $u$, luego
\[
\bigvee\{\,t\in T : t\le u\,\}\;\le\; u.
\]
De ambas desigualdades concluimos
\[
u \;=\; \bigvee\{\,t\in T : t\le u\,\}.
\]
Así $T$ es $\vee$-denso en $\mathbb L$. El caso dual se prueba de manera análoga.
\end{proof}


Ahora si, procedemos con la extensión canónica:
\begin{theorem}\label{teo:ext-canonica-sigma}
Sea $\mathbb{L}$ un retículo distributivo acotado. Entonces, el par $\left(\Up(\mathcal{X}(\mathbb{L})),\sigma_{\mathbb{L}} \right)$, donde
\[ \Up(\mathcal{X}(\mathbb{L})) = \left(\mathrm{Up}(\mathcal X(\mathbb L)),\cap,\cup,\mathcal X(\mathbb L),\emptyset \right), \]
y $\sigma_{\mathbb{L}}$ es la aplicación de Stone definida por  
\begin{align*}
  \sigma_{\mathbb{L}}&\colon \mathbb{L}\hookrightarrow\Up(\mathcal{X}(\mathbb{L}))\\
  \sigma_{\mathbb{L}}&(a)=\{\,P\in\mathcal{X}(\mathbb{L}):a\in P\},
\end{align*}
extiende canónicamente a \(\mathbb{L}\).
\end{theorem}

\begin{proof}[\textbf{Demostración:}]
Ya se observó que $\Up(\mathcal{X}(\mathbb L))$ es un retículo completo. Mostraremos entonces que el par $\left(\Up(\mathcal{X}(\mathbb{L})),\sigma_{\mathbb{L}} \right)$ verifica las tres condiciones de la definición \ref{Def:Ext canonica}.

\begin{enumerate}
    \item Doblemente algebraico.
    
Teniendo en cuenta la observación~\ref{rem:join-meet-densos}, los elementos completamente $\vee$-irreducibles son $\vee$-densos y los completamente $\wedge$-irreducibles son $\wedge$-densos. Además, por Teorema~\ref{thm:J-en-K} se tiene que
$\mathcal J^\infty(\Up(\mathcal{X}(\mathbb L)))\subseteq\mathcal K(\Up(\mathcal{X}(\mathbb L)))$  y $
\mathcal M^\infty(\Up(\mathcal{X}(\mathbb L)))\;\subseteq\; \mathcal{K}(\Up(\mathcal{X}(\mathbb L))^{-1})$. Se sigue entonces del Lema \ref{lem:densidad-heredada} que los compactos algebraicos son $\vee$-densos y los co-compactos algebraicos son $\wedge$-densos;
en consecuencia $\Up(\mathcal{X}(\mathbb L))$ es doblemente algebraico.


\item $\sigma_{\mathbb{L}}(\mathbb L)$ es separante en $\Up(\mathcal{X}(\mathbb L))$.

Sean \(U\in\mathcal J^\infty(\Up(\mathcal{X}(\mathbb L)))\) y 
\(V\in\mathcal M^\infty(\Up(\mathcal{X}(\mathbb L)))\) con \(V \subseteq U\).
Por los Lemas~\ref{lem:irreducibles-principales} y
\ref{lem:meet-irreducibles-complementos} existen
\(P,Q\in\mathcal X(\mathbb L)\) tales que
\[
U=[P)
\qquad\text{y}\qquad
V=(Q]^{c}.
\]
La hipótesis \(V\subseteq U\) equivale a
\[
(Q]^{c} \subseteq [P).
\tag{1}
\]

Como \(P\) y \(Q\) son filtros primos, no puede ocurrir que \(P\subseteq Q\);
por lo tanto existe \(a\in P\) con \(a\notin Q\).
Fijemos tal elemento \(a\in L\).
Probaremos que
\[
V \subseteq \sigma_{\mathbb{L}}(a)\subseteq U.
\]

\begin{itemize}
\item[a)] \(V\subseteq \sigma_{\mathbb{L}}(a)\).

Sea \(Z\in V=(Q]^c\).
Entonces \(Z\nsubseteq Q\), es decir,
\(Z\notin (Q]\).
De (1) se sigue que todo elemento que no pertenece a \((Q]\)
pertenece a \([P)\), luego \(Z\in [P)\),
y en particular \(P\subseteq Z\).
Como \(a\in P\), se tiene \(a\in Z\),
y por definición de \(\sigma_{\mathbb{L}}(a)\),
esto implica \(Z\in \sigma_{\mathbb{L}}(a)\).
Por tanto, \(V\subseteq\sigma_{\mathbb{L}}(a)\).

\item[b)] \(\sigma_{\mathbb{L}}(a)\subseteq U\).

Sea \(R\in\sigma_{\mathbb{L}}(a)\).
Entonces \(a\in R\).
Supongamos, con el fin de obtener una contradicción, que
\(R\in(Q]\), es decir \(R\subseteq Q\).
De ello se seguiría \(a\in Q\),
en contradicción con nuestra elección de \(a\notin Q\).
Por tanto, \(R\notin(Q]\).
Usando nuevamente (1), \(R\in [P)=U\).
De aquí se deduce \(\sigma_{\mathbb{L}}(a)\subseteq U\).
\end{itemize}

Esto muestra que existe \(a\in L\) tal que
\(V \subseteq \sigma_{\mathbb{L}}(a)\subseteq U\),
y por lo tanto que \(\sigma_{\mathbb{L}}(\mathbb L)\) es separante en
\(\mathbb{L}^{\delta}\).


\item $\sigma_{\mathbb{L}}(\mathbb L)$ es compacto en $\Up(\mathcal{X}(\mathbb L))$.

Sean $X,Y\subseteq L$ tales que
\[
\bigwedge_{\Up(\mathcal{X}(\mathbb L))}\sigma_{\mathbb{L}}(X)\;\le\;
\bigvee_{\Up(\mathcal{X}(\mathbb L))}\sigma_{\mathbb{L}}(Y)
\]
Es decir, en nuestro contexto:
\[
\bigcap\sigma_{\mathbb{L}}[X]\;\subseteq\;\bigcup\sigma_{\mathbb{L}}[Y]
\]
Consideremos $\Fg^{\mathbb{L}}(X)$ el filtro generado por el conjunto $X$ e $\Ig^{\mathbb{L}}(Y)$ el ideal generado por el conjunto $Y$. Vamos a comprobar que $\Fg^{\mathbb{L}}(X)\cap \Ig^{\mathbb{L}}(Y)\neq\emptyset$. Asumiremos, por absurdo, que 
\[ \Fg^{\mathbb{L}}(X)\cap \Ig^{\mathbb{L}}(Y) = \emptyset. \]
Por el Teorema del filtro primo, existe $P\in\mathcal{X}(\mathbb{L})$ tal que $\Fg^{\mathbb{L}}(X)\subseteq P$ y $P\cap \Ig^{\mathbb{L}}(Y)\neq\emptyset$. Asi puesto que $X\subseteq \Fg^{\mathbb{L}}(X)\subseteq P$ tenemos que para cada $x \in X$,  $x\in P$.  Luego, para cada $x \in X$, se tiene que $P\in \sigma_{\mathbb{L}}(x)$ y por lo tanto $P\in\bigcap\sigma_{\mathbb{L}}[X]\;\subseteq\;\bigcup\sigma_{\mathbb{L}}[Y]$. Así existe $y\in Y$ tal que $P\in \sigma_{\mathbb{L}}(y)$, lo que resulta imposible. Se sigue entonces que $P\cap Y\neq \emptyset$. 

Se sigue entonces que existe  $z\in L$ tal que $z\in \Fg^{\mathbb{L}}(X)\cap \Ig^{\mathbb{L}}(Y)$. Por la definiciones de filtro e ideal generado, tenemos que existen $x_1,x_2,...,x_n\in X$ y $y_1,y_2,...,y_m\in Y$ tales que
\begin{align*}
    x_1\wedge x_2\wedge...\wedge x_n\leq z\leq y_1\vee y_2\vee...\vee y_m
\end{align*}
Tomando $F=\{x_1,x_2,...,x_n\}$ y $G=\{y_1,y_2,...,y_m\}$, se sigue que
\begin{align*}
    \bigwedge_\mathbb L{} F\leq \bigvee_{\mathbb L} G
\end{align*}
Teniendo en cuenta que $\sigma_{\mathbb{L}}$ es un homomorfismo de retículos
\begin{align*}
    \bigwedge_{\Up(\mathcal{X}(\mathbb L))} \sigma_{\mathbb{L}}[F]=\sigma_{\mathbb{L}}\left(\bigwedge_{\mathbb{L}} F\right)\leq \sigma_{\mathbb{L}}\left(\bigvee_{\mathbb{L}} G\right) =\bigvee_{\Up(\mathcal{X}(\mathbb L))} \sigma_{\mathbb{L}}[G]
\end{align*}
Con $F$ y $G$ claramente finitos. Es decir, la condición de la Definición~\ref{def:compacto-en} se satisface.
\end{enumerate}
Por lo que $\left(\Up(\mathcal{X}(\mathbb{L})),\sigma_{\mathbb{L}} \right)$ es \emph{la} extensión canónica de $\mathbb{L}$, y la nombraremos entonces también como $(\mathbb{L}^\delta, \sigma_{\mathbb{L}})$.
\end{proof}


\begin{theorem}\label{thm:Up-distribuida-acotada}
$\Up(\mathcal{X}(\mathbb{L}))$ es un retículo distributivo acotado.
\end{theorem}

\begin{proof}[\textbf{Demostración}]
Mostraremos ambas propiedades del retículo $\Up(\mathcal{X}(\mathbb{L}))$.
\begin{itemize}
    \item Acotación
    
Tanto $\emptyset$ como $\mathcal{X}(\mathbb{L})$ son crecientes.  
Para todo $U\in\Up(\mathcal{X}(\mathbb{L}))$ se tiene
\[
\emptyset \subseteq U \subseteq \mathcal{X}(\mathbb{L}),
\]
luego $\emptyset$ es la cota inferior (mínimo) y $\mathcal{X}(\mathbb{L})$ es la cota superior (máximo) en $\Up(\mathcal{X}(\mathbb{L}))$.

\item Distributividad

Si $U,V,W\in\Up(\mathcal{X}(\mathbb{L}))$, entonces $U\cap V$, $U\cap W$, $V\cup W$ y $(U\cap V)\cup(U\cap W)$ son crecientes.  
Usando sólo cuentas conjuntistas, para todo $x\in \mathcal{X}(\mathbb{L})$:
\[
\begin{aligned}
x\in U\cap(V\cup W)
&\ \Longleftrightarrow\ x\in U\ \text{y}\ \bigl(x\in V\ \text{o}\ x\in W\bigr) \\
&\ \Longleftrightarrow\ \bigl(x\in U\cap V\bigr)\ \text{o}\ \bigl(x\in U\cap W\bigr) \\
&\ \Longleftrightarrow\ x\in (U\cap V)\cup(U\cap W).
\end{aligned}
\]
Por lo tanto,
\[
U\cap(V\cup W)=(U\cap V)\cup(U\cap W).
\]
La distributividad dual se verifica del mismo modo:
\[
U\cup(V\cap W)=(U\cup V)\cap(U\cup W).
\]
\end{itemize}
Concluimos que $\Up(\mathcal{X}(\mathbb{L}))$ es un retículo distributivo acotado.
\end{proof}

\begin{example}\label{ej:M3-no-distributivo}
Aclaración. No todo retículo completo es distributivo. Un ejemplo clásico es el retículo \(\mathbb{M}_3\), cuyos elementos son
\[
0,\ a,\ b,\ c,\ 1
\]
con relaciones de orden dadas únicamente por
\[
0<a<1,\qquad 0<b<1,\qquad 0<c<1,
\]
y sin comparabilidades adicionales entre \(a,b,c\). Este retículo es completo ya que es finito , pero no es distributivo. En efecto, se verifica que
\[
a\wedge(b\vee c)=a\qquad\text{mientras que}\qquad
(a\wedge b)\vee(a\wedge c)=0,
\]
por lo que la distributividad falla. 
\end{example}

\begin{lemma}\label{lem:creciente-cerrado} Sean $\mathbb L$ un retículo distributivo acotado y $\left( \mathbb{L}^{\delta},\sigma_{\mathbb{L}} \right)$ su extensión canónica. Si $Q\in\mathcal X(\mathbb L)$ entonces el creciente principal $[Q)$ es un cerrado en $\mathbb{L}^{\delta}$. En particular, $(Q]^c$ es un abierto.
\end{lemma}

\begin{proof}[\textbf{Demostración:}]
Sea $Z\in[Q)$. Por definición, se sigue que $Q\subseteq Z$, es decir, $a\in Z$ para todo
$a\in Q$. Esta sencilla observación nos dice que
\[ Z \in \bigcap_{\mathbb{L}^{\delta}} \{ \sigma_{\mathbb{L}}(a) \colon a \in Q\}. \]

Se sigue entonces que $[Q) \subseteq \bigcap_{\mathbb{L}^{\delta}} \{ \sigma_{\mathbb{L}}(a) \colon a \in Q\}$. Recíprocamente, si $Z\in\bigcap_{a\in Q}\sigma_{\mathbb L}(a)$, entonces
$a\in Z$ para todo $a\in Q$, es decir, $Q\subseteq Z$, y por definición
$Z\in[Q)$.

Concluimos que
\[
[Q)\;=\;\bigcap_{a\in Q}\sigma_{\mathbb L}(a),
\]
Teniendo en cuenta la Definición \ref{def:cerrados}, esta intersección es un cerrado de $\mathbb L^\delta$.

Sea  ahora $Q \in X(\mathbb{L})$. Vamos a comprobar que 
\[ (Q]^{c}  = \bigcup_{\mathbb{L}^{\delta}} \{ \sigma_{\mathbb{L}}(a) \colon a \notin Q \}. \]
En efecto, notemos que: 
\begin{align*}
P \in (Q]^{c} & \text{ si y solo si } P \not \subseteq Q \\
              & \text{ si y solo si } \exists a \in P\setminus Q \\
              & \text{ si y solo si } P \in \bigcup_{\mathbb{L}^{\delta}} \{\sigma_{\mathbb{L}}(a) \colon a \notin Q \}
\end{align*}
De donde se sigue que $(Q]^{c}$ es un abierto de la extensión canónica.
\end{proof}


\subsection{Filtros, ideales y su representación en la extensión canónica}

En la presente subsección, vamos a estudiar la manera de describir los filtros e ideales de un retículo distributivo dado $\mathbb{L}$, en su extensión canónica $\left( \mathbb{L}^{\delta},\sigma_{\mathbb{L}} \right)$, la cual construimos en el Teorema \ref{teo:ext-canonica-sigma}.


\begin{theorem}\label{thm:cerrados-filtros}
Sea \(\mathbb{L}\) un retículo distributivo acotado y $(\mathbb{L}^{\delta},\sigma_{\mathbb{L}})$ su extensión canónica. Entonces:
\[
U\subseteq\mathcal{X}(\mathbb{L})
\quad
\text{es cerrado en }\mathbb{L}^{\delta}
\quad\Longleftrightarrow\quad
\exists\,F\in\mathrm{Fi}(\mathbb{L}):\;
U \;=\;\bigcap_{a\in F}\sigma_{\mathbb{L}}(a).
\]
Además, esta representación es única: si \(\bigcap_{a\in F}\sigma_{\mathbb{L}}(a)
=\bigcap_{a\in G}\sigma_{\mathbb{L}}(a)\) con \(F,G\in\mathrm{Fi}(\mathbb{L})\), 
entonces \(F=G\).
\end{theorem}

\begin{proof}[\textbf{Demostración:}]
Veamos la doble implicación y luego la unicidad de la representación.
\begin{itemize}
    \item $\Leftarrow$) Es inmediato por la Definición \ref{def:cerrados}. 
    \item $\Rightarrow$) Supongamos \(U\in\mathcal{C}(\mathbb{L}^\delta)\). Por definición existe \(A\subseteq L\) tal que
\[
  U \;=\;\bigcap_{a\in A}\sigma_{\mathbb{L}}(a).
\]
Sea \(\Fg^{\mathbb{L}}(A)\) el filtro \emph{generado} por \(A\), es decir
\[
  \Fg^{\mathbb{L}}(A) \;=\;
  \bigl\{\,z\in L : \exists\,x_1,\dots,x_n\in A,\;
                 x_1\wedge\cdots\wedge x_n\le z
  \bigr\}.
\]
Como \(\sigma_{\mathbb{L}}\) preserva ínfimos finitos,
\[
  \bigcap_{a\in \Fg^{\mathbb{L}}(A)}\sigma_{\mathbb{L}}(a)
  =\bigcap_{\substack{x_1\wedge\cdots\wedge x_n\in A}}
    \sigma_{\mathbb{L}}(x_1\wedge\cdots\wedge x_n)
  =\bigcap_{a\in A}\sigma_{\mathbb{L}}(a)
  =U.
\]
Por tanto \(U=\bigcap_{a\in \Fg^{\mathbb{L}}(A)}\sigma_{\mathbb{L}}(a)\) y \(\Fg^{\mathbb{L}}(A)\) es filtro, como se quería.

\item Unicidad: Supongamos que \(F,G\in\mathrm{Fi}(\mathbb L)\) satisfacen
\[
\bigcap_{a\in F}\sigma_{\mathbb{L}}(a)
=\bigcap_{b\in G}\sigma_{\mathbb{L}}(b).
\]
Demostraremos \(F=G\).

Supongamos que existe \(a\in F\setminus G\). Es decir que 
\begin{align*}
    G\cap(a] = \emptyset
\end{align*}
Puesto que $G\in\mathrm{Fi}(\mathbb L)$ y $(a]$ es un ideal, por el Teorema del filtro primo existe $P\in \mathcal{X}(\mathbb L)$ tal que $G\subseteq P$ y $a\notin P$. De esto último, $P\notin \bigcap_{a\in F} \sigma_{\mathbb{L}}(a)=\bigcap_{b\in G}\sigma_{\mathbb{L}}(b)$. Es decir que existe $b\in G$ tal que $b\notin P$, lo cual contradice que $G\subseteq P$. Este absurdo está dado por suponer que existe \(a\in F\setminus G\). Con un argumento similar podemos mostrar que tampoco exite  \(b\in G\setminus F\), mostrando por lo tanto que \(F=G\).
\end{itemize}
\end{proof}

\begin{theorem}\label{thm:abiertos-ideales}
Sea \(\mathbb{L}\) un retículo distributivo acotado y $(\mathbb{L}^{\delta},\sigma_{\mathbb{L}})$ su extensión canónica. Entonces:
\[
U\subseteq\mathcal{X}(\mathbb{L})
\quad
\text{es abierto en }\mathbb{L}^{\delta}
\quad\Longleftrightarrow\quad
\exists\,I\in\mathrm{Ide}(\mathbb{L}):\;
U \;=\;\bigcup_{a\in I}\sigma_{\mathbb{L}}(a).
\]
Además, esta representación es única: si \(\bigcup_{a\in I}\sigma_{\mathbb{L}}(a)
=\bigcup_{a\in J}\sigma_{\mathbb{L}}(a)\) con \(I,J\in\mathrm{Ide}(\mathbb{L})\), 
entonces \(I=J\).
\end{theorem}

\begin{proof}[\textbf{Demostración:}]
Procederemos como en el caso de los cerrados algebraicos, probando las dos implicaciones y la unicidad de la representación.

\begin{itemize}
\item \(\Leftarrow\)) Es inmediato por la Definición \ref{def:abiertos}. 

\item \(\Rightarrow\)) Sea \(U\subseteq\mathcal{X}(\mathbb L)\) abierto.
Por definición existe \(A\subseteq L\) tal que
\[
U \;=\; \bigcup_{a\in A}\sigma_{\mathbb{L}}(a).
\]
Sea \(I\) el \emph{ideal} generado por \(A\), es decir
\[
I \;=\;
\{\,z\in L : \exists\,x_1,\dots,x_n\in A,\;
                 z \le x_1\vee\cdots\vee x_n\,\}.
\]
Observamos inmediatamente que
\[
\bigcup_{a\in I}\sigma_{\mathbb{L}}(a)
=\bigcup_{\substack{x_1\vee\cdots\vee x_n\in I}}
  \sigma_{\mathbb{L}}(x_1\vee\cdots\vee x_n)
=\bigcup_{a\in A}\sigma_{\mathbb{L}}(a)
=U,
\]
pues cada \(\sigma_{\mathbb{L}}(x_1\vee\cdots\vee x_n)=\bigcup_{i=1}^n\sigma_{\mathbb{L}}(x_i)\) y las uniones finitas
sean absorbidas por la unión sobre \(A\). Así \(U=\bigcup_{a\in I}\sigma_{\mathbb{L}}(a)\) con \(I\) ideal.

\item Unicidad: Supongamos que \(F,G\in\mathrm{Ide}(\mathbb L)\) satisfacen
\[
\bigcup_{a\in F}\sigma_{\mathbb{L}}(a)
=\bigcup_{b\in G}\sigma_{\mathbb{L}}(b).
\]
Demostraremos \(F=G\).

Supongamos que existe \(a\in F\setminus G\). Es decir que 
\begin{align*}
    G\cap[a) = \emptyset
\end{align*}
Puesto que $G\in\mathrm{Ide}(\mathbb L)$ y $[a)$ es un filtro, por el Teorema del filtro primo existe $P\in \mathcal{X}(\mathbb L)$ tal que $G\cap P= \emptyset$ y $a\in P$. De esto último, $P\in \bigcup_{a\in F} \sigma_{\mathbb{L}}(a)=\bigcup_{b\in G}\sigma_{\mathbb{L}}(b)$. Es decir que existe $b\in G$ tal que $P\in\sigma_{\mathbb{L}}(b)$, y por lo tanto $b\in P$, lo cual contradice que $G\cap P= \emptyset$. Este absurdo está dado por suponer que existe \(a\in F\setminus G\). Con un argumento similar podemos mostrar que tampoco exite  \(b\in G\setminus F\), mostrando por lo tanto que \(F=G\).
\end{itemize}
\end{proof}

Sea $\mathbb{L}$ un retículo distributivo acotado y $(\mathbb{L}^{\delta},\sigma_{\mathbb{L}})$ su extensión canónica. Como consecuencia de los Teoremas \ref{thm:cerrados-filtros} y \ref{thm:abiertos-ideales}, 
obtenemos una identificación biunívoca entre filtros e ideales de $\mathbb{L}$, con cerrados algebraicos y abiertos algebraicos de $\mathbb{L}^{\delta}$, respectivamente. Dicha identificación univoca la formalizaremos en los siguiente apartados. Asimismo, estos resultados serán empleados en las siguientes secciones para construir y caracterizar las extensiones canónicas de operadores mediante relaciones sobre el conjunto de puntos $\mathcal X(\mathbb L)$.


\subsubsection{Isomorfismo entre filtros y cerrados}

En el presente apartado, utilizaremos el Teorema \ref{thm:cerrados-filtros}  para establecer un isomorfismo dual de posets, entre los filtros de un retículo distributivo acotado $\mathbb{L}$ y los cerrados algebraicos de su extensión canónica \(\mathbb{L}^\delta\).

\begin{theorem}\label{thm:dualidad-filtros-cerrados}
Sean \(\mathbb{L}\) un retículo distributivo acotado y
\((\Up(\mathcal{X}(\mathbb{L})),\sigma_{\mathbb{L}})\) su extensión canónica. Definamos las aplicaciones
\[
\begin{aligned}
\Psi \; &: \; \mathrm{Fi}(\mathbb{L}) \longrightarrow \mathcal{C}(\mathbb{L}^\delta), 
\qquad &\Psi(F) &\;=\; \bigcap_{a \in F} \sigma_{\mathbb{L}}(a),\\[4pt]
\Phi \; &: \; \mathcal{C}(\mathbb{L}^\delta) \longrightarrow \mathrm{Fi}(\mathbb{L}),
\qquad &\Phi(C) &\;=\; \{\,a \in L : C \subseteq \sigma_{\mathbb{L}}(a)\,\}.
\end{aligned}
\]
Entonces \(\Psi\) y \(\Phi\) son antimonótonas y verifican
\(\Phi\circ\Psi=\Id_{\mathrm{Fi}(\mathbb{L})}\) y
\(\Psi\circ\Phi=\Id_{\mathcal{C}(\mathbb{L}^\delta)}\). En particular,
\[
  \bigl(\mathrm{Fi}(\mathbb{L}),\subseteq\bigr)^{\!op}
  \;\cong\;
  \bigl(\mathcal{C}(\mathbb{L}^\delta),\subseteq\bigr).
\]
\end{theorem}

\begin{proof}[\textbf{Demostración:}]
Procederemos en cuatro pasos.

\begin{enumerate}
\item $\Psi$ está bien definida y es antimonótona.

Por el Teorema \ref{thm:cerrados-filtros}, para cada filtro \(F\) de \(\mathbb L\) el conjunto
\(\bigcap_{a\in F}\sigma_{\mathbb{L}}(a)\) es un cerrado algebraico, por lo que \(\Psi(F)\in\mathcal C(\mathbb L^\delta)\).
Si \(F_1\subseteq F_2\) entonces claramente
\[\Psi(F_2)=\bigcap_{a\in F_2}\sigma_{\mathbb{L}}(a)\subseteq\bigcap_{a\in F_1}\sigma_{\mathbb{L}}(a)=\Psi(F_1)\]
por lo que \(\Psi\) es antimonótona.

\item $\Phi$ está bien definida y es antimonótona

Sea \(C\in\mathcal C(\mathbb L^\delta)\). Definimos
\(\Phi(C)=\{a\in L : C\subseteq\sigma_{\mathbb{L}}(a)\}\). Si \(a\le b\) y \(a\in\Phi(C)\)
entonces \(C\subseteq\sigma_{\mathbb{L}}(a)\subseteq\sigma_{\mathbb{L}}(b)\), así \(b\in\Phi(C)\). Si
\(a,b\in\Phi(C)\) entonces \(C\subseteq\sigma_{\mathbb{L}}(a)\cap\sigma_{\mathbb{L}}(b)=\sigma_{\mathbb{L}}(a\wedge b)\),
por lo que \(a\wedge b\in\Phi(C)\). Luego \(\Phi(C)\) es un filtro de \(L\).
Además, si \(C_1\subseteq C_2\) entonces \(\Phi(C_2)\subseteq\Phi(C_1)\),
por lo que \(\Phi\) es antimonótona.

\item $\Phi\circ\Psi=\Id_{\mathrm{Fi}(\mathbb{L})}$

Sea \(F\in\mathrm{Fi}(\mathbb L)\). Por definición
\[
\Phi(\Psi(F))
=\{a\in L:\bigcap_{x\in F}\sigma_{\mathbb{L}}(x)\subseteq\sigma_{\mathbb{L}}(a)\}.
\]
Mostramos que este conjunto es exactamente \(F\).

- Si \(a\in F\) entonces \(\bigcap_{x\in F}\sigma_{\mathbb{L}}(x)\subseteq\sigma_{\mathbb{L}}(a)\) es inmediato, por lo que
  \(F\subseteq \Phi(\Psi(F))\).

- Supongamos \(a\notin F\). Como \(F\) es filtro y \(a\notin F\), los conjuntos \(F\) y \((a]\)
  son disjuntos. Por el Teorema del filtro primo existe un filtro primo \(P\) que contiene a \(F\)
  y no contiene a \(a\). Entonces \(P\in\bigcap_{x\in F}\sigma_{\mathbb{L}}(x)\) pero \(P\notin\sigma_{\mathbb{L}}(a)\),
  es decir \(\bigcap_{x\in F}\sigma_{\mathbb{L}}(x)\not\subseteq\sigma_{\mathbb{L}}(a)\). Así \(a\notin\Phi(\Psi(F))\).

Combinando ambas inclusiones obtenemos \(\Phi(\Psi(F))=F\).

\item $\Psi\circ\Phi=\Id_{\mathcal{C}(\mathbb{L}^\delta)}$

Sea \(C\in\mathcal C(\mathbb L^\delta)\). Por definición
\[
\Psi(\Phi(C))
=\bigcap_{\substack{a\in L\\ C\subseteq\sigma_{\mathbb{L}}(a)}}\sigma_{\mathbb{L}}(a).
\]
Como \(C\) es un cerrado algebraico, existe \(A\subseteq L\) con \(C=\bigcap_{a\in A}\sigma_{\mathbb{L}}(a)\).
En particular, todo \(a\in A\) satisface \(C\subseteq\sigma_{\mathbb{L}}(a)\), por lo que
\(\bigcap_{a\in A}\sigma_{\mathbb{L}}(a)\supseteq\bigcap_{C\subseteq\sigma_{\mathbb{L}}(b)}\sigma_{\mathbb{L}}(b)\).
Por otro lado, la intersección sobre todos los \(\sigma_{\mathbb{L}}(b)\) que contienen
\(C\) está contenida en cada uno de ellos, y por tanto en \(C\). En consecuencia
\(\Psi(\Phi(C))=C\).
\end{enumerate}
\medskip
Así vimos que \(\Psi\) y \(\Phi\) son antimonótonas mutuamente inversas,
lo que exhibe el anti-isomorfismo de posets anunciado.
\end{proof}

\begin{corollary}\label{cor:caracterizacion-C-F}
Sean $\mathbb{L}$ un retículo distributivo acotado y 
$(\mathbb{L}^{\delta},\sigma_{\mathbb{L}})$ su extensión canónica.
Para todo $C\in\mathcal{C}(\mathbb{L}^\delta)$ y todo filtro primo 
$Q\in\mathcal{X}(\mathbb{L})$ se cumple
\[
Q \in C 
\quad\Longleftrightarrow\quad 
\Phi(C)\subseteq Q.
\]
\end{corollary}

\begin{proof}[\textbf{Demostración:}]
Por el Teorema~\ref{thm:dualidad-filtros-cerrados}, 
\[
C=\Psi(\Phi(C))
=\bigcap_{a\in \Phi(C)}\sigma_{\mathbb{L}}(a).
\]
Por tanto, para todo \(Q\in\mathcal{X}(\mathbb L)\),
\[
Q\in C
\ \Longleftrightarrow\
\forall a\in\Phi(C),\ Q\in\sigma_{\mathbb{L}}(a)
\ \Longleftrightarrow\
\forall a\in\Phi(C),\ a\in Q
\ \Longleftrightarrow\
\Phi(C)\subseteq Q.
\]
\end{proof}

\subsubsection{Isomorfismo entre ideales y abiertos}

De manera análoga al caso de filtros y cerrados algebraicos, describimos a continuación el isomorfismo entre el poset de ideales de \(\mathbb{L}\) y el de abiertos de \(\mathbb{L}^\delta\).

\begin{theorem}\label{thm:dualidad-ideales-abiertos}
Sea \(\mathbb{L}\) un retículo distributivo acotado y 
\((\mathbb{L}^{\delta},\sigma_{\mathbb{L}})\) su extensión canónica. 
Entonces las aplicaciones
\[
\begin{aligned}
\Theta \; &: \; \mathrm{Ide}(\mathbb{L}) \longrightarrow \mathcal{O}(\mathbb{L}^\delta), 
\qquad &\Theta(I) &\;=\; \bigcup_{a \in I} \sigma_{\mathbb{L}}(a),\\[4pt]
\Gamma \; &: \; \mathcal{O}(\mathbb{L}^\delta) \longrightarrow \mathrm{Ide}(\mathbb{L}),
\qquad &\Gamma(U) &\;=\; \{\,a \in L : \sigma_{\mathbb{L}}(a)\subseteq U\,\}.
\end{aligned}
\]
son monótonas y satisfacen
\(\Gamma\circ\Theta=\Id_{\mathrm{Ide}(\mathbb{L})}\) y
\(\Theta\circ\Gamma=\Id_{\mathcal{O}(\mathbb{L}^\delta)}\).
Por lo tanto,
\[
  \bigl(\mathrm{Ide}(\mathbb{L}),\subseteq\bigr)
  \;\cong\;
  \bigl(\mathcal{O}(\mathbb{L}^\delta),\subseteq\bigr).
\]
\end{theorem}

\begin{proof}[\textbf{Demostración:}]
Haremos la prueba de manera análoga a la prueba del Teorema \ref{thm:dualidad-filtros-cerrados}.

\begin{enumerate}
\item \(\Theta\) bien definida y monótona.  

Por el Teorema~\ref{thm:abiertos-ideales}, para cada ideal \(I\subseteq L\),
\(\Theta(I)=\bigcup_{a\in I}\sigma_{\mathbb{L}}(a)\) es un abierto.  
Además, si \(I_1 \subseteq I_2\), entonces
\[
  \Theta(I_1)=\bigcup_{a\in I_1}\sigma_{\mathbb{L}}(a)
  \;\subseteq\;
  \bigcup_{a\in I_2}\sigma_{\mathbb{L}}(a)=\Theta(I_2),
\]
de modo que \(\Theta\) es monótona.

\medskip
\item \(\Gamma\) bien definida y monótona.  

Sea \(U\in\mathcal{O}(\mathbb{L}^\delta)\). Definimos
\[
\Gamma(U)=\{a\in L : \sigma_{\mathbb{L}}(a)\subseteq U\}.
\]

- Si \(a\le b\) y \(b\in\Gamma(U)\), entonces
  \(\sigma_{\mathbb{L}}(a)\subseteq\sigma_{\mathbb{L}}(b)\subseteq U\), de modo que \(a\in\Gamma(U)\).
- Si \(a,b\in\Gamma(U)\), entonces
  \(\sigma_{\mathbb{L}}(a\vee b)=\sigma_{\mathbb{L}}(a)\cup\sigma_{\mathbb{L}}(b)\subseteq U\), luego \(a\vee b\in\Gamma(U)\).

Así, \(\Gamma(U)\) es un ideal de \(L\).  
Además, si \(U_1\subseteq U_2\), entonces 
\(\Gamma(U_1)\subseteq\Gamma(U_2)\), lo que prueba que \(\Gamma\) es monótona.


\item $\Gamma\circ\Theta=\Id_{\mathrm{Ide}(\mathbb{L})}$.  

Sea $I\in\mathrm{Ide}(\mathbb{L})$. Entonces
\[
\Gamma(\Theta(I))
=\{\,a\in L:\sigma_{\mathbb{L}}(a)\subseteq\bigcup_{x\in I}\sigma_{\mathbb{L}}(x)\,\}.
\]
Si $a\in I$, la inclusión es inmediata; veamos la recíproca.  
Supongamos $a\notin I$. Como $I$ es ideal, $I\cap[a)=\emptyset$.  
Por el teorema del filtro primo, existen un filtro primo $P$ con 
$[a)\subseteq P$ y $I\cap P=\emptyset$.  
De $a\in P$ se sigue $P\in\sigma_{\mathbb{L}}(a)$, mientras que
$I\cap P=\emptyset$ implica $P\notin\sigma_{\mathbb{L}}(x)$ para todo $x\in I$.  
Así $P\in\sigma_{\mathbb{L}}(a)\setminus\bigcup_{x\in I}\sigma_{\mathbb{L}}(x)$, 
por lo que $\sigma_{\mathbb{L}}(a)\nsubseteq\bigcup_{x\in I}\sigma_{\mathbb{L}}(x)$
y $a\notin\Gamma(\Theta(I))$.  
Concluimos que $\Gamma(\Theta(I))=I$.


\medskip
\item \(\Theta\circ\Gamma=\Id_{\mathcal{O}(\mathbb{L}^\delta)}\).  

Sea \(U\) un abierto. Entonces
\[
\Theta(\Gamma(U))
=\bigcup_{a\in\Gamma(U)}\sigma_{\mathbb{L}}(a)
=\bigcup_{\substack{a\in L\\\sigma_{\mathbb{L}}(a)\subseteq U}}\sigma_{\mathbb{L}}(a)=U.
\]
\end{enumerate}

\medskip
Los cuatro pasos anteriores muestran que \(\Theta\) y \(\Gamma\) son aplicaciones monótonas mutuamente inversas, lo que concluye el isomorfismo de posets.
\end{proof}

\begin{corollary}\label{cor:caracterizacion-U-Gamma}
Sea $\mathbb{L}$ un retículo distributivo acotado y 
$(\mathbb{L}^{\delta},\sigma_{\mathbb{L}})$ su extensión canónica.
Para todo abierto $O\in\mathcal{O}(\mathbb{L}^\delta)$ y todo $a\in L$ se cumple
\[
\sigma_{\mathbb{L}}(a)\subseteq O
\quad\Longleftrightarrow\quad
a\in\Gamma(O).
\]
\end{corollary}


\begin{proof}[\textbf{Demostración:}]
Para un $a\in L$ se tiene:
\[
a\in\Gamma(O)
\;\Longleftrightarrow\;
\sigma_{\mathbb L}(a)\subseteq O
\]
por definición misma de $\Gamma$.
\end{proof}
Para finalizar esta sección, haremos una conexión de las nociones de abiertos y cerrados de la extensión canónica $(\mathbb{L}^{\delta},\sigma_{\mathbb{L}})$ para un retículo disitributivo acotado dado $\mathbb{L}$, con las nociones de conjuntos abiertos y cerrados del espacio de Priestley dual de $\mathbb{L}$.

\subsection{Conexiones con nociones topológicas}

En la presente subsección, haremos una conexión entre las estructuras algebraicas y las estructuras topológicas que hemos trabajo a lo largo de esta tesis. En particular, dado un retículo distributivo acotado $\mathbb{L}$, vamos a trabajar simultáneamente con su extensión canónica $(\mathbb{L}^{\delta},\sigma_{\mathbb{L}})$ y su espacio de Priestley dual $(\mathcal{X}(\mathbb{L}),\tau_{\mathbb{L}},\subseteq)$.

\vspace{0.2cm}

Comenzaremos haciendo una observación que será de gran interés en la presente subsección y a lo largo de este capítulo.

\begin{remark} Sean $\mathbb{L}$ un retículo distributivo acotado, $(\mathcal{X}(\mathbb{L}),\tau_{\mathbb{L}},\subseteq)$ su espacio de Priestley dual y $\left(\mathbb{L}^{\delta},\sigma_{\mathbb{L}}\right)$ su extensión canónica. Recordemos que, para $a \in L$, el conjunto 
\[ \sigma_{\mathbb{L}}(a) = \{ P \in \mathcal{X}(\mathbb{L}) \colon a \in P \},\]
es un conjunto clopen creciente del espacio de Priestley asociado a $\mathbb{L}$. Asimismo, teniendo en cuenta la Proposición \ref{prop:clopen-is-image}, el conjunto $\sigma_{\mathbb{L}}(a)$ puede interpretarse como un conjunto abierto en $\mathbb{L}^{\delta}$ y cerrado en $\mathbb{L}^{\delta}$, simultáneamente. Por lo tanto, el conjunto $\sigma_{\mathbb{L}}(a)$ es un clopen, ya sea desde el punto de vista topológico sobre el espacio de Priestley $(\mathcal{X}(\mathbb{L}),\tau_{\mathbb{L}},\subseteq)$ como asi también desde el punto de vista algebraico en $\mathbb{L}^{\delta}$.   
\end{remark}

A continuación, enunciaremos un resultado que permite establecer conexiones entre conceptos algebraicos y conceptos topológicos.


\begin{theorem} Sean $\mathbb{L}$ un retículo distributivo acotado, $(\mathcal{X}(\mathbb{L}),\tau_{\mathbb{L}},\subseteq)$ su espacio de Priestley dual y $\left( \mathbb{L}^{\delta},\sigma_{\mathbb{L}}\right)$ su estensión canónica. Entonces
\begin{itemize}
\item[1.] $U \in \mathcal{C}(\mathbb{L}^{\delta})$ si y sólo si $U$ es un conjunto cerrado creciente de $(\mathcal{X}(\mathbb{L}),\tau_{\mathbb{L}},\subseteq)$
\item[2.] $U \in \mathcal{O}(\mathbb{L}^{\delta})$ si y sólo si $U$ es un conjunto abierto creciente de $(\mathcal{X}(\mathbb{L}),\tau_{\mathbb{L}},\subseteq)$
\end{itemize}
\end{theorem}
\begin{proof}[\textbf{Demostración:}] Vamos a comprobar solamente (1), ya que un razonamiento análogo puede hacerse para $(2)$. En efecto, consideremos $U \in \mathcal{C}(\mathbb{L}^{\delta})$. Teniendo en cuenta el Teorema \ref{thm:dualidad-filtros-cerrados}, se sigue que existe un filtro $F$ de $\mathbb{L}$ tal que \[ U = \bigcap\{ \sigma_{\mathbb{L}}(a) \colon a \in F \}. \]
Por el Corolario \ref{cor:Priestley_decomp}, se sigue que $U$ es un conjunto cerrado creciente de $(\mathcal{X}(\mathbb{L}),\tau_{\mathbb{L}},\subseteq)$. Recíprocamente, asumiremos que $U$ es un conjunto cerrado y creciente del espacio de Priestley asociado a $\mathbb{L}$. Nuevamente, por el Corolario \ref{cor:Priestley_decomp}, se sigue que $U$ es intersección de clopens crecientes de $(\mathcal{X}(\mathbb{L}),\tau_{\mathbb{L}},\subseteq)$. Esto es, 
\[  U = \bigcap \{\sigma_{\mathbb{L}}(a) \colon a \in A\}. \]
Luego, por el Teorema \ref{thm:cerrados-filtros}, se sigue que $U$ es cerrado en $\mathbb{L}^{\delta}$.
\end{proof}





\section{Extensión canónica de operadores unarios}

En la presente sección, vamos a dar algunas definiciones y resultados que serán de gran utilidad para el estudio de la extensión canónica para $(i,j)$-retículos distributivos acotados. En esta primera parte, hablaremos en términos generales, sin hacer mención a la extensión canónica $(\mathbb{L}^{\delta},\sigma_{\mathbb{L}})$ construida en la sección \ref{S4.2}. Además, adoptaremos la técnica para extender el operador unario utilizando la  formulación adoptada en la literatura (ver \cite{BDL_Expansions, BDL_Operators, GehrkeHarding2001, TesisLeo, TesisPaula}). y será fundamental para llevar adelante el estudio de los operadores normales de tipo $(i,j)$ en la sección siguiente.





\begin{definition}
Sean $\mathbb{L}$ y $\mathbb{M}$ dos retículos distributivos acotados, 
$e_\mathbb{L}\colon\mathbb{L}\hookrightarrow\mathbb{L}^\delta$ y 
$e_\mathbb{M}\colon\mathbb{M}\hookrightarrow\mathbb{M}^\delta$ sus extensiones canónicas, 
y sea $m\colon \mathbb{L}\to\mathbb{M}$ un operador monótono.  
Definimos las extensiones \emph{canónicas} $m^{\sigma},m^{\pi}\colon 
\mathbb{L}^\delta\to\mathbb{M}^\delta$ de la siguiente manera:
\begin{align*}
  m^{\sigma}(u)
  &=\bigvee_{\mathbb{M}^\delta}
     \Bigl\{
        \bigwedge_{\mathbb{M}^\delta}
        \{\, e_\mathbb{M}(m(a)) : a\in L,\; C\le e_\mathbb{L}(a)\,\}
      \;\Big|\;
        C\in\mathcal{C}(\mathbb{L}^\delta),\; C\le u
     \Bigr\},\\[3mm]
  m^{\pi}(u)
  &=\bigwedge_{\mathbb{M}^\delta}
     \Bigl\{
        \bigvee_{\mathbb{M}^\delta}
        \{\, e_\mathbb{M}(m(a)) : a\in L,\; e_\mathbb{L}(a)\le O\,\}
      \;\Big|\;
        O\in\mathcal{O}(\mathbb{L}^\delta),\; u\le O
     \Bigr\},
\end{align*}
para todo $u\in\mathbb{L}^\delta$.
\end{definition}
Notemos que las descripciones de $m^\sigma$ y $m^\pi$ se simplifican
considerablemente cuando se aplican a elementos cerrados o abiertos de
$\mathbb L^\delta$. Esto permite, en particular, verificar que ambas extensiones
restringen al operador original $m$ en la copia embebida $e[\mathbb L]$.
\begin{lemma}\label{lem:lambda-en-cerrado-general}
Sea $C\in\mathcal{C}(\mathbb{L}^\delta)$ un cerrado algebraico. Entonces
\[
  m^{\sigma}(C)
  \;=\;
  \bigwedge_{\mathbb{M}^\delta}\{\,e_\mathbb{M}(m(a)):\;a\in L,\; C\le e_\mathbb{L}(a)\,\}.
\]
\end{lemma}

\begin{proof}[\textbf{Demostración:}]
En la definición de $m^\sigma(u)$, al tomar $U=C$ cerrado algebraico, la unión externa sobre los
$x\le C$ se colapsa en el término correspondiente a $x=C$, ya que los demás son mayores.
\end{proof}

\begin{lemma}\label{lem:pi-en-abierto-general}
Sea $O\in\mathcal{O}(\mathbb{L}^\delta)$ un abierto. Entonces
\[
  m^{\pi}(O)
  \;=\;
  \bigvee_{\mathbb{M}^\delta}\{\,e_\mathbb{M}(m(a)):\;a\in L,\; e_\mathbb{L}(a)\le O\,\}.
\]
\end{lemma}

\begin{proof}[\textbf{Demostración:}]
La demostración es dual: en la definición de $m^\pi(u)$ basta considerar el abierto $O$ mismo en la familia de ínfimos.
\end{proof}
En otras palabras, las fórmulas para $m^\sigma$ y $m^\pi$ se simplifican en cerrados y abiertos algebraicos. Esto nos permite verificar que, en efecto, ambas construcciones merecen el nombre de \emph{extensiones}, ya que coinciden con $m$ en los elementos de $\mathbb L$ vistos dentro de $\mathbb L^\delta$.
\begin{theorem}\label{thm:extension-respects-original-general}
Para todo $a\in L$ se cumple
\[
  m^{\sigma}\bigl(e_\mathbb{L}(a)\bigr)
  = e_\mathbb{M}\bigl(m(a)\bigr)
  = m^{\pi}\bigl(e_\mathbb{L}(a)\bigr).
\]
En particular, ambas extensiones restringen a $m$ sobre la copia de $\mathbb{L}$ en $\mathbb{L}^\delta$.
\end{theorem}

\begin{proof}[\textbf{Demostración:}]
Sea $p=e_\mathbb{L}(a)$. Como $p$ es simultáneamente cerrado y abierto algebraico (Proposición ~\ref{prop:clopen-is-image}),
los Lemas \ref{lem:lambda-en-cerrado-general} y \ref{lem:pi-en-abierto-general} aplicados a $p$ dan
\[
m^\sigma(p)
  =\bigwedge\{e_\mathbb{M}(m(b)):a\le b\}
  =e_\mathbb{M}(m(a)),
\qquad
m^\pi(p)
  =\bigvee\{e_\mathbb{M}(m(b)):b\le a\}
  =e_\mathbb{M}(m(a)),
\]
usando la monotonía de $m$ y de las inclusiones $e_\mathbb{L},e_\mathbb{M}$.
\end{proof}

Ahora mencionaremos tres resultados generales sobre las extensiones canónicas de operadores unarios, cuyas demostraciones pueden consultarse en \cite[pp.~349–352]{GehrkeHarding2001}.
Dichos resultados establecen propiedades fundamentales de las extensiones
$m^{\sigma}$ y $m^{\pi}$, y nos permitirán posteriormente concluir que si se preservan los supremos, estas extensiones coinciden.


\begin{lemma}\label{lem:sigma-leq-pi}\cite[Lema~4.2]{GehrkeHarding2001}
Para todo $u\in\mathbb{L}^\delta$ se cumple $m^{\sigma}(u)\le m^{\pi}(u)$.
\end{lemma}

\begin{lemma}\label{lem:igualdad-por-preservacion}\cite[Lema~4.4]{GehrkeHarding2001}
Si $m^{\sigma}$ preserva supremos arbitrarios, entonces $m^{\sigma}=m^{\pi}$.
Dualmente, si $m^{\pi}$ preserva ínfimos arbitrarios, entonces $m^{\sigma}=m^{\pi}$.
\end{lemma}

\begin{lemma}\label{lem:preserva-supremos}\cite[Lema~4.2]{GehrkeHarding2001}
Si $m$ preserva supremos finitos, entonces $m^{\sigma}$ preserva supremos arbitrarios.
Dualmente, si $m$ preserva ínfimos finitos, entonces $m^{\pi}$ preserva ínfimos arbitrarios.
\end{lemma}

A partir de estos lemas podemos plantear el siguiente teorema:

\begin{theorem}\label{thm:suavidad-generalizada}
Sean $\mathbb{L},\mathbb{M}$ retículos distributivos acotados, 
$e_\mathbb{L},e_\mathbb{M}$ sus extensiones canónicas y 
$m\colon\mathbb{L}\to\mathbb{M}$ un operador monótono.
Si $m$ preserva supremos finitos, entonces $m^{\sigma}=m^{\pi}$.
Dualmente, si $m$ preserva ínfimos finitos, entonces $m^{\sigma}=m^{\pi}$.
\end{theorem}

\begin{proof}[\textbf{Demostración:}]
Del Lema~\ref{lem:preserva-supremos} se deduce que $m^{\sigma}$ preserva supremos arbitrarios.
Por el Lema~\ref{lem:igualdad-por-preservacion}, se sigue que $m^{\sigma}=m^{\pi}$.
El caso dual es análogo.
\end{proof}

De este modo, cuando $m$ preserva un tipo de finitud (supremos o ínfimos) —como sucede en los operadores normales de tipo $(1,1)$ o $(-1,-1)$—,
ambas extensiones coinciden. En tales casos escribiremos simplemente $m^{\delta}$ para designar la \emph{extensión canónica} de $m$.

A partir de este marco general, en la siguiente sección particularizaremos la construcción 
a los \emph{operadores normales de tipo} $(i,j)$ introducidos en el Capítulo~\ref{Ch 3}. 
Dado un operador
\[
m:\mathbb{L}^i \longrightarrow \mathbb{L}^j,
\]
su extensión canónica será el morfismo
\[
m^{\delta}:(\mathbb{L}^i)^\delta \longrightarrow (\mathbb{L}^j)^\delta.
\]
la cual permitirá conectar algebraicamente las relaciones inducidas en el espacio dual de Priestley.


\subsection{Caso particular: la extensión \(\sigma_{\mathbb{L}}\) sobre 
\(\Up(\mathcal{X}(\mathbb L))\) y operadores \((i,j)\)}\label{extension-operadores-sigma}

En el marco de la sección anterior, estudiamos la extensión canónica
$m^\sigma,m^\pi\colon \mathbb L^\delta\to\mathbb M^\delta$ de un
operador monótono $m\colon\mathbb L\to\mathbb M$.
Recordemos ahora que, de acuerdo con la Observación~\ref{rem:interpretacion-ij},
todo operador normal de tipo $(i,j)$ puede verse como un 
\(\vee\)-homomorfismo
\[
  m\colon \mathbb{L}^i \longrightarrow \mathbb{L}^j.
\]
En particular, un \(\vee\)-homomorfismo es automáticamente monótono, de modo que
la teoría general de extensiones canónicas desarrollada en la sección anterior
puede aplicarse directamente en este caso.

\smallskip

Para cada índice\(i\in\{-1,1\}\), denotamos por
\[
  \sigma_{\mathbb{L}^i}\colon
  \mathbb{L}^i\longrightarrow
  \Up(\mathcal{X}(\mathbb{L}^i))
\]
la función de Stone asociada al retículo \(\mathbb{L}^i\), definida por
\[
  \sigma_{\mathbb{L}^i}(a)
  =\{\,P\in\mathcal{X}(\mathbb{L}^i): a\in P\,\}.
\]
Obsérvese que esta aplicación depende del orden y de la topología de 
\(\mathbb{L}^i\); por lo tanto, \(\sigma_{\mathbb L}\) y 
\(\sigma_{\mathbb L^{-1}}\) no coinciden.

\smallskip

Aplicando ahora la construcción general de la extensión canónica al caso
particular en que los operadores actúan sobre los retículos de conjuntos
crecientes
\(
\Up(\mathcal{X}(\mathbb{L}^i))
\)
y
\(
\Up(\mathcal{X}(\mathbb{L}^j)),
\)
obtenemos que las extensiones canónicas de un operador normal de tipo $(i,j)$,
es decir de un \(\vee\)-homomorfismo
\[
m\colon \mathbb L^i \longrightarrow \mathbb L^j,
\]
adoptan en este contexto la siguiente forma explícita. Dados $(i,j)\in\{-1,1\}^2$ y
\(m\colon \mathbb L^i \to \mathbb L^j\)
un operador normal de tipo \((i,j)\).
Definimos sus extensiones
\[
  m^{\sigma},\,m^{\pi} :
  \Up(\mathcal{X}(\mathbb{L}^i))
  \longrightarrow
  \Up(\mathcal{X}(\mathbb{L}^j))
\]
por
\[
  m^{\sigma}(U)
  \;=\;
  \bigcup_{\Up(\mathcal{X}(\mathbb L^j))}
  \Bigl\{
    \bigcap_{\Up(\mathcal{X}(\mathbb L^j))}
      \{\sigma_{\mathbb L^j}(m(a)):\;C\subseteq^{i}\sigma_{\mathbb L^i}(a)\}
    \;\Big|\;
    C\in\mathcal C({\Up(\mathcal{X}(\mathbb L^i))}),\;
    C\subseteq^{i} U
  \Bigr\},
\]
\[
  m^{\pi}(U)
  \;=\;
  \bigcap_{\Up(\mathcal{X}(\mathbb L^j))}
  \Bigl\{
    \bigcup_{\Up(\mathcal{X}(\mathbb L^j))}
      \{\sigma_{\mathbb L^j}(m(a)):\;\sigma_{\mathbb L^i}(a)\subseteq^{i} O\}
    \;\Big|\;
    O\in\mathcal O(\Up(\mathcal{X}(\mathbb L^i))),\;
    U\subseteq^{i} O
  \Bigr\}.
\]

\noindent
Recordemos, teniendo en cuenta la Definición \ref{def:orden-i}, que el símbolo $\subseteq^i$ denota inclusión directa si $i=1$
e inclusión inversa si $i=-1$.

Esta formulación explicita las dos representaciones de Stone
involucradas: una para el dominio \(\mathbb{L}^i\) y otra para el codominio \(\mathbb{L}^j\). 
Los signos determinan, en cada caso, si se trabaja sobre
\(\Up(\mathcal{X}(\mathbb{L}))\) o sobre su dual
\(\Up(\mathcal{X}(\mathbb{L}^{-1}))\). 

\begin{theorem}\label{thm:extiende-en-basicos-ij-directo}
Sea $(i,j)\in\{-1,1\}^2$ y
\(m\colon\mathbb L^i\to\mathbb L^j\) un operador normal.
Entonces, para todo \(a\in L\),
\[
  m^{\sigma}\bigl(\sigma_{\mathbb{L}^i}(a)\bigr)
  \;=\;
  \sigma_{\mathbb{L}^j}\bigl(m(a)\bigr)
  \;=\;
  m^{\pi}\bigl(\sigma_{\mathbb{L}^i}(a)\bigr).
\]
\end{theorem}

\begin{proof}[\textbf{Demostración:}]
Sea \(U=\sigma_{\mathbb{L}^i}(a)\).
Dado que \(\sigma_{\mathbb{L}^i}(a)\) es clopen, es simultáneamente cerrado y abierto
en \(\Up(\mathcal{X}(\mathbb{L}^i))\).
Aplicando los Lemas~\ref{lem:lambda-en-cerrado-general}
y~\ref{lem:pi-en-abierto-general} obtenemos
\[
  m^{\sigma}(\sigma_{\mathbb{L}^i}(a))
  = \bigcap\{
     \sigma_{\mathbb{L}^j}(m(b))
     :\;\sigma_{\mathbb{L}^i}(a)\subseteq^{i}\sigma_{\mathbb{L}^i}(b)
   \}.
\]
Como \(\sigma_{\mathbb{L}^i}(a)\subseteq^{i}\sigma_{\mathbb{L}^i}(b)\)
si y sólo si \(a\le^{i} b\),
la intersección puede escribirse
\[
  \bigcap_{b\ge^{i} a} \sigma_{\mathbb{L}^j}(m(b)).
\]
Por monotonía de \(m\) se cumple
\(m(a)\le^{j} m(b)\) para todo \(b\ge^{i}a\),
y entonces
\(\sigma_{\mathbb{L}^j}(m(a))
 \subseteq^{j}
 \sigma_{\mathbb{L}^j}(m(b))\).
De ello se deduce
\[
  m^{\sigma}(\sigma_{\mathbb{L}^i}(a))
  = \sigma_{\mathbb{L}^j}(m(a)).
\]
El argumento dual, aplicado a \(m^{\pi}\),
da
\[
  m^{\pi}(\sigma_{\mathbb{L}^i}(a))
  = \bigcup_{b\le^{i} a}\sigma_{\mathbb{L}^j}(m(b))
  = \sigma_{\mathbb{L}^j}(m(a)).
\]
\end{proof}

\begin{theorem}\label{teo:suavidad-ij}
Sea $(i,j)\in\{-1,1\}^2$ y
\(m\colon\mathbb L^i\to\mathbb L^j\)
un operador normal de tipo \((i,j)\).
Si \(m\) preserva supremos (o ínfimos) finitos, entonces
\[
  m^{\sigma} \;=\; m^{\pi}.
\]
En consecuencia, existe una única extensión canónica
\[
  m^\delta:
  \Up(\mathcal{X}(\mathbb{L}^i))
  \longrightarrow
  \Up(\mathcal{X}(\mathbb{L}^j)).
\]
\end{theorem}

\begin{proof}[\textbf{Demostración:}]
La afirmación es consecuencia directa del
Teorema~\ref{thm:suavidad-generalizada}.
Los operadores normales de tipo \((1,1)\)
preservan supremos finitos, y los de tipo \((-1,-1)\)
preservan ínfimos finitos.
En el caso mixto \((1,-1)\), al ver a 
\(m\colon\mathbb L\to\mathbb L^{-1}\) como un $\vee$-homomorfismo,
por lo que los supremos en el dominio se corresponden con supremos en el codominio,
\[
  m(a\vee_{\mathbb L}b)
= m(a)\wedge_{\mathbb L}m(b)
= m(a)\vee_{\mathbb L^{-1}}m(b)
\]
de modo que \(m\) preserva supremos.
El caso \((-1,1)\) es dual.
\end{proof}

\begin{remark}
En los casos en que \(m^{\sigma}=m^{\pi}\),
escribiremos simplemente \(m^{\delta}\) para denotar la
extensión canónica única.
No obstante, conservaremos las notaciones
\(m^{\sigma}\) y \(m^{\pi}\)
cuando sea necesario referirse explícitamente
a sus definiciones en términos de abiertos y cerrados de $\mathbb{L}^{\delta}$.
\end{remark}








\section{La relación dual inducida por la extensión canónica}

El objetivo de esta sección es establecer, para los distintos tipos de
operadores normales $(i,j)\in\{-1,1\}^2$, el vínculo entre la extensión
canónica y la relación que caracteriza al funtor dual introducido en el
Capítulo~\ref{Ch 3}.
Mostraremos que, en todos los casos, la extensión canónica induce de manera
natural la misma relación $R_m$ que se utiliza en la dualidad tipo Priestley desarrollada en dicho Capítulo.

\vspace{0.2cm}

La técnica que utilizaremos en cada caso será la siguiente: analizaremos la acción puntual del
operador extendido $m^{\sigma}$ sobre los cerrados (o abiertos) de la extensión
canónica, describiendo su comportamiento mediante una multirrelación, la cual relaciona filtro primos con abiertos o cerrados de la extensión canónica, y luego identificaremos dicha
multirrelación con la relación $R_m$ correspondiente en cada caso.

Comenzaremos entonces por definir que entenderemos por una multirrelación.

\begin{definition}\label{def:multirrelacion}
Sea $X$ un conjunto.
Una \emph{multirrelación} sobre $X$ es un subconjunto
\[
N \subseteq X \times \mathcal{P}(X),
\]
\end{definition}
A continuación estudiaremos cada caso por separado.
Comenzaremos con el tipo $(1,1)$, que servirá de modelo general, pues
los demás casos se obtendrán a partir de él mediante dualizaciones o
simetrías apropiadas.

\subsection*{El caso $(1,1)$}

Comenzaremos estudiando cómo el operador $m^\sigma$ se define sobre un cerrado $C\in\mathcal{C}(\mathbb{L}^\delta)$.  Por la definición particular de $m^\sigma$ sobre un cerrado algebraico para el caso $(1,1)$, se tiene que:
\[
m^\sigma(C)=\bigcap_{\Up(\mathcal{X}(\mathbb L))}\{\sigma_{\mathbb L}(m(a)):C\subseteq
\sigma_{\mathbb L}(a)\}.
\]
Por lo tanto, para un filtro primo $P$, se tiene que:
\[
P\in m^\sigma(C)
\;\Longleftrightarrow\;
\forall\,a\in L,\ 
\bigl(C\subseteq\sigma_{\mathbb L}(a)\ \Rightarrow\ m(a)\in P\bigr).
\]
Por la caracterización del Teorema~\ref{cor:caracterizacion-C-F},
podemos reescribir la condición anterior como
\[
P\in m^\sigma(C)
\;\Longleftrightarrow\;
\forall a\in\Phi(C),\ m(a)\in P.
\]
Se sigue entonces que
\[
\Phi(C)\subseteq m^{-1}[P],
\]
o, de manera equivalente,
\[
\Phi(C)\cap m^{-1}[P]^c=\varnothing.
\]


Esta equivalencia sugiere que la pertenencia
$P\in m^\sigma(C)$ puede expresarse de forma puramente relacional,
definiendo la siguiente multirrelación.

\begin{definition}
\label{def:multirrelacion-inducida-(1,1)}
Sean $\mathbb{L}$ un retículo distributivo acotado y $m : \mathbb{L} \to \mathbb{L}$ un operador normal de tipo $(1,1)$. La \emph{multirrelación inducida por $m$} se define de la siguiente manera:
\[
(P,C) \in N_m 
\iff 
\Phi(C) \cap m^{-1}[P]^c = \emptyset.
\]
\end{definition}

Recordemos que, por el Lema \ref{Ideales}, $m^{-1}[P]^c$ es un ideal, puesto que $m$ es un operador de tipo $(1,1)$ y $P$ es un filtro primo. De este modo, la condición anterior
expresa que el filtro $\Phi(C)$ y el ideal $m^{-1}[P]^c$ son disjuntos.

\begin{remark}\label{Obs:multi-msigma}
La equivalencia antes presentada puede leerse ahora así:
\[
P\in m^\sigma(C)
\quad\text{si y sólo si}\quad
(P,C)\in N_m.
\]
Es decir, el operador canónico $m^\sigma$ queda determinado por la multirrelación $N_m$.
\end{remark}
En lo que sigue, vamos a identificar esta multirrelación con la relación dual $R_m$ introducida en el Capítulo~\ref{Ch 3}. Recordemos que dicha relación estaba definida como $(P,Q)\in R_m\iff Q^i\subseteq m^{-1}[P]^j$. Es decir que para este caso particular, tenemos que  $(P,Q)\in R_m\iff Q\subseteq m^{-1}[P]$.

\begin{lemma}\label{lem:N-R}
Sea \(m\) un operador de tipo \((1,1)\). 
Para todo \(P\in\mathcal{X}(\mathbb L)\) y todo \(C\in\mathcal C(\mathbb L^\delta)\) se verifican las siguientes proposiciones equivalentes:
\begin{enumerate}
    \item $(P,C)\in N_m$
    \item $\exists\, Q\in C \text{ tal que } (P,Q)\in R_m$
    \item $(P,[Q))\in N_m$
\end{enumerate}
donde \([Q)=\{R\in \mathcal{X}(\mathbb L): Q\subseteq R\}\) es el creciente principal de \(Q\).
\end{lemma}


\begin{proof}[\textbf{Demostración:}]
Probaremos $(1)\Rightarrow(2)\Rightarrow (3)\Rightarrow (1)$:
\begin{itemize}
    \item $(1)\Rightarrow(2)$:
    
Supongamos $(P,C)\in N_m$.  
Por definición, $\Phi(C)\subseteq m^{-1}[P]$, es decir,
$\Phi(C):=\Phi(C)$ es un filtro tal que
$\Phi(C)\cap m^{-1}[P]^c=\emptyset$.
Por el Lema~\ref{Ideales}, $m^{-1}[P]^c$ es un ideal.
Aplicando el teorema del filtro primo, existe un filtro primo
$Q\in\mathcal{X}(\mathbb L)$ con
\[
\Phi(C)\subseteq Q
\qquad\text{y}\qquad
Q\cap m^{-1}[P]^c=\emptyset.
\]
De la segunda condición se deduce $Q\subseteq m^{-1}[P]$,
es decir $(P,Q)\in R_m$.
Como $\Phi(C)\subseteq Q$, por el Corolario~\ref{cor:caracterizacion-C-F} se tiene $Q\in C$.
Hemos probado así $(P,C)\in N_m\Rightarrow\exists\,Q\in C$ tal que $(P,Q)\in R_m$.

\item $(2)\Rightarrow (3)$:

Sea $Q\in C$ tal que $(P,Q)\in R_m$, o sea $Q\subseteq m^{-1}[P]$.
Consideremos el creciente principal de $Q$,
\[
[Q)=\{R\in\mathcal{X}(\mathbb L):Q\subseteq R\}.
\]
Por la definición de la correspondencia $\Phi$, se tiene
$\Phi([Q))=Q$.
Como $Q\subseteq m^{-1}[P]$, resulta entonces
$\Phi([Q))\subseteq m^{-1}[P]$, 
y en consecuencia $(P,[Q))\in N_m$.

\item $(3)\Rightarrow (1)$: 

Si $(P,[Q))\in N_m$, entonces 
$\Phi([Q))=Q\subseteq m^{-1}[P]$.
Puesto que $Q\in [Q)\subseteq C$,
se tiene también $\Phi(C)\subseteq Q\subseteq m^{-1}[P]$,
lo que muestra $(P,C)\in N_m$.
\end{itemize}
\end{proof}
El lema anterior muestra que, gracias a las propiedades de normalidad del operador $(1,1)$, la multirrelación $N_m$ coincide punto a punto con la relación $R_m$ introducida en el
Capítulo~\ref{Ch 3}.

\vspace{0.2cm}

Para finalizar esta sección, vamos a comprobar que el operador $m^{\sigma}$ extiende de buena manera a el operador $m_{R_{m}}$ definido en la dualidad de tipo Priestley (ver Capitulo \ref{Ch 3}) i.e., ambos coinciden sobre los clopen. Recordemos que este operador estaba definido sobre los clopens como $m_{R_m}(U)=R_m^{-1}(U^i)^j$, es decir que en este caso particular, $m_{R_m}(U)=R_m^{-1}(U)$. Antes de probar dicho resultado, vamos a probar un lema previo que será útil.

\begin{lemma}\label{lem:cerrado-vs-punto}
Sean $\mathbb{L}$ un retículo distributivo acotado, $(\mathbb{L}^{\delta},\sigma_{\mathbb{L}})$ su extension canónica, $U$ un clopen de $\mathbb{L}^{\delta}$ y $R\subseteq\mathcal X(\mathbb L)\times\mathcal X(\mathbb L)$ una relación binaria. Entonces, para todo $P\in\mathcal X(\mathbb L)$, son equivalentes:
\begin{itemize}
\item[1] Existe un cerrado algebraico $C \subseteq U$ y un filtro primo $Q \in C$ tal que $(P,Q) \in R$
\item[2] Existe un filtro primo $Q \in U$ tal que $(P,Q) \in R$
\end{itemize}
\end{lemma}

\begin{proof}[\textbf{Demostración:}]
Probaremos ambas implicaciones:
\begin{itemize}
\item  $ (1) \Rightarrow (2)$ Si existen $C\subseteq U$ cerrado algebraico y un punto $Q\in C$ con $(P,Q)\in R$,
entonces necesariamente $Q\in U$.

\item $(2)\Rightarrow (1)$ Recíprocamente, si existe $Q\in U$ con $(P,Q)\in R$,  
podemos tomar simplemente $C=U$.  
Como $U$ es clopen (y en particular cerrado),
se tiene $C\subseteq U$, $Q\in C$ y $(P,Q)\in R$
\end{itemize}
\end{proof}

A continuación, vamos a comprobar que los operadores \(m^\sigma\) y \(m_{R_m}\) coinciden sobre los clopen de la extensión canónica de cualquier retículo distributivo acotado dado.


\begin{theorem}\label{teo:igualdad-operadores}
Sean $\mathbb{L}$ un retículo distributivo acotado, $m$ un operador de tipo $(1,1)$ definido sobre $\mathbb{L}$ y $U$ un clopen de la extensión canónica. Entonces, para todo
$P\in\mathcal X(\mathbb L)$,
\[
P\in m^\sigma(U)
\quad\Longleftrightarrow\quad
P\in m_{R_m}(U).
\]
\end{theorem}

\begin{proof}[\textbf{Demostración:}]
Sea $P\in\mathcal X(\mathbb L)$. Vamos a comprobar que, para cada $U$ clopen de $\mathbb{L}^{\delta}$,
se verifica:
\[
P\in m^\sigma(U)
\quad\Longleftrightarrow\quad
P\in m_{R_m}(U).
\]
Recordemos que $m^\sigma$ se define a partir de
los cerrados algebraicos $C\subseteq U$. Por lo tanto, se tiene que
\begin{align*}
P \in m^\sigma(U)
&\Longleftrightarrow 
\exists\,C\subseteq U \text{ cerrado tal que } P\in m^\sigma(C)\\
&\Longleftrightarrow 
\exists\,C\subseteq U \text{ tal que } (P,C)\in N_m & \text{(Observación \ref{Obs:multi-msigma})}\\
&\Longleftrightarrow 
\exists\,C\subseteq U,\,\exists\,Q\in C \text{ con } (P,Q)\in R_m & \text{(Lema \ref{lem:N-R})}\\
&\Longleftrightarrow 
\exists\,Q\in U \text{ tal que } (P,Q)\in R_m & \text{(Lema \ref{lem:cerrado-vs-punto})} \\
&\Longleftrightarrow 
P\in R_m^{-1}(U)=m_{R_m}(U).
\end{align*}
En consecuencia, $m^\sigma$ y $m_{R_m}$ definen el mismo operador sobre los clopens de la extensión canónica.
\end{proof}

Es importante mencionar que el operador $m^\sigma$ está definido sobre todo el retículo $\mathbb{L}^{\delta}$, sin embargo, coincide con $m_{R_m}$ sobre los clopens de la extensión canónica.

\subsection*{El caso $(1,-1)$}

Consideremos ahora el caso en donde $(\mathbb{L},m)$ es un retículo distributivo acotado y $m$ es un operador normal de tipo $(1,-1)$. Se sigue entonces que $m$ es un operador que en $\mathbb L$ que es antimonótono y satisface  $m(a\vee b)=m(a)\wedge m(b)$ y $m(0)=1$.

\vspace{0.2cm}

Nuestro objetivo será análogo al caso $(1,1)$, a saber: describir $m^\sigma$ sobre la extensión canónica en términos relacionales y comprobar que $m^{\sigma}$ es igual al operador inducido por la relación dual $R_m$.

\vspace{0.2cm}

Comenzaremos esta sección con el siguiente resultado.

\begin{lemma}\label{lem:equivalencia-sigma-op} Sean $\mathbb L$ un retículo distributivo acotado y
$P\in\mathcal X(\mathbb L)$. Entonces para todo $a\in L$ se verifica la siguiente igualdad:
\[
P\in \sigma_{\mathbb L^{-1}}(a)
\quad\Longleftrightarrow\quad
P^c \in \sigma_{\mathbb L}(a)^c.
\]
\end{lemma}

\begin{proof}[\textbf{Demostración:}]
Por definición de la aplicación de Stone, tanto sobre $\mathbb{L}$ como sobre $\mathbb{L}^{-1}$:
\begin{align*}
P\in\sigma_{\mathbb L^{-1}}(a)
& \Longleftrightarrow a\in P \text{ y } P \in \mathcal{X}(\mathbb L^{-1})\\
& \Longleftrightarrow a\notin P^c \text{ y } P^c \in \mathcal{X}(\mathbb L)\qquad\text{(Lema \ref{lem:primos-op-complemento})}\\
& \Longleftrightarrow P^c\in\sigma_{\mathbb L}(a)^c
\end{align*}
\end{proof}
A continuación, daremos una descripción del operador $m^\sigma$ sobre los elementos cerrados de $\mathbb{L}^{\delta}$.

\vspace{0.2cm}

Sea $C\in\mathcal C(\mathbb L^\delta)$ un cerrado y $P\in\mathcal X(\mathbb L)$. Por la definición particular de las extensiones dada en la subsección \ref{extension-operadores-sigma} y por el Lema \ref{lem:lambda-en-cerrado-general}, se tiene que
\[
m^\sigma(C)=
\bigcap_{\Up(\mathcal X(\mathbb L^{-1}))}
\{\ \sigma_{\mathbb L^{-1}}(m(a)):\ C\subseteq\sigma_{\mathbb L}(a)\ \}.
\]
Notemos que, la operación de ínfimo arbitrario está siendo tomada sobre el retículo $\left(\mathbb{L}^{\delta}\right)^{-1}$. Por lo tanto, expresando dicha operación en términos de $\mathbb{L}^{\delta}$, se tiene que
\[
m^\sigma(C)=
\bigcup_{\Up(\mathcal X(\mathbb L))}
\{\ \sigma_{\mathbb L^{-1}}(m(a)):\ C\subseteq\sigma_{\mathbb L}(a)\ \}.
\]

Por tanto, un punto $P$ pertenece a $m^\sigma(C)$ si
\[
P\in m^\sigma(C)
\;\Longleftrightarrow\;
\exists\,a\in L\ \text{tal que } 
C\subseteq\sigma_{\mathbb L}(a)\ \text{y}\ 
P\in\sigma_{\mathbb L^{-1}}(m(a)).
\]

Notemos que, la condicion $C\subseteq\sigma_{\mathbb L}(a)$ puede intepretarse por medio del filtro asociado $\Phi(C)$ (Teorema~\ref{cor:caracterizacion-C-F}), obteniendo
\[
P\in m^\sigma(C)
\;\Longleftrightarrow\;
\exists\,a\in\Phi(C)\ \text{tal que } 
P\in\sigma_{\mathbb L^{-1}}(m(a)).
\]

Teniendo en cuenta el Lema~\ref{lem:equivalencia-sigma-op}, se sigue que
\[
P\in\sigma_{\mathbb L^{-1}}(m(a))
\;\Longleftrightarrow\;
P^c\in\sigma_{\mathbb L}(m(a))^c
\;\Longleftrightarrow\;
m(a)\notin P^c,
\]

de donde se sigue que:
\[
P\in m^\sigma(C)
\;\Longleftrightarrow\;
\exists\,a\in\Phi(C)\ \text{tal que }
m(a)\notin P^c.
\]

Finalmente, como $m^{-1}[P]=\{\,a\in L:m(a)\in P\,\}$ y 
$P^c$ es un ideal, la condición 
$m(a)\notin P^c$ es equivalente a $a\in m^{-1}[P]$.  
De este modo llegamos a la siguiente formulación:
\[
P\in m^\sigma(C)
\;\Longleftrightarrow\;
\exists\,a\in\Phi(C)\ \text{tal que } a\in m^{-1}[P]
\;\Longleftrightarrow\;
\Phi(C)\cap m^{-1}[P]\neq\varnothing.
\]


\noindent
La última equivalencia sintetiza la idea central:
en el caso $(1,-1)$ la pertenencia de $P$ a $m^\sigma(C)$
se expresa mediante la existencia de un elemento del filtro asociado
$\Phi(C)$ que no es excluido por el ideal $m^{-1}[P]$.
De este modo, la relación entre $P$ y $C$ refleja la
antimonotonía (en $\mathbb{L}$)  del operador de tipo $(1,-1)$.
En consecuencia,
\[
P\in m^\sigma(C)
\quad\Longleftrightarrow\quad
\Phi(C)\cap m^{-1}[P]\neq\emptyset,
\]
donde
\(
m^{-1}[P]
\) es un ideal por el Lema~\ref{Ideales} aplicado al tipo $(1,-1)$.

\medskip

Motivados por la caracterización anterior, definimos:

\begin{definition}\label{def:Nm-1-1}
Sean $\mathbb{L}$ un retículo distributivo acotado y $m:\mathbb L\to\mathbb L^{-1}$ un operador normal de tipo $(1,-1)$. Definimos la multirrelación:
\[
(P,C)\in N_m
\quad\Longleftrightarrow\quad
\Phi(C)\cap m^{-1}[P]=\emptyset,
\]
para $P\in\mathcal X(\mathbb L)$ y $C\in\mathcal C(\mathbb L^\delta)$.
\end{definition}
\noindent
\begin{remark} \label{obs:C-noN}
Obsérvese que esta definición reproduce la misma multirrelación
introducida en el caso $(1,1)$, pero aquí la condición de pertenencia
a $m^\sigma(C)$ se formula en términos de los pares que \emph{no}
pertenecen a $N_m$. Con esta notación, para todo $C\in\mathcal C(\mathbb L^\delta)$ y $P\in\mathcal X(\mathbb L)$ se tiene
\[
P\in m^\sigma(C)
\quad\Longleftrightarrow\quad
(P,C)\notin N_m.
\]
\end{remark}

Nuevamente, el paso siguiente consiste en identificar esta multirrelación con la relación dual $R_m$ introducida en el Capítulo~\ref{Ch 3}. Recordemos que dicha relación, para este caso particular, está definida como $(P,Q)\in R_m\iff Q\subseteq m^{-1}[P]^c$.

\begin{lemma}\label{lem:N-R-1-1}
Sean $\mathbb{L}$ un retículo distributivo acotado y \(m\) un operador de tipo \((1,-1)\). Entonces, para todo \(P\in\mathcal X(\mathbb L)\) y todo \(C\in\mathcal C(\mathbb L^\delta)\), las siguiente afirmaciones son equivalentes:
\begin{enumerate}
    \item $(P,C)\notin N_m$.
    \item $\forall\, Q\in C: (P,Q)\notin R_m$.
\end{enumerate}
\end{lemma}

\begin{proof}[\textbf{Demostración:}]

$(1)\Rightarrow(2)$ Supongamos $(P,C)\notin N_m$. Por definición de $N_m$, se sigue que 
\[ \Phi(C)\cap m^{-1}[P]\neq\emptyset. \]
Sea $a\in \Phi(C)\cap m^{-1}[P]$. Como $a\in \Phi(C)$, se tiene que $C\subseteq\sigma_{\mathbb L}(a)$.
Ahora bien, para cualquier $Q\in C$, de $a\in m^{-1}[P]$ se sigue $m(a)\in P$,
y por la definición de la relación dual para el tipo $(1,-1)$,
\[
(P,Q)\notin R_m
\quad\Longleftrightarrow\quad
Q\not\subseteq m^{-1}[P]^c
\quad\Longleftrightarrow\quad
\exists\,a\in Q\cap m^{-1}[P].
\]
Pero precisamente $a\in Q\cap m^{-1}[P]$, ya que
$C\subseteq\sigma_{\mathbb L}(a)$ implica $a\in Q$.
Por tanto, $(P,Q)\notin R_m$ para todo $Q\in C$.\\

\vspace{0.2cm}

$(2)\Rightarrow(1)$ Supongamos que $\forall\,Q\in C,\ (P,Q)\notin R_m$.
Entonces para cada $Q\in C$ existe $a_Q\in Q\cap m^{-1}[P]$.
Consideremos el conjunto
\[
F=\{\,a_Q : Q\in C\,\}.
\]
Como cada $a_Q$ pertenece a $\Phi(C)$ (pues $a_Q\in Q$ para todo $Q\in C$)
y a $m^{-1}[P]$, se concluye que
$F\subseteq \Phi(C)\cap m^{-1}[P]$.
En particular, la intersección $\Phi(C)\cap m^{-1}[P]$ es no vacía,
y por tanto $(P,C)\notin N_m$.
\end{proof}
Veamos por último que lo operadores $m^\sigma$ y $m_{R_m}$ coinciden sobre los clopens de la extensión canónica. Recordemos que en este caso particular, si $U$ es un clopen de la extensión canónica, $m_{R_m}(U)=R_m^{-1}(U)^c$.
\begin{theorem} \label{teo:igualdad-operadores1-1}
Sean $\mathbb{L}$ un retículo distributivo acotado, $m$ un operador de tipo $(1,-1)$ definido sobre $\mathbb{L}$ y $U$ un clopen de la extensión canónica. Entonces, para todo
$P\in\mathcal X(\mathbb L)$,
\[
P\in m^\sigma(U)
\quad\Longleftrightarrow\quad
P\in m_{R_m}(U).
\]
\end{theorem}

\begin{proof}[\textbf{Demostración:}]
Concatenamos las equivalencias clave, como en el caso $(1,1)$. Recordemos que, en este caso, $m^\sigma$ se define como una
intersección (indexada por los cerrados contenidos en $U$);
por tanto, la pertenencia $P\in m^\sigma(U)$ se traduce en una
condición universal sobre los cerrados $C\subseteq U$:
\begin{align*}
P \in m^\sigma(U)
&\Longleftrightarrow 
\forall\,C\subseteq U \text{ cerrado}, P\in m^\sigma(C)\\
&\Longleftrightarrow 
\forall\,C\subseteq U: (P,C)\notin N_m & \text{(Observación~\ref{obs:C-noN})}\\
&\Longleftrightarrow 
\forall\,C\subseteq U,\,\forall\,Q\in C: (P,Q)\notin R_m
& \text{(Lema~\ref{lem:N-R-1-1})}\\
&\Longleftrightarrow \forall\,Q\in U: (P,Q)\notin R_m & \text{(Lema~\ref{lem:cerrado-vs-punto})}\\
&\Longleftrightarrow 
R_m(P)\subseteq U^c\\
&\Longleftrightarrow 
U\subseteq R_m(P)^c\\
&\Longleftrightarrow 
P\in R_m^{-1}(U)^c=m_{R_m}(U).
\end{align*}
En consecuencia, $m^\sigma$ y $m_{R_m}$ definen el mismo operador sobre los clopens de la extensión canónica.
\end{proof}
\noindent

\subsection*{El caso $(-1,-1)$}

En esta subsección consideraremos el caso en donde $(\mathbb{L},m)$ es un retículo distributivo acotado y $m$ es un operador normal de tipo $(-1,-1)$. Sin embargo aquí no utilizaremos la Observación \ref{rem:interpretacion-ij}, sino que tomaremos a
\[
  m \colon \mathbb{L} \longrightarrow \mathbb{L}
\]
es decir, un operador que es preserva ínfimos finitos y el elemento superior de $\mathbb{L}$, y por lo tanto es monótono. Es decir que la extensión canónica que podemos definir está dada de $\Up(\mathcal{X}(\mathbb{L}))$ a $\Up(\mathcal{X}(\mathbb{L}))$. Además, recordemos que por el Teorema~\ref{teo:suavidad-ij}, la extensión canónica de $m$ satisface
$m^\sigma = m^\pi$, de modo que podemos trabajar indistintamente con
una u otra. En lo que sigue utilizaremos $m^\pi$, ya que su definición
se formula directamente en términos de abiertos, lo que permite
describirla mediante ideales de $\mathbb{L}$.

Comencemos observando cómo el operador $m^\pi$ actúa sobre un abierto $O\in\mathcal{O}(\mathbb{L}^\delta)$ y un punto $P\in\mathcal{X}(\mathbb L)$.
Por la definición particular de $m^\pi$ sobre un abierto para el caso $(-1,-1)$:
\[
m^\pi(O)=\bigcup_{\Up(\mathcal{X}(\mathbb L))}\{\sigma_{\mathbb L}(m(a)):\sigma_{\mathbb L}(a)\subseteq O\}.
\]
Por lo tanto, la pertenencia puntual se expresa como
\[
P\in m^\pi(O)
\;\Longleftrightarrow\;
\exists\,a\in L\ 
\bigl(\sigma_{\mathbb L}(a)\subseteq O
\ \text{ y }\ 
P\in\sigma_{\mathbb L}(m(a))\bigr).
\]

Usando ahora la caracterización del Teorema~\ref{cor:caracterizacion-U-Gamma},
la condición $\sigma_{\mathbb L}(a)\subseteq O$ es equivalente a
$a\in\Gamma(O)$, donde $\Gamma(O)$ es el ideal asociado a $O$.
De este modo obtenemos
\[
P\in m^\pi(O)
\;\Longleftrightarrow\;
\exists\,a\in\Gamma(O)\ \text{con}\ m(a)\in P.
\]

Dicho de otra manera,
\[
\exists\,a\in\Gamma(O)\cap m^{-1}[P],
\]
lo que nos lleva a la caracterización fundamental:
\[
P\in m^\pi(O)
\quad\Longleftrightarrow\quad
\Gamma(O)\cap m^{-1}[P]\neq\varnothing.
\]
Definiendo así la siguiente multirrelación.

\begin{definition}
\label{def:multirrelacion-inducida-(-1,-1)}
Sean $\mathbb{L}$ un retículo distributivo acotado y $m : \mathbb{L} \to \mathbb{L}$ un operador normal de tipo $(-1,-1)$. Definimos la \emph{multirrelación inducida por $m$} como
\[
(P,O) \in M_m 
\iff 
\Gamma(O)\cap m^{-1}[P]\neq\varnothing
\]
\end{definition}

\noindent
\begin{remark}
Recordemos además del Lema~\ref{Ideales} que, si
$m$ es un operador de tipo $(-1,-1)$ y $P$ es un filtro primo, entonces $m^{-1}[P]$ es un filtro.  De este modo, la condición anterior expresa que el filtro $m^{-1}[P]$ y el ideal $\Gamma(O)$ no son disjuntos.
\end{remark}

\begin{remark}\label{obs:O-M}
La equivalencia anterior puede leerse ahora así:
\[
P\in m^\pi(O)
\quad\text{si y sólo si}\quad
(P,O)\in M_m.
\]
Es decir, el operador canónico $m^\pi$ queda determinado por la multirrelación $M_m$.
\end{remark}
Así, por tercera vez, el paso siguiente consiste en identificar esta multirrelación con la relación dual $R_m$ introducida en el Capítulo~\ref{Ch 3}. Recordemos que, al ser $m$ un operador de tipo \((-1,-1)\), $(P,Q)\in R_m\iff Q^c\subseteq m^{-1}[P]^c$ o lo que es lo mismo $m^{-1}[P]\subseteq Q$
\begin{lemma}\label{lem:M-R}
Sea \(m\) un operador de tipo \((-1,-1)\). 
Para todo \(P\in\mathcal{X}(\mathbb L)\) y todo \(O\in\mathcal O(\mathbb L^\delta)\) se verifican las siguientes proposiciones equivalentes:
\begin{enumerate}
    \item $(P,O)\in M_m$
    \item $\forall\, Q\in R_m(P) \Rightarrow Q\in O$
\end{enumerate}
\end{lemma}


\begin{proof}[\textbf{Demostración:}]
Probaremos las dos implicaciones:
\begin{itemize}
    \item $(1)\Rightarrow(2)$
    
Supongamos $(P,O)\in M_m$ y  sea  $Q\in R_m(P)$. Por definición de $M_m$, $\Gamma(O)\cap m^{-1}[P]\neq\varnothing$. Pero puesto que $Q\in R_m(P)$, es decir, $m^{-1}[P]\subseteq Q$, se cumple que $\Gamma(O)\cap Q\neq\varnothing$. Luego $Q\in O$

\item $(2)\Rightarrow (1)$

Supongamos (2) y que $(P,O)\notin M_m$ para llegar a un absurdo. Se debe cumplir así que $\Gamma(O)\cap m^{-1}[P]=\varnothing$. Como vimos, $\Gamma(O)$ es un abierto y $m^{-1}[P]$ un filtro, por lo que aplicando el Teorema del Filtro Primo existe un $Q\in\mathcal{X}(\mathbb{L})$ tal que $m^{-1}[P]\subseteq Q$ (es decir que $Q\in R_m(P)$) y $\Gamma(O)\cap Q=\varnothing$. Pero si $\Gamma(O)\cap Q=\varnothing$ entonces $Q\notin O$, lo cual es un absurdo.
\end{itemize}
\end{proof}
El lema anterior muestra que la multirrelación $M_m$ coincide punto a punto con la relación $R_m$. Antes de probar el teorema final que conecta los operadores \(m^\pi\) y \(m_{R_m}\), debemos mencionar un lema previo:
\begin{lemma}\label{lem:cerrado-interseccion-abiertos}
Sea $\mathbb{L}$ un retículo distributivo acotado y 
$(\mathbb{L}^{\delta},\sigma_{\mathbb{L}})$ su extensión canónica.
Sea $U\in\mathcal{C}(\mathbb{L}^\delta)$ un cerrado y
$S\subseteq \mathcal{X}(\mathbb{L})$ un subconjunto cualquiera. Entonces son equivalentes:
\begin{itemize}
  \item[1)] Para todo abierto $O\in\mathcal{O}(\mathbb{L}^\delta)$, si $U\subseteq O$ entonces $S\subseteq O$.
  \item[2)] $S\subseteq U$.
\end{itemize}
\end{lemma}

\begin{proof}[\textbf{Demostración:}]
Probaremos ambas implicaciones:
\begin{itemize}
\item  $(2)\Rightarrow(1)$ es inmediato: si $S\subseteq U$ y $U\subseteq O$ para algún
$O\in\mathcal{O}(\mathbb{L}^\delta)$, entonces $S\subseteq O$ por transitividad de la inclusión.

\item $(1)\Rightarrow(2)$.
Supongamos que $S\not\subseteq U$. Entonces existe $Q_{0}\in S$ tal que $Q_{0}\notin U$.
Como $U$ es un cerrado de $\mathbb{L}^\delta$, por el Teorema~\ref{thm:cerrados-filtros}
existe un filtro $F\in\mathrm{Fi}(\mathbb{L})$ tal que
\[
U \;=\; \bigcap_{a\in F}\sigma_{\mathbb{L}}(a).
\]
De $Q_{0}\notin U$ se sigue que
\[
Q_{0}\notin\bigcap_{a\in F}\sigma_{\mathbb{L}}(a),
\]
y por lo tanto existe $a_{0}\in F$ tal que $Q_{0}\notin\sigma_{\mathbb{L}}(a_{0})$,
es decir, $a_{0}\notin Q_{0}$.

Como $U=\bigcap_{a\in F}\sigma_{\mathbb{L}}(a)$, en particular
\[
U\;\subseteq\;\sigma_{\mathbb{L}}(a_{0}).
\]
Definimos entonces el abierto
\[
O:=\sigma_{\mathbb{L}}(a_{0})\in\mathcal{O}(\mathbb{L}^\delta).
\]
Se tiene $U\subseteq O$, pero $Q_{0}\notin O$ porque $a_{0}\notin Q_{0}$.
Por lo tanto, $S\nsubseteq O$.

Esto contradice la condición (1), según la cual para todo abierto $O$ con $U\subseteq O$
debería cumplirse $S\subseteq O$. Concluimos que nuestra suposición $S\not\subseteq U$
es imposible, y por lo tanto $S\subseteq U$.
\end{itemize}
\end{proof}


Ahora sí, veamos como los operadores \(m^\pi\) y \(m_{R_m}\) coinciden sobre los clopens de la extensión canónica. Recordemos que en este caso particular, si $U$ es un clopen de la extensión canónica, $m_{R_m}(U)=R_m^{-1}(U^c)^c$.
\begin{theorem} \label{teo:igualdad-operadores-1-1}
Sean $\mathbb{L}$ un retículo distributivo acotado, $m$ un operador de tipo $(-1,-1)$ definido sobre $\mathbb{L}$ y $U$ un clopen de la extensión canónica. Entonces, para todo
$P\in\mathcal X(\mathbb L)$,
\[
P\in m^\pi(U)
\quad\Longleftrightarrow\quad
P\in m_{R_m}(U).
\]
\end{theorem}

\begin{proof}[\textbf{Demostración:}]
Concatenamos las equivalencias clave, como en los casos previos.
\begin{align*}
P \in m^\pi(U)&\Longleftrightarrow \bigcap_{\Up(\mathcal{X}(\mathbb L))}\{m^\pi (O):U\subseteq O\}\\
&\Longleftrightarrow 
\forall\,O \in\mathcal O(\mathbb L^\delta) :U\subseteq O\Rightarrow P\in m^\pi(O)\\
&\Longleftrightarrow 
\forall\,O\in\mathcal O(\mathbb L^\delta) :U\subseteq O\Rightarrow (P,O)\in M_m & \text{(Observación \ref{obs:O-M}})\\
&\Longleftrightarrow 
\forall\,O\in\mathcal O(\mathbb L^\delta) :U\subseteq O\Rightarrow\forall Q\in R_m(P): Q\in O & \text{(Lema \ref{lem:M-R}})\\
&\Longleftrightarrow 
\forall Q\in R_m(P): Q\in U  & \text{(Lema \ref{lem:cerrado-interseccion-abiertos}})\\
&\Longleftrightarrow 
R_m(P)\subseteq U\\
&\Longleftrightarrow 
R_m(P)\cap U^c=\emptyset\\
&\Longleftrightarrow 
P\in R_m^{-1}(U^c)^c=m_{R_m}(U)
\end{align*}
En consecuencia, $m^\pi$ y $m_{R_m}$ definen el mismo operador sobre los clopens de la extensión canónica.
\end{proof}





\subsection*{El caso $(-1,1)$}

Consideremos un operador normal de tipo $(-1,1)$,
\[
m:\mathbb L^{-1} \longrightarrow \mathbb L,
\]
lo que significa que, interpretado en $\mathbb L$, el operador $m$ es
antimonótono y satisface
\[
m(a\wedge b)=m(a)\vee m(b), \qquad m(1)=0.
\]
Nuestro objetivo es describir la acción puntual de $m^\sigma$ sobre los cerrados de la extensión canónica, expresar dicha
acción en términos de una multirrelación, y finalmente identificarla con el operador relacional
$m_{R_m}$ asociado a la relación dual $R_m$ del tipo $(-1,1)$.

\medskip

En esta sección introducimos un operador auxiliar 
\((-)^{\dagger}\) que permite traducir condiciones sobre 
cerrados de $\Up(\mathcal X(\mathbb L^{-1}))=(\mathbb{L}^{-1})^\delta$ 
a condiciones equivalentes sobre abiertos de 
$\mathbb{L}^\delta$, y con ello expresar 
la caracterización puntual de $m^\sigma$ en términos 
de la relación dual $R_m$ (que está definida en $\mathcal X(\mathbb L)\times \mathcal X(\mathbb L)$). Lo definiremos y probaremos algunos sobre su aplicación que nos serán útiles.

\begin{definition}\label{def:operador-dagger}
Sea $U\subseteq\mathcal X(\mathbb L^{-1})$. 
Definimos
\[
U^{\dagger}
:= 
\{\,P^{c}\in\mathcal X(\mathbb L)\mid 
P\in\mathcal X(\mathbb L^{-1}),\ P\notin U\,\}.
\]
\end{definition}

\begin{lemma}
Sea $\mathbb L$ un retículo distributivo acotado y
$C$ un cerrado de $(\mathbb{L}^{-1})^\delta$.
Entonces $C^{\dagger}$ es un abierto de $\mathbb{L}^\delta$.
\end{lemma}

\begin{proof}[\textbf{Demostración:}]
Como $C$ es cerrado en la extensión canónica de $\mathbb L^{-1}$,
por el Teorema~\ref{thm:cerrados-filtros} aplicado al retículo
$\mathbb L^{-1}$, existe un filtro $F$ de $\mathbb L^{-1}$ tal que
\[
C
=
\bigcap_{a\in F}\sigma_{\mathbb L^{-1}}(a).
\]
Tomando complementos en $\mathcal X(\mathbb L^{-1})$ se obtiene
\[
\mathcal X(\mathbb L^{-1})\setminus C
=
\bigcup_{a\in F}\sigma_{\mathbb L^{-1}}(a)^{c}.
\]

Por definición de $(-)^{\dagger}$,
\[
C^{\dagger}
=
\{\,P^{c}\in\mathcal X(\mathbb L) : P\notin C\,\}
=
\{\,P^{c} : P\in\mathcal X(\mathbb L^{-1})\setminus C\,\},
\]
luego
\[
C^{\dagger}
=
\bigcup_{a\in F}\{\,P^{c} : P\notin\sigma_{\mathbb L^{-1}}(a)\,\}.
\]

Ahora usamos el Lema~\ref{lem:equivalencia-sigma-op}. Para cada $a\in L$,
\[
P\notin\sigma_{\mathbb L^{-1}}(a)
\;\Longleftrightarrow\;
P^{c}\in\sigma_{\mathbb L}(a),
\]
de donde
\[
\{\,P^{c} : P\notin\sigma_{\mathbb L^{-1}}(a)\,\}
=
\sigma_{\mathbb L}(a).
\]
Sustituyendo en la expresión anterior,
\[
C^{\dagger}
=
\bigcup_{a\in F}\sigma_{\mathbb L}(a).
\]

Observe que $F$, visto en $\mathbb L$, es un ideal
(pues es un filtro en $\mathbb L^{-1}$). Por el
Teorema~\ref{thm:abiertos-ideales}, aplicado ahora al retículo
$\mathbb L$, toda unión de la forma
$\bigcup_{a\in I}\sigma_{\mathbb L}(a)$ con $I$ ideal de $\mathbb L$
es un abierto de $\mathbb L^{\delta}$. Por lo tanto,
$C^{\dagger}$ es un abierto en $\mathbb{L}^\delta$, como queríamos.
\end{proof}



\begin{lemma}\label{lem:dagger-antimonotono}
Sea $\mathbb L$ un retículo distributivo acotado y sean 
$U,V\subseteq\mathcal X(\mathbb L^{-1})$.  
Entonces el operador $(-)^{\dagger}$ es antimonótono, es decir:
\[
U \subseteq V 
\quad\Longrightarrow\quad 
V^{\dagger} \subseteq U^{\dagger}.
\]
\end{lemma}

\begin{proof}[\textbf{Demostración:}]
Supongamos $U\subseteq V$ y sea $Q\in V^{\dagger}$.  
Por definición de $(-)^{\dagger}$, existe $P\in\mathcal X(\mathbb L^{-1})$
tal que $Q=P^{c}$ y $P\notin V$.
Como $U\subseteq V$, de $P\notin V$ se sigue en particular que 
$P\notin U$.  
Por lo tanto, nuevamente por la definición de $(-)^{\dagger}$, se tiene
$Q=P^{c}\in U^{\dagger}$.
Así obtenemos $V^{\dagger}\subseteq U^{\dagger}$, como se quería.
\end{proof}


\begin{lemma}\label{lem:inclusion-estrella}
Sea $a\in L$ y $C\in\mathcal{C}((\mathbb{L}^{-1})^\delta)$.
Entonces
\[
C\subseteq\sigma_{\mathbb L^{-1}}(a)
\quad\Longleftrightarrow\quad
\sigma_{\mathbb L}(a)\subseteq C^{\dagger}.
\]
\end{lemma}

\begin{proof}[\textbf{Demostración:}] Probaremos ambas implicaciones

\noindent\((\Rightarrow)\) 
Supongamos $C\subseteq\sigma_{\mathbb L^{-1}}(a)$.
Sea $P\in\sigma_{\mathbb L}(a)$, esto es, $a\in P$.
Queremos ver que $P\in C^{\dagger}$, es decir, que $P^{c}\notin C$.

Si ocurriera $P^{c}\in C$, como $C\subseteq\sigma_{\mathbb L^{-1}}(a)$,
tendríamos $a\in P^{c}$, lo que contradice $a\in P$ (pues $P$ y $P^{c}$ son disjuntos).
Por tanto $P^{c}\notin C$ y, por definición de $C^{\dagger}$, se sigue
$P\in C^{\dagger}$. Como $P\in\sigma_{\mathbb L}(a)$ fue arbitrario,
concluimos $\sigma_{\mathbb L}(a)\subseteq C^{\dagger}$.

\medskip
\noindent\((\Leftarrow)\) 
Supongamos ahora $\sigma_{\mathbb L}(a)\subseteq C^{\dagger}$ y sea
$I\in C$.  
Por el Lema~\ref{lem:primos-op-complemento}, existe un filtro primo
$P\in\mathcal X(\mathbb L)$ tal que $I=P^{c}$.

Si $a\notin I$, entonces $a\in P$, es decir $P\in\sigma_{\mathbb L}(a)$.
Por hipótesis, $\sigma_{\mathbb L}(a)\subseteq C^{\dagger}$ implica
$P\in C^{\dagger}$, esto es, $P^{c}\notin C$, es decir $I\notin C$,
lo cual contradice $I\in C$.
Por lo tanto debe cumplirse $a\in I$, y así $I\in\sigma_{\mathbb L^{-1}}(a)$. Como $I\in C$ fue arbitrario, se concluye
$C\subseteq\sigma_{\mathbb L^{-1}}(a)$.
\end{proof}


Además, para cada filtro primo $Q\in\mathcal X(\mathbb L)$ vale:

\begin{lemma}\label{lem:QnotinCstar}
Sea $C\in\mathcal{C}((\mathbb{L}^{-1})^\delta)$ y $Q\in\mathcal X(\mathbb L)$.
Entonces
\[
Q\notin C^{\dagger}
\quad\Longleftrightarrow\quad
Q^{c}\in C.
\]
\end{lemma}

\begin{proof}[\textbf{Demostración:}]
Recordemos que, por definición de $C^{\dagger}$,
\[
C^{\dagger}
=
\{\,P^{c} : P\notin C\,\}
=
\{\,R\in\mathcal X(\mathbb L) : R^{c}\notin C\,\}.
\]
En particular, para un filtro primo fijo $Q\in\mathcal X(\mathbb L)$,
\[
Q\in C^{\dagger}
\quad\Longleftrightarrow\quad
Q^{c}\notin C.
\]
Tomando negaciones a ambos lados, obtenemos inmediatamente
\[
Q\notin C^{\dagger}
\quad\Longleftrightarrow\quad
Q^{c}\in C,
\]
como queríamos.
\end{proof}
Ahora si estamos en condiciones de comenzar un desarrollo similar al de los tres casos anteriores.
Comenzaremos estudiando cómo el operador $m^\sigma$ se define sobre un cerrado $C\in\mathcal{C}((\mathbb{L}^{-1})^\delta)$.  
Para un operador normal de tipo $(-1,1)$, la extensión canónica
\[
m^\sigma:\Up(\mathcal{X}(\mathbb{L}^{-1}))
\longrightarrow
\Up(\mathcal{X}(\mathbb{L}))
\]
viene dada por
\[
m^\sigma(C)
=
\bigcap_{\Up(\mathcal X(\mathbb L))}
\Bigl\{
\sigma_{\mathbb L}(m(a)) 
:\;
C\subseteq_{\mathbb L^{-1}}\sigma_{\mathbb L^{-1}}(a)
\Bigr\}.
\]

Para un punto $P\in\mathcal X(\mathbb L)$ obtenemos:
\begin{align*}
P\in m^\sigma(C)
&\;\Longleftrightarrow\;
\forall\,a\in L\,
\Bigl(
C\subseteq_{\mathbb L^{-1}}\sigma_{\mathbb L^{-1}}(a)
\Rightarrow
P\in\sigma_{\mathbb L}(m(a))
\Bigr).
\end{align*}

Usando el 
Lema~\ref{lem:inclusion-estrella}, la condición
\[
C\subseteq_{\mathbb L^{-1}}\sigma_{\mathbb L^{-1}}(a)
\]
es equivalente a
\[
\sigma_{\mathbb L}(a)\subseteq_{\mathbb L} C^{\dagger}.
\]
Por lo tanto:
\begin{align*}
P\in m^\sigma(C)
&\;\Longleftrightarrow\;
\forall\,a\in L\,
\Bigl(
\sigma_{\mathbb L}(a)\subseteq_{\mathbb L}C^{\dagger}
\Rightarrow
m(a)\in P
\Bigr)\\[0.4em]
&\;\Longleftrightarrow\;
\forall\,a\in\Gamma(C^{\dagger}),\; m(a)\in P
\qquad\text{(por definición de }\Gamma(C^{\dagger})\text{)}\\[0.4em]
&\;\Longleftrightarrow\;
\Gamma(C^{\dagger})
\subseteq_{\mathbb L} m^{-1}[P]\\[0.4em]
&\;\Longleftrightarrow\;
\Gamma(C^{\dagger})\cap m^{-1}[P]^{c}=\varnothing.
\end{align*}

Aquí $\Gamma(C^{\dagger})$ es el ideal asociado al abierto $C^{\dagger}$:
\[
\Gamma(C^{\dagger})
=\{\,a\in L : \sigma_{\mathbb L}(a)\subseteq_{\mathbb L} C^{\dagger}\,\}.
\]

Esto motiva la siguiente definición.

\begin{definition}
Sea $m:\mathbb L^{-1}\to\mathbb L$ un operador normal de tipo $(-1,1)$.
Definimos la multirrelación inducida por $m$ como
\[
(P,C)\in H_m
\quad\Longleftrightarrow\quad
\Gamma(C^{\dagger})\cap m^{-1}[P]^{c}=\varnothing.
\]
\end{definition}
Ahora identificaremos esta multirrelación con la relación dual $R_m$ introducida en el Capítulo~\ref{Ch 3}. Al ser $m$ un operador de tipo \((-1,1)\), $(P,Q)\in R_m\iff Q^c\subseteq m^{-1}[P]$.
\begin{lemma}\label{lem:M-R-estrella}
Sea $m:\mathbb L^{-1}\to\mathbb L$ un operador normal de tipo $(-1,1)$.
Sea $C\in\mathcal{C}((\mathbb{L}^{-1})^\delta)$ y 
$P\in\mathcal X(\mathbb L)$.
Entonces son equivalentes:
\begin{enumerate}
    \item $(P,C)\in H_m$.
    \item Existe $Q\in\mathcal X(\mathbb L)$ tal que
    \[
    (P,Q)\in R_m
    \qquad\text{y}\qquad
    Q\notin C^{\dagger}.
    \]
\end{enumerate}
\end{lemma}

\begin{proof}[\textbf{Demostración:}]
Recordemos que $C^{\dagger}$ es un abierto de $\mathbb{L}^\delta$
y que, para todo abierto $U$ y todo filtro primo $Q\in\mathcal X(\mathbb L)$,
se cumple
\[
Q\in U
\quad\Longleftrightarrow\quad
Q\cap\Gamma(U)\neq\varnothing.
\]

\medskip
\noindent
$(1)\Rightarrow(2)$.
Supongamos $(P,C)\in H_m$, esto es,
\[
\Gamma(C^{\dagger})\cap m^{-1}[P]^{c}=\varnothing.
\]
Como $\Gamma(C^{\dagger})$ es un ideal de $\mathbb L$ y $m^{-1}[P]^{c}$ es un
filtro de $\mathbb L$, por el Teorema del Filtro Primo existe un filtro
primo $Q\in\mathcal X(\mathbb L)$ tal que
\[
m^{-1}[P]^{c}\subseteq Q
\qquad\text{y}\qquad
Q\cap\Gamma(C^{\dagger})=\varnothing.
\]

De la inclusión $m^{-1}[P]^{c}\subseteq Q$, al tomar complementos, se obtiene
\[
Q^{c}\subseteq m^{-1}[P],
\]
lo que equivale a $(P,Q)\in R_m$ por la definición de la relación dual
para el tipo $(-1,1)$.

Por otra parte, de $Q\cap\Gamma(C^{\dagger})=\varnothing$ se deduce
\[
Q\notin C^{\dagger},
\]
ya que, si $Q\in C^{\dagger}$, por la caracterización anterior tendríamos
$Q\cap\Gamma(C^{\dagger})\neq\varnothing$, contradicción.  
Así obtenemos un filtro primo $Q$ con $(P,Q)\in R_m$ y $Q\notin C^{\dagger}$.

\medskip
\noindent
$(2)\Rightarrow(1)$.
Supongamos ahora que existe $Q\in\mathcal X(\mathbb L)$ tal que
\[
(P,Q)\in R_m
\qquad\text{y}\qquad
Q\notin C^{\dagger}.
\]
De $(P,Q)\in R_m$ se sigue, por definición de $R_m$, que
\[
Q^{c}\subseteq m^{-1}[P],
\quad\text{es decir}\quad
m^{-1}[P]^{c}\subseteq Q.
\]

Como $Q\notin C^{\dagger}$, por la caracterización de pertenencia a abiertos,
se tiene
\[
Q\cap\Gamma(C^{\dagger})=\varnothing.
\]

Supongamos, hacia una contradicción, que
\(\Gamma(C^{\dagger})\cap m^{-1}[P]^{c}\neq\varnothing\).
Entonces existe $a\in\Gamma(C^{\dagger})\cap m^{-1}[P]^{c}$.  
Pero $m^{-1}[P]^{c}\subseteq Q$ implica $a\in Q$, y a la vez
$a\in\Gamma(C^{\dagger})$, con lo cual
\[
a\in Q\cap\Gamma(C^{\dagger}),
\]
contradiciendo que $Q\cap\Gamma(C^{\dagger})=\varnothing$. Por lo tanto,
\[
\Gamma(C^{\dagger})\cap m^{-1}[P]^{c}=\varnothing,
\]
es decir, $(P,C)\in H_m$.
\end{proof}
Veamos por último que lo operadores $m^\sigma$ y $m_{R_m}$ coinciden sobre los clopens $U$ de la extensión canónica de $\mathbb{L}^{-1}$, recordando que en este caso particular, $m_{R_m}(U)=R_m^{-1}(U^c)$.
\begin{theorem}\label{teo:igualdad-operadores-estrella}
Sea $m:\mathbb L^{-1}\to\mathbb L$ un operador normal de tipo $(-1,1)$.
Entonces, para todo $U$ clopen de la extensión canónica de $\mathbb L^{-1}$ y todo
$P\in\mathcal X(\mathbb L)$, se verifica
\[
P\in m^\sigma(U)
\quad\Longleftrightarrow\quad
P\in m_{R_m}(U^{\dagger}),
\]
\end{theorem}

\begin{proof}[\textbf{Demostración:}]
Recordemos que $m^\sigma$ en el caso $(-1,1)$ es determinada por la unión de sus valores sobre los cerrados contenidos en $U$, es decir que
\[
P\in m^\sigma(U)
\;\Longleftrightarrow\;
\exists\, C\subseteq U\ \text{cerrado tal que } P\in m^\sigma(C).
\]

Por la caracterización puntual obtenida para el caso $(-1,1)$,
\[
P\in m^\sigma(C)
\;\Longleftrightarrow\;
(P,C)\in H_m,
\]
donde $H_m$ es la multirrelación inducida por
\[
(P,C)\in H_m
\;\Longleftrightarrow\;
\Gamma(C^{\dagger})\cap m^{-1}[P]^{c}=\varnothing.
\]

Aplicando el Lema~\ref{lem:M-R-estrella}, obtenemos:
\[
(P,C)\in H_m
\;\Longleftrightarrow\;
\exists\,Q\in\mathcal X(\mathbb L)\ 
\text{tal que }(P,Q)\in R_m\ \text{y}\ Q\notin C^{\dagger}.
\]

Como $U$ es cerrado y $C\subseteq U$, por el Lema \ref{lem:dagger-antimonotono} tenemos que $U^{\dagger}\subseteq C^{\dagger}$, y por contrarrecíproca,
$Q\notin C^{\dagger}$ implica $Q\notin U^{\dagger}$.
Tomando $C=U$ se obtiene la equivalencia
\[
\exists\,C\subseteq U\ \text{cerrado con }Q\notin C^{\dagger}
\;\Longleftrightarrow\;
Q\notin U^{\dagger}.
\]

Sustituyendo:
\[
P\in m^\sigma(U)
\;\Longleftrightarrow\;
\exists\,Q\in\mathcal X(\mathbb L)\ 
\text{tal que }(P,Q)\in R_m\ \text{y}\ Q\notin U^{\dagger}.
\]

Por el Lema \ref{lem:QnotinCstar},
\[
Q\notin U^{\dagger}
\;\Longleftrightarrow\;
Q^{c}\in U,
\]
y, equivalentemente,
\[
(U^{\dagger})^{c}
=
\{\,Q\in\mathcal X(\mathbb L):Q^{c}\in U\,\}.
\]

De este modo,
\[
P\in m^\sigma(U)
\;\Longleftrightarrow\;
\exists\,Q\in (U^{\dagger})^{c}\ 
\text{con }(P,Q)\in R_m.
\]

Pero esto significa exactamente
\[
P\in R_m^{-1}\bigl((U^{\dagger})^{c}\bigr)
=
m_{R_m}(U^{\dagger}),
\]
pues $m_{R_m}(V)=R_m^{-1}(V^{c})$ para todo $V\subseteq\mathcal X(\mathbb L)$. Se concluye así
\[
P\in m^\sigma(U)
\quad\Longleftrightarrow\quad
P\in m_{R_m}(U^{\dagger}),
\]
como se quería.
\end{proof}

En esta sección hemos analizado, de manera exhaustiva, la relación entre la extensión canónica de los operadores normales y la correspondencia relacional que aparece en la dualidad de Priestley
para las clases $\BDL(i,j)$.

Para cada uno de los cuatro tipos posibles $(i,j)\in\{-1,1\}^2$,
hemos mostrado que la acción puntual del operador extendido
$m^\delta$ (ya que en los cuatro casos vimos que $m^\pi=m^\sigma$)
puede describirse mediante una multirrelación natural sobre 
el espacio de Priestley correspondiente.
Dicha multirrelación, obtenida directamente de la definición de
la extensión canónica, coincide en todos los casos con la relación 
dual $R_m$ caracterizada abstractamente en el 
Capítulo~\ref{Ch 3}.

En otras palabras, la extensión canónica recupera, de manera uniforme y transparente,  
la misma relación $R_m$ que determina el funtor dual
\[
\BDL(i,j)\;\simeq\;\PS(i,j)^{op}.
\]
Esto establece una compatibilidad fundamental entre la teoría de
extensiones canónicas y la dualidad de Priestley con operadores,
garantizando que ambos enfoques describen el mismo comportamiento
semántico de los operadores normales sobre retículos distributivos. Este resultado completa el desarrollo del capítulo.






\chapter{Aplicaciones} \label{Ch 5}




En los capítulos anteriores desarrollamos el marco teórico que articula,
por un lado, las categorías de retículos distributivos acotados con
operadores normales de tipo $(i,j)$ y sus correspondientes
espacios de Priestley relacionales, y por otro, la extensión canónica
como mecanismo de completitud algebraica.
%
En este capítulo mostraremos cómo estos resultados se aplican a
situaciones concretas, tanto en el plano categórico como en el
topológico. Haremos uso principalmente de las dualidades presentadas en el Capítulo 3.


\section{Aplicación del isomorfismo
$\BDL(i,j)\cong\BDL(-i,-j)$}

En el Capítulo~3 se estableció que, para cada par 
\((i,j)\in\{-1,1\}^2\), existe un isomorfismo de categorías
\[
  (-)^*\colon \BDL(i,j)\;\longrightarrow\;\BDL(-i,-j),
\]
obtenido invirtiendo el orden subyacente del retículo.  
Dicho isomorfismo refleja una simetría fundamental: toda identidad algebraica que caracteriza a un operador de tipo \((i,j)\) tiene una contraparte exacta para un operador de tipo \((-i,-j)\).

En esta sección aplicaremos esta simetría al estudio de una familia particular de retículos con operador que resulta
estable bajo dicha dualidad: las \emph{Álgebras de Ockham}.


\subsection{Álgebras de Ockham como objetos de
$\BDL(1,-1)$ y $\BDL(-1,1)$}

\begin{definition}
Una \emph{álgebra de Ockham} es una estructura
\[
  (\mathbb{L},f)=(L,\vee,\wedge,0,1,f)
\]
que satisface las siguientes condiciones:
\begin{enumerate}
  \item $\mathbb{L}=(L,\vee,\wedge,0,1)$ es un retículo distributivo acotado.
  \item El operador unario $f\colon L\to L$ cumple
  \[
    f(0)=1,\qquad f(x\vee y)=f(x)\wedge f(y),
  \]
  y además
  \[
    f(1)=0,\qquad f(x\wedge y)=f(x)\vee f(y).
  \]
\end{enumerate}
\end{definition}
Así vemos que cada álgebra de Ockham parece que puede verse como un retículo distributivo acotado con un operador normal de tipo $(1,-1)$ y $(-1,1)$ simultáneamente. Esto lo desarrollaremos más adelante.  

\begin{proposition}\label{CategoriaOckham}
La clase de las álgebras de Ockham, con los homomorfismos de retículos acotados que preservan el operador $f$, forma una categoría, cuya composición de flechas está dada por la composición usual de funciones y la flecha identidad está dada por la función identidad.
\end{proposition}

\begin{proof}[\textbf{Demostración}]
Dado que la composición de flechas se define como la composición usual de funciones, bastará comprobar que la composición de dos homomorfismos de álgebras de Ockham es nuevamente un homomorfismo de álgebras de Ockham, y  la identidad de cada álgebra de Ockham es un homomorfismo de álgebras de Ockham. En efecto:

\begin{itemize}
\item[1.] \textbf{Composición de morfismos.}  

Sean $h\colon (\mathbb{L}_1,f_1)\to(\mathbb{L}_2,f_2)$ y $g\colon (\mathbb{L}_2,f_2)\to(\mathbb{L}_3,f_3)$ homomorfismos de álgebras de Ockham.  
\begin{itemize}
  \item Como $h$ y $g$ son homomorfismos de retículos acotados, su composición $g\circ h$ también lo es.
  \item Además, para todo $a\in L_1$:
  \[
    (g\circ h)\bigl(f_1(a)\bigr)
    = g\bigl(h(f_1(a))\bigr)
    = g\bigl(f_2(h(a))\bigr)
    = f_3\bigl(g(h(a))\bigr)
    = f_3\bigl((g\circ h)(a)\bigr).
  \]
  Por tanto, $g\circ h$ conmuta con los operadores.
\end{itemize}
Así, $g\circ h$ es un homomorfismo de álgebras de Ockham.

\item[2.] \textbf{Identidad.}  

Sea $(\mathbb{L},f)$ un álgebra de Ockham y definamos la función
\begin{align*}
\Id_L &\colon L\to L,\\
\Id_L&(a) = a.
\end{align*}

\noindent\emph{(a) Homomorfismo de retículos.}  
Para todo $a,b\in L$,
\begin{align*}
\Id_L(a\vee b)&=a\vee b=\Id_L(a)\vee\Id_L(b),\\
\Id_L(a\wedge b)&=a\wedge b=\Id_L(a)\wedge\Id_L(b),\\
\Id_L(0)&=0,\qquad \Id_L(1)=1.
\end{align*}

\noindent\emph{(b) Conmutatividad con el operador.}  
Para todo $a\in L$,
\[
  \Id_L\bigl(f(a)\bigr)=f(a)=f\bigl(\Id_L(a)\bigr).
\]

\noindent\emph{(c) Ley de identidad.}  
Sea $h\colon (\mathbb{L}_1,f_1)\to(\mathbb{L}_2,f_2)$ un homomorfismo de álgebras de Ockham. Entonces
\[
h\circ\Id_{L_1}=h,\qquad
\Id_{L_2}\circ h=h.
\]
De (a), (b) y (c) se concluye que para cada objeto $(\mathbb{L},f)$ la flecha
\[
\Id_{(\mathbb{L},f)}:=\Id_L
\]
es la flecha identidad.
\end{itemize}

\noindent De (1) y (2) se concluye que las álgebras de Ockham, con los homomorfismos que conmutan con su operador, forman efectivamente una categoría que nombraremos $\Ock$.
\end{proof}
Hemos mencionado que las álgebras de Ockham pueden interpretarse, indistintamente, como objetos de $\BDL(1,-1)$ o de $\BDL(-1,1)$, dependiendo de qué par de leyes se enfatice.
Esto sugiere que la categoría $\Ock$ se encuentra naturalmente “entre” ambas categorías, compartiendo su estructura básica pero añadiendo las dos condiciones duales de manera simultánea.
Con el fin de formalizar esta idea, introducimos a continuación los funtores de inclusión que expresan de manera precisa cómo $\Ock$ se encuentra en $\BDL(1,-1)$ y $\BDL(-1,1)$.


\subsection{Los funtores de inclusión y su relación}

Comenzaremos definiendo los pares de asignaciones
\[
I \colon \Ock \longrightarrow \BDL(1,-1),
\qquad
J \colon \Ock \longrightarrow \BDL(-1,1),
\]
que permitirán la inmersión de la categoría $\Ock$ en
$\BDL(1,-1)$ y $\BDL(-1,1)$.

\noindent\textbf{Sobre Objetos:}
\begin{itemize}
  \item Dado $(\mathbb{L},f)$ un álgebra de Ockham, definimos
  \[
  I(\mathbb{L},f) = (\mathbb{L},f)
  \qquad
  J(\mathbb{L},f) = (\mathbb{L},f)
  \]
  Es decir, cada álgebra de Ockham es vista como un mismo retículo
  con operador, interpretado bajo las leyes de tipo $(1,-1)$ o $(-1,1)$ respectivamente.
\end{itemize}

\noindent\textbf{Sobre Flechas:}
\begin{itemize}
  \item Dado un morfismo $h\colon(\mathbb{L}_1,f_1)\to(\mathbb{L}_2,f_2)$ en $\Ock$
  (homomorfismo de retículos acotados tal que $h\circ f_1 = f_2\circ h$),
  definimos
  \[
  I(h)=h,\qquad J(h)=h.
  \]
  La acción sobre flechas coincide con la función subyacente, pues
  las condiciones de homomorfismo son idénticas en las tres categorías.
\end{itemize}

\begin{lemma}\label{lem:funtores-I-J}
Las asignaciones anteriores determinan funtores covariantes
\[
I\colon \Ock\to\BDL(1,-1)
\quad\text{y}\quad
J\colon \Ock\to\BDL(-1,1),
\]
\[
\begin{tikzcd}[column sep=large]
& \Ock \arrow[dl,hook',swap,"I"] \arrow[dr,hook,"J"] & \\
\BDL(1,-1) \arrow[rr,"(-)^*"'] && \BDL(-1,1)
\end{tikzcd}
\]
\end{lemma}

\begin{proof}[\textbf{Demostración}]
Verifiquemos las propiedades funtoriales para $I$; el caso de $J$ es análogo.

\begin{enumerate}
  \item Buena definición sobre objetos.
  
  Si $(\mathbb{L},f)$ es un álgebra de Ockham, por definición $\mathbb{L}$ es un retículo distributivo acotado y $f$ cumple simultáneamente las leyes
  \[
  f(0)=1,\ f(x\vee y)=f(x)\wedge f(y),
  \qquad
  f(1)=0,\ f(x\wedge y)=f(x)\vee f(y).
  \]
  En particular, la primera pareja de ecuaciones muestra que
  $(\mathbb{L},f)$ es un $(1,-1)$-retículo, por lo que $I(\mathbb{L},f)$
  está bien definido como objeto de $\BDL(1,-1)$.
  Análogamente, la segunda pareja asegura que
  $J(\mathbb{L},f)$ es un objeto de $\BDL(-1,1)$.

  \item Buena definición sobre flechas.
  
  Sea $h\colon(\mathbb{L}_1,f_1)\to(\mathbb{L}_2,f_2)$ un morfismo en $\Ock$.
  Como $h$ es homomorfismo de retículos acotados que conmuta con $f$,
  también es un homomorfismo en $\BDL(1,-1)$ y en $\BDL(-1,1)$.
  Por tanto, $I(h)$ y $J(h)$ están bien definidos.

  \item Preservación de la identidad.
  
  Para toda álgebra de Ockham $(\mathbb{L},f)$,
  \[
  I(\Id_{(\mathbb{L},f)}) = \Id_{(\mathbb{L},f)} =
  \Id_{I(\mathbb{L},f)},
  \qquad
  J(\Id_{(\mathbb{L},f)}) = \Id_{J(\mathbb{L},f)}.
  \]

  \item Preservación de la composición.
  
  Si $h\colon(\mathbb{L}_1,f_1)\to(\mathbb{L}_2,f_2)$ y
  $g\colon(\mathbb{L}_2,f_2)\to(\mathbb{L}_3,f_3)$ son morfismos en $\Ock$,
  entonces
  \[
  I(g\circ h)=g\circ h = I(g)\circ I(h),
  \qquad
  J(g\circ h)=g\circ h = J(g)\circ J(h).
  \]
\end{enumerate}
Por tanto, $I$ y $J$ son funtores covariantes.
\end{proof}
\begin{lemma}\label{lem:Ock-plena}
$\Ock$ es una subcategoría plena de $\BDL(1,-1)$ y de $\BDL(-1,1)$.
\end{lemma}
\begin{proof}[\textbf{Demostración}]
Se sigue de que $I$ y $J$ actúan como la identidad sobre morfismos.
\end{proof}

Las inclusiones anteriores permiten describir la posición de $\Ock$ respecto de las categorías $\BDL(1,-1)$ y $\BDL(-1,1)$.
No obstante, el isomorfismo de categorías $(-)^*$ definido en el Capítulo 3, actúa sobre ambas categorías y, por tanto, debería tener un efecto previsible sobre los objetos de $\Ock$.
El siguiente resultado va más allá y muestra que el isomorfismo de categorías $(-)^*$ definido en el Capítulo 3 preserva exactamente la categoría de las álgebras de Ockham, lo que explica por qué éstas pueden considerarse “estables” bajo el isomorfismo $(-)^*$.


\begin{theorem}\label{thm:Ock-intermedia-formal}
Sea $(-)^*\colon \BDL(1,-1)\to\BDL(-1,1)$ el isomorfismo de categorías definido en el Capítulo~3.
Entonces $(-)^*$ preserva la subcategoría de las álgebras de Ockham, es decir,
\[
(\Ock)^*=\Ock.
\]
\end{theorem}

\begin{proof}[\textbf{Demostración}]
Recordemos la acción de $(-)^*$:
\begin{itemize}
  \item Sobre objetos: $(\mathbb{L},m)\mapsto (\mathbb{L}^{-1},m^*)$ con $m^*=m$.
  \item Sobre flechas: $h\mapsto h$ (misma función subyacente).
\end{itemize}
\begin{enumerate}

\item  $(-)^*$ envía objetos Ockham a objetos Ockham.

Sea $(\mathbb{L},f)$ un objeto de $\Ock$.
Por definición, $f$ satisface simultáneamente
\[
    f(0)=1,\qquad f(x\vee y)=f(x)\wedge f(y),
  \]
  y
  \[
    f(1)=0,\qquad f(x\wedge y)=f(x)\vee f(y).
  \]
Al aplicar $(-)^*$ obtenemos el objeto $(\mathbb{L}^{-1},f^*)=(\mathbb{L}^{-1},f)$ en $\BDL(-1,1)$ y se verifica, para todo $a,b\in L$,
\[
f(0^{\mathbb{L}^{-1}})=f(1^{\mathbb{L}})=0^{\mathbb{L}}=1^{\mathbb{L}^{-1}},
\qquad
f(a\vee^{\mathbb{L}^{-1}} b)=f(a\wedge^{\mathbb{L}} b)
= f(a)\vee^{\mathbb{L}} f(b)=f(a)\wedge^{\mathbb{L}^{-1}} f(b),
\]
y
\[
f(1^{\mathbb{L}^{-1}})=f(0^{\mathbb{L}})=1^{\mathbb{L}}=0^{\mathbb{L}^{-1}},
\qquad
f(a\wedge^{\mathbb{L}^{-1}} b)=f(a\vee^{\mathbb{L}} b)
= f(a)\wedge^{\mathbb{L}} f(b)=f(a)\vee^{\mathbb{L}^{-1}} f(b).
\]
Es decir, $(\mathbb{L}^{-1},f)$ vuelve a satisfacer ambas familias de ecuaciones de Ockham.
Por consiguiente, la imagen de cualquier objeto de $\Ock$ por $(-)^*$ es nuevamente un objeto de $\Ock$.

\smallskip
\item $(-)^*$ envía morfismos Ockham a morfismos Ockham.

Sea $h:(\mathbb{L}_1,f_1)\to(\mathbb{L}_2,f_2)$ un morfismo en $\Ock$; por definición,
$h$ es homomorfismo de retículos acotados y $h\circ f_1=f_2\circ h$.
La imagen por $(-)^*$ es la misma función $h$, ahora vista entre $(\mathbb{L}_1^{-1},f_1)$ y $(\mathbb{L}_2^{-1},f_2)$.
Como $h$ preserva $\vee,\wedge,0,1$ en $\mathbb{L}_1,\mathbb{L}_2$, preserva también $\vee,\wedge,0,1$ en los órdenes invertidos, y la conmutatividad con el operador se conserva:
\[
h\circ f_1=f_2\circ h.
\]
Luego $h$ es un morfismo entre objetos Ockham en la imagen.

\smallskip
\item Equivalencia sobre $\Ock$.

Por el mismo argumento de los Pasos 1 y 2 aplicado a $(-)_*$ (definido también en el Capítulo~3), éste también se puede restringir a un funtor $\Ock\to\Ock$.
Las identidades naturales de la equivalencia se restringen asimismo, de modo que
\[
(-)^* \;:\; \Ock \;\longrightarrow\; \Ock\qquad(-)_* \;:\; \Ock \;\longrightarrow\; \Ock
\]
y por lo tanto $(-)^*$ es una equivalencia.

\end{enumerate}
\end{proof}

El resultado anterior establece que las álgebras de Ockham permanecen invariantes bajo la equivalencia $(-)^*$.
Sin embargo, esta estabilidad no se extiende a toda la categoría $\BDL(1,-1)$, ya que existen retículos con operador que satisfacen únicamente una de las dos familias de ecuaciones.
El siguiente ejemplo ilustra precisamente este fenómeno, mostrando un objeto de $\BDL(1,-1)$ que no pertenece a $\Ock$ ni es isomorfo a ninguno de sus objetos.

\begin{example}\label{ex:bdl11-no-ock}
Sea $\mathbb{L}=(\{0<1\},\vee,\wedge,0,1)$ el retículo booleano de dos elementos y defínase
$m\colon \mathbb{L}\to \mathbb{L}$ por
\[
m(0)=1,\qquad m(1)=1.
\]
Entonces $(\mathbb{L},m)$ es un objeto de la categoría $\BDL(1,-1)$ pero $(\mathbb{L},m)$ no de la categoría $\Ock$, y no es isomorfo a ninguna álgebra de Ockham.
\end{example}

\begin{proof}[\textbf{Demostración}]
Analicemos parte por parte.
\begin{enumerate}
    \item $(\mathbb{L},m)$ es un $(1,-1)$-retículo.
    
Claramente $m(0)=1$. Además, para cualesquiera $a,b\in L$, se tiene
\[
m(a\vee b)=1
\qquad\text{y}\qquad
m(a)\wedge m(b)=1\wedge 1=1,
\]
luego $m(a\vee b)=m(a)\wedge m(b)$. Por tanto $m$ es de tipo $(1,-1)$.

\item $(\mathbb{L},m)$ no es un álgebra de Ockham.

En una álgebra de Ockham debe valer, además, $f(1)=0$ y $f(x\wedge y)=f(x)\vee f(y)$.
En nuestro caso $m(1)=1\neq 0$, por lo que falla la primera ecuación.

\item No es isomorfo a ninguna álgebra de Ockham.

Sea $(\mathbb{C},f)$ un álgebra de Ockham y supongamos que existe un isomorfismo de álgebras con operador
\[
\varphi\colon (\mathbb{L},m)\xrightarrow{\;\cong\;}(\mathbb{C},f).
\]
Como $\mathbb{L}$ es el único retículo distributivo acotado de dos elementos, necesariamente
$\mathbb{C}$ es también de dos elementos y $\varphi$ es un isomorfismo de retículos acotados.
Todo automorfismo de un retículo acotado preserva $0$ y $1$, de donde en el caso de dos elementos
$\varphi=\Id$. La compatibilidad con el operador exige entonces
\[
f = \varphi\circ m\circ \varphi^{-1} = m,
\]
pero esto es imposible porque en toda álgebra de Ockham de dos elementos se tiene $f(1)=0$,
mientras que $m(1)=1$. Concluimos que $(\mathbb{L},m)$ no es isomorfo a ninguna álgebra de Ockham.
\end{enumerate}
\end{proof}

\begin{remark}
El ejemplo muestra que la inclusión $I\colon \Ock\hookrightarrow \BDL(1,-1)$ es propia. Se puede plantear un contrajemplo análogo para mostrar lo mismo de la inclusión $J\colon \Ock\hookrightarrow \BDL(-1,1)$.
\end{remark}
El análisis precedente muestra que las álgebras de Ockham constituyen el punto de equilibrio entre las categorías $\BDL(1,-1)$ y $\BDL(-1,1)$. Ambas categorías resultan equivalentes por el funtor $(-)^*$, y $\Ock$ aparece como su subcategoría plena común, estable bajo dicha equivalencia. De este modo, la simetría $(i,j)\leftrightarrow(-i,-j)$ encuentra en $\Ock$ su expresión más exacta: el lugar donde ambas leyes, duales entre sí, se satisfacen simultáneamente.

\section{Aplicación de la dualidad de Priestley
$\BDL(1,1)\cong\PS(1,1)^{op}$}

La dualidad de Priestley estudiada en el Capítulo \ref{Ch 3} permite interpretar operadores algebraicos de tipo $(1,1)$ como relaciones binarias sobre el espacio de Priestley asociado. Esto nos proporciona una traducción entre ecuaciones algebraicas y propiedades topológicas. En esta sección estudiaremos estas correspondencias en detalle, mostrando cómo las propiedades algebraicas de un operador $m$ se reflejan en condiciones topológicas de la relación $R_m$ en el espacio dual.

\subsection{El marco de la correspondencia dual}

Sea $(\mathbb{L},m)$ un $(1,1)$–retículo distributivo acotado y $(\mathcal{X}(\mathbb{L}),\tau,\subseteq,R_m)$ su espacio de Priestley asociado.  Por construcción, la relación
$R_m\subseteq \mathcal{X}(\mathbb{L})\times \mathcal{X}(\mathbb{L})$ se define por
\[
(P,Q)\in R_m
\quad\text{si y sólo si}\quad
Q\subseteq m^{-1}[P],
\]
y cumple, por el Lema \ref{Lema de Representacion de BDL(i,j)}, que para todo conjunto clopen creciente $U\subseteq X$
\[
\sigma_{\mathbb{L}}(m(a))=R_m^{-1}(\sigma_{\mathbb{L}}(a)), \qquad
m_\sigma(U)=\{x\in X: R_m^{-1}(x)\subseteq U\},
\]
donde $\sigma_{\mathbb{L}}:L\to \mathcal{D}(\mathcal{X}(\mathbb{L}))$ es el isomorfismo ya trabajado.

\subsection{Algunas correspondencias relacionales}

A continuación presentamos, una a una, las principales equivalencias
entre propiedades algebraicas del operador $m$ y propiedades topológicas
de la relación $R_m$.
En cada caso se ofrece la traducción y su demostración formal.

Las dos primeras propiedades que presentaremos son satisfechas por \emph{todo} operador normal de tipo $(1,1)$, mientras que las siguientes son propiedades adicionales que pueden o no verificarse, y cuya interpretación dual resulta especialmente significativa.


\begin{proposition}[Preservación de uniones finitas]
Sea $(\mathbb{L},m)$ un $(1,1)$–retículo y $R_m$ su relación dual. Entonces
\[
m(a\vee b)=m(a)\vee m(b)
\quad\Longleftrightarrow\quad
R_m^{-1}(U\cup V)=R_m^{-1}(U)\cup R_m^{-1}(V)
\]
para todo par $U,V\in \mathcal{D}(\mathcal{X}(\mathbb{L}))$. Es decir, la preservación de uniones finitas por $m$
equivale a que $R_m^{-1}[\cdot]$ preserve uniones de clopens crecientes, o lo que es lo mismo, $R_m$ es una relación clopen.
\end{proposition}

\begin{proof}[{\textbf{Demostración:}}] Probaremos las dos implicaciones.

\begin{itemize}
    \item $(\Rightarrow)$
Si $m(a\vee b)=m(a)\vee m(b)$, aplicando $\sigma_{\mathbb{L}}$ y la igualdad
$\sigma_{\mathbb{L}}(m(x))=R_m^{-1}[\sigma_{\mathbb{L}}(x)]$ obtenemos:
\begin{align*}
R_m^{-1}(\sigma_{\mathbb{L}}(a\vee b))
&=\sigma_{\mathbb{L}}(m(a\vee b))
=\sigma_{\mathbb{L}}(m(a)\vee m(b))\\
&=\sigma_{\mathbb{L}}(m(a))\cup\sigma_{\mathbb{L}}(m(b))
=R_m^{-1}(\sigma_{\mathbb{L}}(a))\cup R_m^{-1}(\sigma_{\mathbb{L}}(b)).
\end{align*}
Dado que $\sigma_{\mathbb{L}}(a\vee b)=\sigma_{\mathbb{L}}(a)\cup\sigma_{\mathbb{L}}(b)$, la igualdad vale para
$U=\sigma_{\mathbb{L}}(a)$ y $V=\sigma_{\mathbb{L}}(b)$.

\item $(\Leftarrow)$
Recíprocamente, si $R_m^{-1}(U\cup V)=R_m^{-1}(U)\cup R_m^{-1}(V)$
para todos los clopens crecientes $U,V$, tomando
$U=\sigma_{\mathbb{L}}(a)$ y $V=\sigma_{\mathbb{L}}(b)$ se sigue:
\[
\sigma_{\mathbb{L}}(m(a\vee b))
=R_m^{-1}(\sigma_{\mathbb{L}}(a\vee b))
=R_m^{-1}(\sigma_{\mathbb{L}}(a))\cup R_m^{-1}(\sigma_{\mathbb{L}}(b))
=\sigma_{\mathbb{L}}(m(a)\vee m(b)).
\]
Puesto que $\sigma_{\mathbb{L}}$ es un isomorfismo, $m(a\vee b)=m(a)\vee m(b)$.
\end{itemize}
\end{proof}


\begin{proposition}[Caso $m(0)=0$]
Sea $(\mathbb{L},m)$ un $(1,1)$–retículo y $R_m$ su relación dual. Entonces
\[
m(0)=0
\quad\Longleftrightarrow\quad
R_m^{-1}(\emptyset)=\emptyset.
\]
Es decir, la condición $m(0)=0$ significa que $R_m$ no tiene preimágenes vacías.
\end{proposition}

\begin{proof}[\textbf{Demostración}]
Como $\sigma_{\mathbb{L}}(0)=\emptyset$, se tiene
$\sigma_{\mathbb{L}}(m(0))=R_m^{-1}(\sigma_{\mathbb{L}}(0))=R_m^{-1}(\emptyset)$.
Si $m(0)=0$, la imagen es vacía, y viceversa.
Por la inyectividad de $\sigma_{\mathbb{L}}$, las condiciones son equivalentes.
\end{proof}

\begin{proposition}[Idempotencia]
Sea $(\mathbb{L},m)$ un $(1,1)$–retículo y $R_m$ su relación dual. Entonces
\[
m(m(a))=m(a)
\quad\Longleftrightarrow\quad
(R_m\circ R_m)^{-1}=R_m^{-1},
\]
donde $R_m\circ R_m=\{(P,R):\exists Q\ (P,Q)\in R_m,\ (Q,R)\in R_m\}$. Es decir que la idempotencia algebraica 
se traduce en que el operador $R_m^{-1}$ es idempotente.
\end{proposition}

\begin{proof}[{\textbf{Demostración:}}]Probaremos las dos implicaciones.

\begin{itemize}
    \item \((\Rightarrow)\) Si \(m(m(a))=m(a)\) para todo \(a\), aplicando \(\sigma_{\mathbb{L}}\) se obtiene
\[
R_m^{-1}\big(R_m^{-1}(\sigma_{\mathbb{L}}(a))\big)
=\sigma_{\mathbb{L}}(m(m(a)))
=\sigma_{\mathbb{L}}(m(a))
=R_m^{-1}(\sigma_{\mathbb{L}}(a)).
\]
Como los conjuntos de la forma \(\sigma_{\mathbb{L}}(a)\) generan \(\mathcal D(\mathcal{X}(\mathbb L))\),
se tiene la igualdad de operadores \((R_m\circ R_m)^{-1}=R_m^{-1}\) sobre todo
\(\mathcal D(\mathcal{X}(\mathbb L))\).

\item \((\Leftarrow)\) Si \((R_m\circ R_m)^{-1}=R_m^{-1}\), entonces para todo \(a\in L\)
\[
\sigma_{\mathbb{L}}(m(m(a)))= (R_m\circ R_m)^{-1}(\sigma_{\mathbb{L}}(a)) = R_m^{-1}(\sigma_{\mathbb{L}}(a)) = \sigma_{\mathbb{L}}(m(a)).
\]
Puesto que \(\sigma_{\mathbb{L}}\) es un isomorfismo se concluye \(m(m(a))=m(a)\).
\end{itemize}
\end{proof}

\begin{proposition}[Extensividad]
Sea $(\mathbb{L},m)$ un $(1,1)$–retículo distributivo acotado y $R_m$ su relación dual. Entonces
\[
a\leq m(a)
\quad\Longleftrightarrow\quad
(\forall P\in\mathcal{X}(\mathbb{L}))\ (P,P)\in R_m,
\]
Es decir, la extensividad algebraica $a\le m(a)$ se corresponde con la \emph{reflexividad} de la relación dual $R_m$.
\end{proposition}

\begin{proof}[{\textbf{Demostración:}}]Probaremos las dos implicaciones.
\begin{itemize}
    \item $(\Rightarrow)$
Supongamos $a\leq m(a)$ para todo $a\in\mathbb{L}$.
Sea $P\in\mathcal{X}(\mathbb{L})$. Queremos ver que $(P,P)\in R_m$,
es decir, $P\subseteq m^{-1}[P]$.  
Si $a\in P$, entonces $a\le m(a)$ implica $m(a)\in P$
(porque $P$ es filtro primo). Por tanto $a\in m^{-1}[P]$, y de aquí
$P\subseteq m^{-1}[P]$.

\item $(\Leftarrow)$
Supongamos que $R_m$ es reflexiva, es decir,
$(P,P)\in R_m$ para todo $P\in\mathcal{X}(\mathbb{L})$.
Entonces $P\subseteq m^{-1}[P]$ para todo $P$.  
Si existiera $a$ con $a\not\le m(a)$, habría un filtro primo $P$
con $a\in P$ y $m(a)\notin P$, contradiciendo la inclusión anterior.
Por tanto $a\le m(a)$ para todo $a$.
\end{itemize}
\end{proof}
Las equivalencias obtenidas resumen el modo en que las identidades
algebraicas de un operador $(1,1)$ se reflejan en propiedades topológicas
de su relación dual. Es así como la dualidad de Priestley hace visible la
estructura subyacente a los operadores normales.




%\nocite{*}
\bibliographystyle{acm}
\bibliography{bibfile}


\end{document}